\documentclass[12pt]{article}
\usepackage[utf8]{inputenc}
\usepackage[T1]{fontenc}
\usepackage{amsmath, amssymb, amsthm, bm}
\usepackage{geometry}
\usepackage{booktabs}
\usepackage{array}
\newcolumntype{L}[1]{>{\raggedright\arraybackslash}p{#1}}
\usepackage{hyperref}

\geometry{margin=1.05in}
\hypersetup{colorlinks=true, linkcolor=blue, citecolor=blue, urlcolor=blue}

\title{Ballistic Evidence Accumulation with Decay:\\
       From the Linear Ballistic Accumulator to Six Decaying Extensions}
\author{}
\date{\today}

\newcommand{\dd}{\,\mathrm{d}}
\newcommand{\Normal}{\mathcal{N}}
\newcommand{\PhiN}{\Phi}
\newcommand{\phiN}{\phi}
\newcommand{\Phibar}{\bar{\Phi}}
\newcommand{\Gbar}{\bar{G}}
\newcommand{\e}{\mathrm{e}}

\theoremstyle{remark}
\newtheorem{remark}{Remark}
\newtheorem{proposition}{Proposition}

\begin{document}
\maketitle

\begin{abstract}
This note derives, from first principles, the response-time laws of seven
ballistic evidence-accumulation models: the linear ballistic accumulator (LBA)
and six decaying extensions of it. Five of the extensions differ in \emph{what}
the decay attacks. In the ballistic accumulator with leak (BAwL) it attacks the
accumulated state; in the ballistic accumulator with drive decay (BAwD) it
attacks the input, which fades while an independently scaled clearance drains
the accumulator; in the proportional-clearance variant (BAwDp) that clearance
is tied to the trialwise drive, so every trial integrates one common waveform;
in the ballistic accumulator with global fading (BAwF) it attacks the readout,
attenuating the whole trace including the start point; and in the
ballistic accumulator with ramping clearance (BAwR) it is exogenous, a drain
that grows with elapsed time and consults neither the state nor the drive. The
sixth, the ballistic transient/sustained local race with leak (BTAwL), comes
from Smith's (1995) sustained-and-transient detector under the ballistic
simplification.  Its transient and sustained subraces both leak, while their
stimulus drives are respectively a pulse and a rising sustained profile.
BAwD and BAwF are treated as fixed-index families whose kernels range from
power-law decay to an exponential limit. All seven share one generative idea:
within each subrace the trajectory is a deterministic function of a launch
strength and a start point, so the likelihood follows from a change of
variables rather than from a first-passage argument; the full BTAwL local race
then races two independent such pairs. We show how the leaky
variant sits between the LBA and the leaky competing accumulator, with the leak
term retained and lateral inhibition deliberately absent, and how the BTAwL
subraces stand in the same relation to single-boundary
Ornstein--Uhlenbeck detectors driven by filtered stimulus profiles. We then derive the
live and frozen branches of each finite-endpoint model, its intrinsic
omissions, and the closed positive-launch frozen integrals that the power-law
and exponential kernels induce (ordinary and split-lognormal through Gaussian
tails, Weibull through incomplete gamma functions; the transient BTAwL frozen
branch is the explicit quadrature exception). Each model is shown to be an
ordinary LBA racing a moving boundary in a warped clock, and the induced
(clock, boundary) pair is what distinguishes the variants: the leak in BAwL,
BAwF, and both BTAwL subraces attack the start point, only BAwDp leaves the
boundary alone, and the finite endpoint belongs to the transient BTAwL
subrace (not to the full local race when its sustained allocation is positive).
We
catalogue reductions,
defective components, scale invariances, and the precise reason that normal,
log-scale, and Weibull launches differ on the staggered frozen branch (and, for
proportional clearance, when negative launches are retained). The exposition is
self-contained. A common subject-level operational-time warp can be applied to
each model: \(s=c_\eta(u)\), with \(\eta=0\) recovering the parent model
exactly and the response density carrying the corresponding Jacobian.
\end{abstract}

\tableofcontents
\newpage

%======================================================================
\section{Scope and Notation}
%======================================================================
Let $T$ be raw response time, $t_0$ a non-decision time, and
\begin{equation}
  u = T - t_0
\end{equation}
the evidence-accumulation time. Densities and distribution functions are
written in $u$ throughout; converting to $T$ is a shift.

Each accumulator carries a scalar state $X(u)$ that starts at a random point
and rises toward a response boundary. Following the usual convention we write
\begin{equation}
  z \sim U(0,A), \qquad b = B + A, \qquad 0 \le A \le b,
\end{equation}
so that $B$ is the distance from the top of the start-point range to the
boundary and $b$ is the distance from the bottom of it. The trialwise launch
strength (the drive, or drift rate) is a random variable $V$ drawn
independently of $z$, with density $g$, survivor $\Gbar(w) = P(V > w)$, and
support in $\mathbb{R}$ or $(0,\infty)$ depending on the launch law.

Standard normal density and distribution functions are $\phiN$ and $\PhiN$,
with $\Phibar = 1 - \PhiN$. We use four launch laws:
\begin{align}
  \text{normal:}\quad & V \sim \Normal(v, s_v^2),
    & \Gbar(w) &= \PhiN\!\left(\frac{v-w}{s_v}\right), \\
  \text{lognormal:}\quad & \log V \sim \Normal(\mu, \sigma^2),
    & \Gbar(w) &= \PhiN\!\left(\frac{\mu - \log w}{\sigma}\right).
\end{align}
For ordinary lognormal launches in BAwL, BAwD, BAwDD, BAwDp, BAwF, and BAwR,
the public launch pair is $(\texttt{mean},\texttt{cv})$, where
$m_L=E[V]$ is the arithmetic mean and $c_L=\operatorname{SD}(V)/E[V]$ is
the coefficient of variation. The kernels use the equivalent log-scale
coordinates
\begin{equation}
  \sigma^2 = \log(1+c_L^2),
  \qquad
  \mu = \log m_L - \tfrac12\sigma^2,
  \label{eq:lognormal-mean-cv}
\end{equation}
so the inverse map is $m_L=\exp(\mu+\sigma^2/2)$ and
$c_L=\sqrt{\operatorname{expm1}(\sigma^2)}$. Thus, in the formulas below,
$\mu$ and $\sigma$ are internal coordinates for these BAw* models. BTAwL is
the exception: its ordinary lognormal launch exposes $(\mu,\sigma)$ directly,
and its split-lognormal launch exposes $(\mu,\sigma,\delta)$.

For the mathematics below, the Weibull launch is written with the pair
$(m_W,r)=(\text{mean},\text{shape})$. The public model parameter vector names
these entries $\texttt{shape}$ and $\texttt{mean}$, in that order. Write $r>0$
for the shape,
$m_W>0$ for the arithmetic mean, and
\begin{equation}
  \theta_W = \frac{m_W}{\Gamma(1+1/r)}
\end{equation}
for the corresponding conventional scale. Below, $\mathcal W(m,r)$ denotes
a Weibull law with arithmetic mean $m$ and shape $r$ (scale
$m/\Gamma(1+1/r)$). Thus, for $w>0$,
\begin{align}
  V &\sim \mathcal W(m_W,r), \\
  \Gbar_W(w) &= \exp\!\left[-\left(\frac{w}{\theta_W}\right)^r\right],
  &g_W(w) &= \frac{r}{\theta_W}\left(\frac{w}{\theta_W}\right)^{r-1}
             \exp\!\left[-\left(\frac{w}{\theta_W}\right)^r\right].
  \label{eq:weibull-launch}
\end{align}
For $w\le0$, set $\Gbar_W(w)=1$. The mean parameter is therefore the
arithmetic mean itself, $E[V]=m_W$, rather than the conventional scale.
Let $\Gamma(s,x)=\int_x^\infty t^{s-1}\e^{-t}\dd t$ denote the upper
incomplete gamma function. For $x>0$ the Weibull tail moments and stop-loss
transform are
\begin{align}
  M_{q,W}(x)
    &\equiv E\!\left[V^q\mathbf 1\{V>x\}\right]
      =\theta_W^q\Gamma\!\left(1+\frac{q}{r},
                         \left(\frac{x}{\theta_W}\right)^r\right), \\
  C_W(x)
    &\equiv E[(V-x)_+]
      =M_{1,W}(x)-x\Gbar_W(x).
  \label{eq:weibull-tail-moments}
\end{align}
For $x\le0$, use $M_{q,W}(x)=\theta_W^q\Gamma(1+q/r)$ when the moment
exists and $C_W(x)=m_W-x$. The power and logarithmic primitives needed on
staggered frozen branches are, for $x>0$,
\begin{align}
  P_{\alpha,W}(x)
    &\equiv \int_x^\infty w^{-\alpha}\Gbar_W(w)\dd w
      =\frac{\theta_W^{1-\alpha}}{r}
        \Gamma\!\left(\frac{1-\alpha}{r},
                    \left(\frac{x}{\theta_W}\right)^r\right), \\
  L_W(x;\ell)
    &\equiv \int_x^\infty\log(w/\ell)\Gbar_W(w)\dd w
      =-\left.\frac{\partial P_{\alpha,W}(x)}{\partial\alpha}\right|_{\alpha=0}
        -\log\ell\,P_{0,W}(x).
  \label{eq:weibull-power-primitives}
\end{align}
The last line is a derivative of the same incomplete-gamma expression, so it
is closed in incomplete gamma functions (and their parameter derivative).
The split-lognormal launch is continuous on the log scale. Let $Y=\log V$, with
$\mu\in\mathbb{R}$, $\sigma>0$, and $\delta\in\mathbb{R}$, and define
\begin{equation}
  s_L=\sigma\e^{\delta/2},\qquad
  s_R=\sigma\e^{-\delta/2},\qquad
  a=\frac{s_L}{s_L+s_R}=\frac{1}{1+\e^{-\delta}}.
  \label{eq:split-shape}
\end{equation}
The split point $c$ is chosen so that the location parameter $\mu$ is the
exact median of $Y$:
\begin{equation}
  c=\begin{cases}
    \mu, & \delta=0,\\[2pt]
    \mu-s_L\PhiN^{-1}\!\left(\dfrac{1}{4a}\right), & \delta>0,\\[8pt]
    \mu-s_R\PhiN^{-1}\!\left(\frac{3}{4}-\frac{1}{4}\e^{\delta}\right), & \delta<0.
  \end{cases}
  \label{eq:split-median-point}
\end{equation}
Its density on the log scale is
\begin{equation}
 f_Y(y)=
 \begin{cases}
   \dfrac{2a}{s_L}\phiN\!\left(\dfrac{y-c}{s_L}\right), & y\le c,\\[8pt]
   \dfrac{2(1-a)}{s_R}\phiN\!\left(\dfrac{y-c}{s_R}\right), & y>c,
 \end{cases}
 \qquad
 g_{\rm sp}(w)=\frac{f_Y(\log w)}{w},\quad w>0.
 \label{eq:split-density}
\end{equation}
Equivalently, its log-scale distribution function is
\begin{equation}
 F_Y(y)=
 \begin{cases}
   2a\PhiN\!\left(\dfrac{y-c}{s_L}\right), & y\le c,\\[8pt]
   2(1-a)\PhiN\!\left(\dfrac{y-c}{s_R}\right)+2a-1, & y>c.
 \end{cases}
 \label{eq:split-cdf}
\end{equation}
The two half-normal pieces have masses $a$ and $1-a$, and
\eqref{eq:split-median-point} gives $P(Y\le\mu)=1/2$. Thus $\delta$ controls
skewness without moving the median; $\delta=0$ gives the ordinary lognormal
exactly in these log-scale coordinates. For an ordinary BAw* launch, the
corresponding public pair is $m_L=\exp(\mu+\sigma^2/2)$ and
$c_L=\sqrt{\operatorname{expm1}(\sigma^2)}$. The split launch is positive by
construction, so its truncation normalizer is $Z_+=1$.

The normal law may optionally be truncated to $V > 0$, in which case every
expression below acquires the normalizer $Z_+ = \PhiN(v/s_v)$. The ordinary
lognormal, split-lognormal, and Weibull laws are positive by construction, so
truncation is vacuous for all three and $Z_+ = 1$.
For the lognormal we will need the two arguments
\begin{equation}
  d_1(x) = \frac{\mu + \sigma^2 - \log x}{\sigma},
  \qquad
  d_2(x) = \frac{\mu - \log x}{\sigma},
  \label{eq:d1d2}
\end{equation}
so that, for $x>0$, $\Gbar(x) = \PhiN(d_2(x))$ and the partial expectation is
$E[V\,\mathbf{1}\{V>x\}] = \e^{\mu+\sigma^2/2}\PhiN(d_1(x))$. For $x\le0$,
$\Gbar(x)=1$ and the partial expectation is $\e^{\mu+\sigma^2/2}$.

For the split-lognormal launch it is useful to collect all required tail
moments in one notation. For $x>0$, put $y_x=\log x$ and, for any real $r$,
\begin{align}
 M_{r,L}(x) &\equiv E\!\left[V^r\mathbf{1}\{V>x,Y\le c\}\right]
 =\begin{cases}
  2a\e^{rc+r^2s_L^2/2}
  \left[\PhiN(-rs_L)-
  \PhiN\!\left(\dfrac{y_x-c-rs_L^2}{s_L}\right)\right],&y_x<c,\\[6pt]
  0,&y_x\ge c,
 \end{cases} \\
 M_{r,R}(x) &\equiv E\!\left[V^r\mathbf{1}\{V>x,Y>c\}\right]
 =\begin{cases}
  2(1-a)\e^{rc+r^2s_R^2/2}\Phibar(-rs_R),&y_x\le c,\\[6pt]
  2(1-a)\e^{rc+r^2s_R^2/2}
  \Phibar\!\left(\dfrac{y_x-c-rs_R^2}{s_R}\right),&y_x>c,
 \end{cases} \\
 M_r(x) &= M_{r,L}(x)+M_{r,R}(x).
 \label{eq:split-tail-moments}
\end{align}
The split survivor and its total moments are
\begin{equation}
  \Gbar_{\rm sp}(x)=M_0(x)\quad(x>0),\qquad
  M_r(0)=2a\e^{rc+r^2s_L^2/2}\PhiN(-rs_L)
       +2(1-a)\e^{rc+r^2s_R^2/2}\Phibar(-rs_R).
  \label{eq:split-survivor-moment}
\end{equation}
For $x\le0$, set $M_r(x)=M_r(0)$ and $\Gbar_{\rm sp}(x)=1$. In particular, the split stop-loss
transform and lower partial expectation are
\begin{align}
 C_{\rm sp}(x)&=E[(V-x)_+]
 =\begin{cases}M_1(x)-xM_0(x),&x>0,\\ M_1(0)-x,&x\le0,\end{cases}
 \label{eq:split-stoploss}\\
 P_{\rm sp}(x)&=E[(x-V)_+]
 =\begin{cases}x[1-M_0(x)]-[M_1(0)-M_1(x)],&x>0,\\ 0,&x\le0.
 \end{cases}
 \label{eq:split-put}
\end{align}
Thus the additive start-point integrals used below remain closed in normal
distribution functions. The power and logarithmic primitives needed by
staggered frozen branches can be written without introducing a new special
function:
\begin{align}
 P_\alpha^{\rm sp}(x)
   &\equiv \int_x^\infty w^{-\alpha}\Gbar_{\rm sp}(w)\dd w
    =\frac{M_{1-\alpha}(x)-x^{1-\alpha}M_0(x)}{1-\alpha},
    &&\alpha\ne1,\\
 P_1^{\rm sp}(x)
   &\equiv \int_x^\infty \frac{\Gbar_{\rm sp}(w)}{w}\dd w
    =\left.\frac{\partial M_r(x)}{\partial r}\right|_{r=0}
       -\log x\,M_0(x),\\
 L^{\rm sp}(x;\ell)
   &\equiv \int_x^\infty\log(w/\ell)\Gbar_{\rm sp}(w)\dd w
    =-\left.\frac{\partial P_\alpha^{\rm sp}(x)}{\partial\alpha}\right|_{\alpha=0}
       -\log\ell\,P_0^{\rm sp}(x).
 \label{eq:split-power-primitives}
\end{align}
Derivatives in the last two lines are taken from the explicit expressions in
\eqref{eq:split-tail-moments}; they involve only $\PhiN$ and $\phiN$. Also,
$P_0^{\rm sp}=C_{\rm sp}$, and at
$\delta=0$, $M_r$, $C_{\rm sp}$, $P_\alpha^{\rm sp}$ and $L^{\rm sp}$ reduce
to their ordinary lognormal counterparts.

%======================================================================
\section{Canonical Operational-Time Warp}
\label{sec:time-warp}
Each ballistic model admits one common time-transform parameter, \(\eta\),
which changes the clock rather than the dynamics of the accumulator.
Let
\begin{equation}
  u=T-t_0>0,
  \qquad
  \omega=\e^{\eta}>0,
  \qquad
  s=c_\eta(u)=\frac{(1+u)^\omega-1}{\omega},
  \label{eq:time-warp-forward}
\end{equation}
where \(u\) is physical accumulation time and \(s\) is the internal
operational time used by the parent model. Times are measured in seconds; the
fixed one-second reference in \(1+u\) is part of the definition, not an
additional parameter. The derivative and inverse are
\begin{equation}
  c_\eta'(u)=(1+u)^{\omega-1},
  \qquad
  c_\eta^{-1}(s)=(1+\omega s)^{1/\omega}-1.
  \label{eq:time-warp-jacobian}
\end{equation}
At \(\eta=0\) (equivalently \(\omega=1\)),
\(c_0(u)=u\), \(c_0'(u)=1\), and \(c_0^{-1}(s)=s\), so the parent model is
recovered exactly. The clock starts at ordinary speed,
\(c_\eta'(0)=1\): negative \(\eta\) progressively slows operational time
and positive \(\eta\) progressively accelerates it. This is therefore a
curvature parameter, not a constant rescaling of time.
The inverse clock has curvature
\begin{equation}
  \frac{\dd^2 c_\eta^{-1}(s)}{\dd s^2}
  =(1-\omega)(1+\omega s)^{1/\omega-2},
  \label{eq:time-warp-curvature}
\end{equation}
so negative \(\eta\) makes physical time convex in $s$, whereas positive
\(\eta\) makes it concave.

Take \(\eta_i\) to be a subject-level parameter shared across conditions,
accumulators, and (for BTAwL) both local subraces. No
\(\eta\times L\) effect or channel-specific warp is assumed. The parameter is
unconstrained on the real line, and \(\eta=0\) is the parent model.

\subsection{General composition theorem}
Suppose a parent ballistic model is written in operational time as
\begin{equation}
  \frac{\dd X_0}{\dd s}=a(X_0,s;\theta),
  \qquad
  X_0(s;V,z,\theta),
  \label{eq:time-warp-parent}
\end{equation}
and has required launch strength \(V_0^*(s,z)\). The warped trajectory and
physical-time generator are
\begin{equation}
  X_\eta(u;V,z,\theta)=X_0(c_\eta(u);V,z,\theta),
  \qquad
  \frac{\dd X_\eta}{\dd u}
   =c_\eta'(u)\,a\!\left(X_\eta,c_\eta(u);\theta\right),
  \label{eq:time-warp-trajectory}
\end{equation}
and
\begin{equation}
  V_\eta^*(u,z)=V_0^*(c_\eta(u),z).
  \label{eq:time-warp-vstar}
\end{equation}
Since \(c_\eta\) is strictly increasing, the running minimum that defines a
first passage is unchanged after replacing the parent horizon by
\(s=c_\eta(u)\). Thus every live/frozen split, tangency, start-point
integral, and intrinsic omission threshold is inherited from the parent, with
no additional start-point integration or trialwise random variable.

If \(F_0(s)\), \(f_0(s)\), and \(S_0(s)=1-F_0(s)\) denote the (possibly
defective) parent CDF, density, and survivor, then
\begin{equation}
  F_\eta(u)=F_0(c_\eta(u)),
  \qquad
  S_\eta(u)=S_0(c_\eta(u)),
  \qquad
  f_\eta(u)=f_0(c_\eta(u))c_\eta'(u).
  \label{eq:time-warp-law}
\end{equation}
Equivalently,
\begin{equation}
  \log f_\eta(u)=\log f_0(c_\eta(u))
       +(\omega-1)\log(1+u).
  \label{eq:time-warp-log-density}
\end{equation}
As \(c_\eta(u)\to\infty\), intrinsic response probability is invariant:
\begin{equation}
  F_\eta(\infty)=F_0(\infty).
  \label{eq:time-warp-defect}
\end{equation}
For a raw deadline \(U_C\), put \(u_C=U_C-t_0\). If \(u_C>0\), its
omission probability is \(S_0(c_\eta(u_C))\); if \(U_C\le t_0\), it is one.
Parent internal quantiles and finish times transform as
\begin{equation}
  T_q=t_0+c_\eta^{-1}(s_q),
  \qquad
  T=t_0+c_\eta^{-1}(S),
  \label{eq:time-warp-inversion}
\end{equation}
where an infinite parent finish remains infinite. Any external time cutoff is
applied only after this conversion to physical time.

\subsection{Coverage of the ballistic family}
The following substitutions give the explicit warped trajectories for every
ballistic model family. The equations in the model sections
below are the parent equations; when \(\eta\ne0\), their argument \(u\) is
the internal time \(s=c_\eta(u_{\rm phys})\), and the physical density receives
the single Jacobian in \eqref{eq:time-warp-law}.

\begin{align}
 X_{\mathrm{LBA},\eta}(u)&=z+V c_\eta(u),
 &V_{\mathrm{LBA},\eta}^*(u,z)&=\frac{b-z}{c_\eta(u)},
 \label{eq:time-warp-lba}\\
 X_{\mathrm{BAwL},\eta}(u)&=z\e^{-kc_\eta(u)}
       +\frac{V}{k}\left[1-\e^{-kc_\eta(u)}\right],
 &V_{\mathrm{BAwL},\eta}^*(u,z)&=
       V_{\mathrm{BAwL}}^*(c_\eta(u),z),
 \label{eq:time-warp-bawl}\\
 X_{\mathrm{BAwD},\eta}(u)&=z+VQ_\rho(c_\eta(u))
       -\ell R_{\rho,\gamma}(c_\eta(u)),
 &V_{\mathrm{BAwD},\eta}^*(u,z)&=
       V_{\mathrm{BAwD}}^*(c_\eta(u),z),
 \label{eq:time-warp-bawd}\\
 X_{\mathrm{BAwDp},\eta}(u)&=z+Vq_\lambda(c_\eta(u)),
 &V_{\mathrm{BAwDp},\eta}^*(u,z)&=
       V_{\mathrm{BAwDp}}^*(c_\eta(u),z),
 \label{eq:time-warp-bawdp}\\
 X_{\mathrm{BAwF},\eta}(u)&=h_\rho(c_\eta(u))
       \left[z+Vc_\eta(u)\right],
 &V_{\mathrm{BAwF},\eta}^*(u,z)&=
       V_{\mathrm{BAwF}}^*(c_\eta(u),z),
 \label{eq:time-warp-bawf}\\
 X_{\mathrm{BAwR},\eta}(u)&=z+Vc_\eta(u)
       -\frac{\kappa c_\eta(u)^{p+1}}{p+1},
 &V_{\mathrm{BAwR},\eta}^*(u,z)&=
       V_{\mathrm{BAwR}}^*(c_\eta(u),z).
 \label{eq:time-warp-bawr}
\end{align}

For BTAwL the same clock is shared by both independent channels:
\begin{align}
 X_{j,\eta}(u)&=z_j\e^{-kc_\eta(u)}+V_jH_j(c_\eta(u)),
 &W_{j,\eta}(u,z)&=W_j(c_\eta(u),z),
 \qquad j\in\{T,S\}.
 \label{eq:time-warp-btawl}
\end{align}
Consequently, for the transient and sustained component models,
\begin{equation}
  F_{j,\eta}(u)=F_j(c_\eta(u)),
  \qquad
  f_{j,\eta}(u)=f_j(c_\eta(u))c_\eta'(u),
  \qquad j\in\{T,S\},
  \label{eq:time-warp-btawl-components}
\end{equation}
and the full local race obeys
\begin{equation}
  F_{\mathrm{BTAwL},\eta}(u)=F_{\mathrm{BTAwL}}(c_\eta(u)),
  \qquad
  f_{\mathrm{BTAwL},\eta}(u)=
       f_{\mathrm{BTAwL}}(c_\eta(u))c_\eta'(u).
  \label{eq:time-warp-btawl-race}
\end{equation}
The common Jacobian appears once, after the race has been formed.

More generally, if \(n\) independent parent accumulators have subdensity
\(\mathcal{L}_{i,0}(s)=f_{i,0}(s)\prod_{j\ne i}[1-F_{j,0}(s)]\), a shared warp
preserves the winner and gives
\begin{equation}
  \mathcal{L}_{i,\eta}(u)
   =\mathcal{L}_{i,0}(c_\eta(u))c_\eta'(u).
  \label{eq:time-warp-race}
\end{equation}
Thus the transform is applied once to the race likelihood, not once per
accumulator.

\subsection{Endpoints and critical launches}
Any finite parent endpoint \(s_{\max}\) maps to
\begin{equation}
  u_{\max,\eta}=c_\eta^{-1}(s_{\max})
    =(1+\omega s_{\max})^{1/\omega}-1.
  \label{eq:time-warp-endpoint}
\end{equation}
Thus the finite windows of BAwD (\(0\le\gamma<1\)), positive-launch BAwDp,
BAwF (\(\rho\in\{2,4,\infty\}\)), BAwR, and the pure transient BTAwL component are warped by
\eqref{eq:time-warp-endpoint}. Specifically, their parent endpoints are,
respectively,
\begin{equation}
  s_{\max}=T_{\rm sat}(0),\quad
  \frac{\log(1/\lambda)}{k},\quad
  \frac{x_{\max}}{k},\quad
  \left[\frac{(p+1)b}{\kappa p}\right]^{1/(p+1)},\quad
  \tau_t\psi(k\tau_t),
  \label{eq:time-warp-endpoints-list}
\end{equation}
where the corresponding parent definitions are given in the relevant
sections. LBA, BAwL, matched-decay BAwD, and sustained BTAwL have no finite
parent endpoint and remain unbounded after the warp. In every case the
critical launch \(V_{\rm crit}(z)\) and the intrinsic defect are unchanged;
only their physical-time locations are transformed.

For readability, the derivations that follow use the bare symbol \(u\) for
the parent model's accumulation coordinate. When the eta parameter is active,
replace that argument by \(s=c_\eta(u_{\rm phys})\) and use
\eqref{eq:time-warp-law} for the physical-time density; physical response
times and endpoints are then mapped back with \eqref{eq:time-warp-inversion}
or \eqref{eq:time-warp-endpoint}.

%======================================================================
\section{The Ballistic Idea}
%======================================================================
\subsection{Why There Is No Reflection Term}
\label{sec:ballistic-idea-recipe}
A diffusion accumulator has a new source of randomness at every instant, so its
first-passage law must be derived from the path measure, and the reflection
principle contributes an image term to the distribution function. A
\emph{ballistic} accumulator has no within-trial noise at all. Once $(V,z)$ are
drawn the trajectory $X(u)$ is a deterministic curve. All response-time
variability is inherited from two draws, and the likelihood is an exercise in
changing variables.

The consequence is a simple recipe, which every model below instantiates.

\begin{enumerate}
  \item \textbf{Required launch strength.} Find $V^{*}(u,z)$, the launch
        strength for which a trajectory starting at $z$ sits exactly at $b$ at
        time $u$.
  \item \textbf{Hit set on a sign-preserving branch.} If the coefficient of
        $V$ in the trajectory remains positive, the accumulator has crossed by
        time $u$ if it has crossed at any earlier time, so the event is
        $V \ge \min_{u' \le u} V^{*}(u',z)$. Conditional on $z$, the
        probability of that event is the launch survivor evaluated at that
        running minimum. If the coefficient changes sign, as in proportional
        clearance with an untruncated normal launch, split the hit set by sign
        intervals; the late negative-launch branch is derived in
        Section~\ref{sec:bawdp}.
  \item \textbf{Start-point average.} Integrate over $z \sim U(0,A)$. Because
        $V^{*}$ is affine or monotone in $z$, this is a one-dimensional
        integral of $\Gbar$ over an interval of launch strengths.
\end{enumerate}

Step 2 is where the models separate. On a sign-preserving branch, if
$V^{*}$ is decreasing in $u$ the running minimum is just $V^{*}(u,z)$ itself,
and waiting always helps. If $V^{*}$ is U-shaped in $u$ the running minimum
freezes once the trajectory has passed its peak, and waiting stops helping.
Sign-changing waveforms require the additional interval decomposition just
described.

\subsection{Two Antiderivatives}
Step 3 needs an antiderivative of the launch survivor. The normal launch uses
one primitive, while positive launches use the stop-loss transform.

For a normal launch, integrating the survivor is integrating $\PhiN$, and
\begin{equation}
  \Psi(x) = \int_{-\infty}^{x} \PhiN(w) \dd w = x\PhiN(x) + \phiN(x),
  \qquad \Psi'(x) = \PhiN(x).
  \label{eq:psi}
\end{equation}

For a positive launch law it is more natural to integrate the survivor from a
point upward, which gives the \emph{stop-loss transform}
\begin{equation}
  C(x) = \int_{x}^{\infty} \Gbar(w) \dd w = E\!\left[(V-x)_+\right],
  \qquad C'(x) = -\Gbar(x).
  \label{eq:stoploss}
\end{equation}
For an ordinary lognormal launch this is the Black--Scholes call price,
\begin{equation}
  C(x) =
  \begin{cases}
    \e^{\mu + \sigma^2/2}\,\PhiN\!\left(d_1(x)\right)
      - x\,\PhiN\!\left(d_2(x)\right), & x>0,\\[4pt]
    \e^{\mu + \sigma^2/2}-x, & x\le0,
  \end{cases}
  \label{eq:bs}
\end{equation}
and for a normal launch it is $C(x) = s_v\,\Psi\!\left((v-x)/s_v\right)$.
For the split-lognormal launch, use $C_{\rm sp}$ in
\eqref{eq:split-stoploss}; it is the same stop-loss object, with the two
piecewise Gaussian tails in \eqref{eq:split-tail-moments}. The normal,
lognormal, and split-lognormal forms are available in closed form. For a
Weibull launch, use $C_W$ from
\eqref{eq:weibull-tail-moments}; it is an incomplete-gamma expression. Whenever
a start-point average can be
written as $\int \Gbar(w)\dd w$ over an interval, \eqref{eq:stoploss} (or
$C_{\rm sp}$ or $C_W$) evaluates it as a difference of two values, and the model has a
closed-form distribution function. Section~\ref{sec:closure} identifies the
one place where power or logarithmic measures appear and shows how the
ordinary and split-lognormal launches close through Gaussian tails while the
Weibull launch closes through incomplete gamma functions.

%======================================================================
\section{The Linear Ballistic Accumulator}
%======================================================================
When the warp is active, the trajectory and likelihood in this section are
evaluated at \(s=c_\eta(u)\) and the density is multiplied by
\(c_\eta'(u)\), as in \eqref{eq:time-warp-lba} and
\eqref{eq:time-warp-law}.
\subsection{Trajectory and Crossing Time}
The LBA takes a constant drive and no decay,
\begin{equation}
  \dot X = V, \qquad X(0) = z, \qquad X(u) = z + Vu.
  \label{eq:lba-traj}
\end{equation}
Setting $X(u) = b$ gives the required launch strength
\begin{equation}
  V^{*}(u,z) = \frac{b-z}{u},
  \qquad
  -\partial_u V^{*} = \frac{b-z}{u^2},
  \label{eq:lba-vstar}
\end{equation}
which is decreasing in $u$ and affine and decreasing in $z$. Both facts are
what make the LBA the easiest member of the family: the hit set never freezes,
and the start-point average never leaves the Gaussian primitives.

Equivalently, and more revealingly, the crossing time is available in closed
form,
\begin{equation}
  u = \frac{b-z}{V}, \qquad V > 0.
  \label{eq:lba-crossing}
\end{equation}
The response-time law is therefore the pushforward of the joint law of $(V,z)$
through \eqref{eq:lba-crossing}. Nothing else happens in the model, and this
observation drives the reductions in Section~\ref{sec:reductions}.

\subsection{Point Start Point}
With $A = 0$ the start point is degenerate at $z = 0$ and $u = b/V$ exactly. The
distribution function and density are the pushforward of the launch law under
$V \mapsto b/V$:
\begin{equation}
  F_{\mathrm{LBA}}(u) = \Gbar\!\left(\frac{b}{u}\right),
  \qquad
  f_{\mathrm{LBA}}(u) = g\!\left(\frac{b}{u}\right)\frac{b}{u^2}.
  \label{eq:lba-point}
\end{equation}
For a normal launch, with $Z_+ = \PhiN(v/s_v)$ under truncation and $Z_+ = 1$
without,
\begin{equation}
  F_{\mathrm{LBA}}(u) = \frac{1}{Z_+}\PhiN\!\left(\frac{v - b/u}{s_v}\right),
  \qquad
  f_{\mathrm{LBA}}(u) = \frac{1}{Z_+}\frac{b}{u^2 s_v}
     \phiN\!\left(\frac{b/u - v}{s_v}\right).
\end{equation}

\subsection{Start-Point Variability}
\label{sec:lba-startpoint}
For $A > 0$, write the standardized threshold offset and slope
\begin{equation}
  c(u) = \frac{v - b/u}{s_v},
  \qquad
  m(u) = \frac{1}{u\,s_v},
  \qquad
  \mathrm{span}(u) = m(u)\,A,
  \label{eq:lba-cm}
\end{equation}
so that the conditional hit probability is $P(V \ge V^{*}) = \PhiN(c + mz)$,
affine inside the probit. Averaging over $z$ and using \eqref{eq:psi},
\begin{equation}
  F_{\mathrm{LBA}}(u)
  = \frac{1}{Z_+}\cdot\frac{1}{A}\int_0^A \PhiN(c + mz)\dd z
  = \frac{1}{Z_+}\cdot\frac{\Psi(c+\mathrm{span}) - \Psi(c)}{\mathrm{span}}.
  \label{eq:lba-cdf}
\end{equation}
Differentiating, the $u$-derivative of $\PhiN(c+mz)$ is
$\phiN(c+mz)\,(b-z)/(u^2 s_v)$, and the substitution $w = c + mz$ turns the
integrand into $\phiN(w)(bm + c - w)$. Since $\int -w\phiN(w)\dd w = \phiN(w)$,
\begin{equation}
  f_{\mathrm{LBA}}(u)
  = \frac{1}{Z_+}\cdot\frac{1}{A}
    \left[v\,\Delta\PhiN + s_v\,\Delta\phiN\right],
  \label{eq:lba-pdf}
\end{equation}
where $\Delta\PhiN = \PhiN(c+\mathrm{span}) - \PhiN(c)$ and
$\Delta\phiN = \phiN(c+\mathrm{span}) - \phiN(c)$, after using
$bm + c = v/s_v$. Equations \eqref{eq:lba-cdf} and \eqref{eq:lba-pdf} are the
standard LBA formulas in a compact parameterization; expanding
$c$ and $\mathrm{span}$ recovers the familiar expressions in
$(b-A-uv)/(u s_v)$ and $(b-uv)/(u s_v)$.

\subsection{Positive Launches}
Replacing the normal launch by a positive launch changes
nothing structural, because \eqref{eq:lba-vstar} is still affine in $z$.
For the ordinary lognormal, substituting
$w = V^{*}(u,z) = (b-z)/u$, so that $\dd z = -u \dd w$,
\begin{equation}
  F(u) = \frac{u}{A}\left[C\!\left(\tfrac{b-A}{u}\right) - C\!\left(\tfrac{b}{u}\right)\right],
\end{equation}
with $C$ the Black--Scholes form \eqref{eq:bs}. For a split-lognormal launch
the identical formula holds with $C$ replaced by $C_{\rm sp}$ from
\eqref{eq:split-stoploss}; its density uses $M_1$ and $M_0$ from
\eqref{eq:split-tail-moments}. Thus no start-point quadrature is introduced by
the asymmetry. For a Weibull launch, the same expression holds with $C_W$:
\begin{equation}
  F_{\mathrm{LBA},W}(u)
   = \frac{u}{A}\left[C_W\!\left(\frac{b-A}{u}\right)-C_W\!\left(\frac{b}{u}\right)\right],
\end{equation}
and, writing $w_{\rm lo}=(b-A)/u$, $w_{\rm hi}=b/u$,
\begin{equation}
  f_{\mathrm{LBA},W}(u)
   = \frac{M_{1,W}(w_{\rm lo})-M_{1,W}(w_{\rm hi})}{A}.
\end{equation}
The Weibull stop-loss and partial moments are closed in incomplete gamma
functions, so no start-point quadrature is introduced there either.

The point-start case deserves separate attention. With $A = 0$,
$u = b/V$, so
\begin{equation}
  \log u = \log b - \log V \sim \Normal\!\left(\log b - \mu,\ \sigma^2\right),
\end{equation}
and the response time is exactly $t_0$ plus a lognormal variate. This is the
lognormal race model, and it is the cleanest instance of the reduction
discussed in Section~\ref{sec:reductions}.

For the split-lognormal launch, the same point-start transformation gives
$\log u=\log b-Y$. It is split normal with median $\log b-\mu$, split point
$\log b-c$, left and right widths $s_R$ and $s_L$, respectively, and width
ratio parameter $-\delta$. Consequently the point-start response is a shifted
split-lognormal law, reducing to the shifted ordinary lognormal when
$\delta=0$.

For a Weibull launch, the point-start response is the reciprocal-Weibull
pushforward
\begin{equation}
  F_{\mathrm{LBA},W}(u)=\exp\!\left[-\left(\frac{b}{u\theta_W}\right)^r\right],
  \qquad
  f_{\mathrm{LBA},W}(u)=g_W(b/u)\frac{b}{u^2},
\end{equation}
which is closed but is not itself a Weibull law on the response-time scale.

\subsection{The Race}
Nothing in this note requires more than an independent race. With $n$
accumulators, response $i$ at time $u$ has defective subdensity
\begin{equation}
  \mathcal{L}_i(u) = f_i(u)\prod_{j \ne i}\left[1 - F_j(u)\right],
  \label{eq:race}
\end{equation}
and the probability that no accumulator ever responds is
\begin{equation}
  P(\text{omission}) = \prod_{j=1}^{n}\left[1 - F_j(\infty)\right].
  \label{eq:race-omission}
\end{equation}
The accumulators are coupled only through the winner-take-all reading of
\eqref{eq:race}, which is a decision rule applied to independent trajectories,
not an interaction in the dynamics. Section~\ref{sec:lca} returns to that
distinction, and Section~\ref{sec:defect} evaluates
\eqref{eq:race-omission} for each model.

%======================================================================
\section{From the Leaky Competing Accumulator to the LBA and Back}
\label{sec:lca}
%======================================================================
\subsection{The Full Generator}
The leaky competing accumulator posits, for accumulator $i$ among $n$,
\begin{equation}
  \dd X_i = \left[\rho_i - k X_i - \beta\sum_{j\ne i} X_j\right]\dd u
            + \varsigma \dd W_i,
  \qquad
  X_i \ge 0,
  \label{eq:lca}
\end{equation}
with input $\rho_i$, leak rate $k$, lateral inhibition $\beta$, within-trial
noise of scale $\varsigma$, and a rectification that forbids negative
activation. Four mechanisms are bundled together here, and the ballistic family
is obtained by switching three of them off.

\begin{itemize}
  \item Set $\varsigma = 0$ and let the input be a single trialwise draw
        $\rho_i = V_i$. This is the ballistic simplification: within-trial noise
        is replaced by between-trial variability in the drive, and the
        trajectory becomes deterministic given the draw. Rectification then
        becomes inert for any trajectory that starts at $z \ge 0$ and has
        positive drive, so it can be dropped rather than assumed away.
  \item Set $\beta = 0$. This removes lateral inhibition, so the generator
        becomes diagonal and the accumulators are dynamically independent.
  \item Set $k = 0$. This removes leak, so the accumulator becomes a perfect
        integrator of a constant input.
\end{itemize}

Applying all three gives $\dot X_i = V_i$: the LBA of
\eqref{eq:lba-traj}. Applying the first two but retaining $k > 0$ gives
\begin{equation}
  \dot X_i = V_i - k X_i,
  \label{eq:bawl-gen}
\end{equation}
which is the ballistic accumulator with leak of Section~\ref{sec:bawl}. In the
historical order the leaky system came first and the LBA was introduced as the
tractable simplification of it, so BAwL is better described as a partial
restoration than as an extension: it puts the leak term back and leaves the
other three simplifications in place.

\subsection{What Omitting Inhibition Buys, and What It Costs}
We do not include $\beta$, and the omission is deliberate rather than
incidental.

The benefit is exactly the factorization in \eqref{eq:race}. With $\beta = 0$
the drift field is diagonal, each $X_i$ depends only on its own $(V_i, z_i)$,
and the joint law of the finishing times is the product of the marginals.
Every closed form in this note rests on that. With $\beta \ne 0$ the states are
coupled, $V^{*}$ for accumulator $i$ depends on the entire trajectory of the
others, and there is no change of variables from a single launch strength to a
crossing time. The corresponding response-time law then requires a
multidimensional first-passage calculation, and the winner's density is no
longer a one-dimensional object.

The cost is that lateral inhibition is omitted. With $\beta\ne0$, losing
accumulators feed back into the winner, the state equations are coupled, and
the one-dimensional change-of-variables derivations below no longer apply.
The leak term alone, without that coupling, still changes the response-time
law while preserving analytic tractability.

The models below do retain a form of competition, but only at the level of
the decision rule. The first accumulator
to reach $b$ determines the response, so accumulators compete for the response
in the sense of \eqref{eq:race}. They do not interact in the sense of
\eqref{eq:lca}. In the rectified leaky competing accumulator the two are
confounded; here they are separated, and the separation is the point.

\subsection{Where the Decay Acts}
\label{sec:decay-targets}
Having restored leak, there are several further decay mechanisms that are not
in \eqref{eq:lca} at all and that are easy to conflate with it and with one
another. What separates them is not the \emph{shape} of the decay, since any
of them can be given an exponential or a power-law profile. It is the object
the decay attacks. There are four such objects: the accumulated state, the
input, the readout, and nothing at all.

\paragraph{The state (leak).} Compare
\begin{align}
 \text{leak:}\quad &\dot X = V-kX,
 &\text{drive decay:}\quad
 \dot X = Vh_\rho(u)-\ell h_\rho(u)^\gamma, \nonumber\\
 & h_\rho(u)=
 \begin{cases}
 (1+ku/\rho)^{-\rho},&\rho<\infty,\\
 \e^{-ku},&\rho=\infty,
 \end{cases}
 \label{eq:two-generators}
\end{align}
The left generator is autonomous: the decay term is proportional to the state,
so the accumulator forgets evidence at a rate set by how much it currently
holds, and the start point is forgotten along with everything else.

\paragraph{The input (drive decay).} The right generator of
\eqref{eq:two-generators} is a perfect integrator driven by a transient input:
nothing already accumulated is forgotten, the start point persists
undiminished, and the clearance fades with the same base kernel raised to the
fixed exponent $\gamma$. The indices $(\rho,\gamma)$ therefore distinguish
temporal decay shape from clearance tracking. Sections~\ref{sec:bawl} and
\ref{sec:bawd} derive the two families, and Section~\ref{sec:clocks} shows how
they differ geometrically.

The right generator admits one further distinction, which is the subject of
Section~\ref{sec:bawdp}. In the BAwD family the amplitude $\ell$ is fixed
across trials, while $\gamma$ prescribes how its clearance waveform tracks the
base kernel. A different construction ties the clearance amplitude to the
trialwise launch strength, $\ell=\lambda V$; with the exponential base this
gives
\begin{equation}
  \text{proportional clearance:}\quad
  \dot X = V\left(\e^{-ku} - \lambda\right)
         = V h(u),
  \qquad
  h(u) = \e^{-ku} - \lambda,
  \label{eq:bawdp-gen}
\end{equation}
which is better read not as a clearance claim at all but as the statement that
every trial integrates the \emph{same} transient waveform $h$, with $V$ setting
only its amplitude. That single change removes the algebraic difficulty of the
drive-decay model without removing either of its two distinctive predictions,
the hard right endpoint and the intrinsic omissions.

\paragraph{The readout (global fading).} A third possibility, the subject of
Section~\ref{sec:bawf}, is that nothing in the accumulation is lost at all and
what fades is how strongly the accumulated quantity remains \emph{available} to
the decision mechanism:
\begin{equation}
  \text{global fading:}\quad
  X(u) = h_\rho(u)\left[z + Vu\right],
  \label{eq:bawf-gen}
\end{equation}
with the same kernel $h_\rho$ as in \eqref{eq:two-generators}. The latent trace
$z+Vu$ is an ordinary LBA; $h_\rho$ is a global gain applied to it. Written as
a generator this is $\dot X = Vh_\rho(u) - kX$ in the exponential case, which
looks like drive decay \emph{and} leak acting together but has only one
primitive rate, because the clearance term $kX$ takes its evidence scale from
the state instead of carrying an independent one. That absence is the point of
the model: the drive-decay pair $(k,\ell)$ sets the response window and the
omission boundary jointly, whereas here the window is fixed by $k$ alone.

\paragraph{Nothing (elapsed-time clearance).} The last possibility is that the
drain consults neither the state nor the drive but only the clock:
\begin{equation}
  \text{ramping clearance:}\quad
  \dot X = V - \kappa u^{p},
  \qquad \kappa > 0,\ p > 0,
  \label{eq:bawr-gen}
\end{equation}
the subject of Section~\ref{sec:bawr}. The drive never fades, the state is
never forgotten, and the opposing term \emph{grows} without bound instead of
decaying. That is why it belongs to none of the families above: in all of them
the clearance waveform $h_\rho^\gamma$ is bounded by one and non-increasing.
Here the opposing term is exogenous to the accumulator and grows with elapsed
time. It is the simplest way to obtain a hard
response window, since $-\partial_u V^{*}$ changes sign for the elementary
reason that a constant is eventually beaten by a growing power.

\paragraph{Both the state and the input.} Nothing forces the targets to be
exclusive.  The BTAwL family uses two local subraces, both with state leak,
and gives them different stimulus profiles:
\begin{equation}
  \begin{aligned}
  \text{transient:}\quad&
    \dot X_T=V_T\eta_T(u)-kX_T,
    &\eta_T(u)&=\frac{u}{\tau_t}\,\e^{-u/\tau_t},\\
  \text{sustained:}\quad&
    \dot X_S=V_S\eta_S(u)-kX_S,
    &\eta_S(u)&=1-\e^{-u/\tau_s}.
  \end{aligned}
  \label{eq:btawl-intro}
\end{equation}
The full model races these independent channels locally; it does not
replace them by a shared trajectory.  The transient profile is the filtered
onset response, while the sustained profile is the filtered step response.

\paragraph{Summary.} Each of these targets produces a different answer to
two structural questions: whether the start point
survives, and whether the response criterion $b$ can move the response window.
Section~\ref{sec:glance} tabulates the answers, and Section~\ref{sec:clocks}
proves them by rewriting every model as an LBA against a moving boundary.

%======================================================================
\section{The Six Variants at a Glance}
\label{sec:glance}
%======================================================================
This section states what each variant \emph{is}, before any of them is derived,
so that the derivations can be read as answers to a fixed list of questions.
Everything asserted here is proved later: the trajectories in
Sections~\ref{sec:bawl}--\ref{sec:btawl}, the clock and boundary equivalences in
Section~\ref{sec:clocks}, the omission thresholds in Section~\ref{sec:defect},
and the closure claims in Section~\ref{sec:closure}.

Throughout, $h_\rho(u)=(1+ku/\rho)^{-\rho}$ with $h_\infty(u)=\e^{-ku}$ is the
decay kernel, $H_\rho(x)=(1+x/\rho)^{\rho}$ is its reciprocal in the
dimensionless phase $x=ku$, and $Q_\rho,R_{\rho,\gamma},q_\lambda$ are the
accumulated clocks defined where they are first needed. The drive-decay family
is shown for $\gamma<1$; its matched corner $\gamma=1$ behaves like BAwL in
having no endpoint, and is treated separately.

\begin{center}
\small
\begin{tabular}{L{0.075\textwidth}L{0.235\textwidth}L{0.115\textwidth}L{0.44\textwidth}}
\toprule
Model & Trajectory or generator & Decay attacks & What it says \\
\midrule
LBA &
$X=z+Vu$ &
--- &
A perfect integrator of a constant trialwise input. No decay of any kind, so
the only randomness is $(V,z)$ and the only dynamics is a straight line. \\
\addlinespace
BAwL &
$\dot X = V-kX$ &
the state &
Evidence already held is forgotten at a rate set by how much is held. The
start point is forgotten with it. Evidence saturates at the plateau $V/k$, so
a trial responds if and only if $V>kb$, whatever its start point. \\
\addlinespace
BAwD &
$\dot X = Vh_\rho(u)-\ell h_\rho(u)^{\gamma}$ &
the input &
The stimulus drive is transient and a separately scaled clearance $\ell$
drains the accumulator while it lasts. Nothing accumulated is forgotten. The
two rates $(k,\ell)$ are independent, which lets the response window and the
omission boundary be set separately, and also entangles both of them with
$b$. \\
\addlinespace
BAwDp &
$\dot X = V\left(\e^{-ku}-\lambda\right)$ &
the input, with clearance tied to it &
Every trial integrates the \emph{same} transient waveform and $V$ sets only
its amplitude. Read as a clearance model, $\ell=\lambda V$; read
geometrically, the accumulation clock simply runs out at a time common to all
trials. \\
\addlinespace
BAwF &
$X=h_\rho(u)\left[z+Vu\right]$,
i.e.\ $\dot X = Vh_\rho(u)-kX$ at $\rho=\infty$ &
the readout &
Nothing is lost from the accumulation; what fades is how strongly the
accumulated trace remains available to the decision mechanism. The implied
clearance is $kX$, so unlike BAwD there is \emph{no independent evidence
scale}: one rate $k$ fixes the response window and $b$ cannot move it. \\
\addlinespace
BAwR &
$\dot X = V-\kappa u^{p}$ &
nothing &
An exogenous drain that consults neither the state nor the drive, only the
clock, and that \emph{grows} rather than decays. The opposing term is imposed
externally to the accumulator. Alone among these models the opposing term
is unbounded. \\
\addlinespace
BTAwL &
$\dot X_j=V_j\eta_j(u)-kX_j$, $j\in\{T,S\}$, raced locally &
the state and the input &
Two independent subraces implement the transient pulse and sustained step
profiles of Smith's detector.  The allocation $\pi$ scales their independent
launch laws; the pure transient and pure sustained components are the endpoint
submodels. \\
\bottomrule
\end{tabular}
\end{center}

The consequences follow. ``Start point survives'' asks whether $z$ still
contributes undiminished at late $u$; ``$b$ moves $T_{\max}$'' asks whether the
boundary level relocates the hard response window. These two structural
questions distinguish the models.

\begin{center}
\small
\begin{tabular}{L{0.07\textwidth}L{0.175\textwidth}L{0.05\textwidth}L{0.195\textwidth}L{0.05\textwidth}L{0.135\textwidth}L{0.05\textwidth}}
\toprule
Model & Equivalent (clock, boundary) & Start pt.\ lives & Endpoint
$T_{\max}$ & $b$ moves it & Omission threshold
$V_{\mathrm{crit}}(z)$ & Decay pars \\
\midrule
LBA & $u$; flat at $b$ & yes & none & --- & $0$ & --- \\
\addlinespace
BAwL & $r=(\e^{ku}-1)/k$ \emph{and} a state rescaling; straight, slope $kb$ &
no & none; $\e^{-ku}$ tail & --- & $kb$, free of $z$ & $k$ \\
\addlinespace
BAwD & $Q_\rho(u)$; $b+\ell R_{\rho,\gamma}$, diverging &
yes & $T_{\rm sat}(0)$, solving
$\Theta_{\rho,\gamma}(kT)=kb/\ell$ &
yes & $V_c(z)$, eq.~\eqref{eq:bawd-crit} & $k,\ell$ \\
\addlinespace
BAwDp & $q_\lambda(u)$, halting; flat at $b$, unchanged &
yes & $u_{*}=\log(1/\lambda)/k$ & no & $(b-z)/q_\lambda^{\max}$, affine in $z$
& $k,\lambda$ \\
\addlinespace
BAwF & $u$, unwarped; $b\,H_\rho(ku)$, multiplicative &
no & $x_{\max}/k$, with $x_{\max}=1$ at $\rho=\infty$ and
$\rho/(\rho-1)$ otherwise &
no & $kbH_\rho'(x_{\rm sat}(z))$, so $V_c(0)=\e kb$ at $\rho=\infty$ & $k$ \\
\addlinespace
BAwR & $u$, unwarped; $b+\kappa u^{p+1}/(p+1)$, additive &
yes & $\bigl[(p+1)b/(\kappa p)\bigr]^{\frac{1}{p+1}}$ & yes &
$\kappa\,u_{\rm sat}(z)^{p}$ & $\kappa,p$ \\
\addlinespace
BTAwL transient & $r=\int_0^u\e^{k\theta}\eta_T$ \emph{and} a state rescaling;
slope $kb/\eta_T(u)$ &
no & $\tau_t\psi(k\tau_t)$, solving $H_T'(T)=0$ &
no & $kb/\eta_T(T_{\rm sat}(z))$, never below $\e kb$ & $k,\tau_t$ \\
BTAwL sustained & $r=\int_0^u\e^{k\theta}\eta_S$ \emph{and a state rescaling};
$b\e^{ku}$, asymptotic slope $kb$ &
no & none & --- & $kb$ (for $k>0$) & $k,\tau_s$ \\
\bottomrule
\end{tabular}
\end{center}

Three things follow from the table.

Three models attack the start point, and they do it for different reasons.
BAwL forgets its start point along with the accumulated state; BAwF stops
reporting it through the readout; and both BTAwL subraces have the same state
leak.  The transient subrace has a U-shaped required-launch curve and a hard
endpoint, whereas the sustained subrace approaches the leak plateau and has
no finite endpoint.  For $k>0$, the full local race therefore has a hard
endpoint only at $\pi=0$.

BAwF and BAwR leave the clock alone and move the boundary, which makes them the
two ``moving boundary in ordinary time'' members. They differ in exactly one
respect. BAwF's boundary moves multiplicatively ($b\mapsto bH_\rho$), which
attenuates the start point; BAwR's moves additively
($b\mapsto b+\kappa u^{p+1}/(p+1)$), which leaves it alone. That single
contrast accounts for everything else that separates them.

The hard endpoints belong to BAwD, BAwDp, BAwF, BAwR, and the pure transient
BTAwL component.  For $k>0$, the full local race loses the transient endpoint as
soon as $\pi>0$ because the sustained subrace can finish later.  For the transient
component the window is $\tau_t\psi(k\tau_t)$ and is independent of $b$; for the
sustained component and the interior local race there is no finite window.

\subsection{A Note on the Names}
\label{sec:names}
The abbreviations are not as orthogonal as the models are, and two of them
deserve a caveat.

First, ``decay'' and ``fading'' are near-synonyms in ordinary use but are used
here for different mechanisms: \emph{decay} attacks the input (BAwD) and
\emph{fading} attacks the readout (BAwF). The distinction is real: one has an
independent clearance scale and the other does not. It is carried entirely by
convention, though, since neither word contains it.

Second, the model in \eqref{eq:bawr-gen} was originally called BAwP, for the
polynomial profile $\kappa u^{p}$. That name had two problems. Its ``P'' and
the ``p'' of BAwDp were unrelated, one abbreviating a polynomial and the other
proportional clearance, and ``polynomial'' named a functional form where every
other abbreviation in the family names a mechanism. The form is also not what
makes the model distinctive: an exponentially growing drain belongs to the same
model class, and a polynomially decaying \emph{drive} already belongs to BAwD
as its finite-$\rho$ member. What is distinctive is that the clearance rate
\emph{ramps up} with elapsed time, so the model is BAwR, and ``p'' is left free
for proportional clearance. The exponent is still called $p$; it is now the
only $p$ in the family that means what it looks like.

Third, BTAwL names the transient/sustained local race.  The ``T'' identifies
the transient subrace and the ``wL'' says that both subraces use ordinary state
leak; the sustained partner is part of the full local race, while the
transient and sustained components expose its endpoint submodels.

%======================================================================
\section{Ballistic Accumulator with Leak}
\label{sec:bawl}
%======================================================================
The clock-free, uncorrelated BAwL applies the common warp through
\eqref{eq:time-warp-bawl}; its critical launch and intrinsic defect are
unchanged.
\subsection{Trajectory and Required Launch Strength}
Solving \eqref{eq:bawl-gen} with $X(0)=z$,
\begin{equation}
  X(u) = z\,E + \frac{V}{k}\left(1 - E\right),
  \qquad
  E = \e^{-ku},
  \qquad
  G = 1 - E.
  \label{eq:bawl-traj}
\end{equation}
The trajectory is monotone in $u$ and approaches the asymptote $V/k$, so the
first crossing of $b$ happens if and only if $X(u) \ge b$ at some point, and
\begin{equation}
  V^{*}(u,z) = \frac{k(b - zE)}{G} = C_1(u) - C_2(u)\,z,
  \qquad
  C_1 = \frac{kb}{G},
  \qquad
  C_2 = \frac{kE}{G}.
  \label{eq:bawl-vstar}
\end{equation}
Both coefficients are positive and decreasing in $u$, and $V^{*}$ is affine and
decreasing in $z$. As in the LBA, the hit set never freezes, and the
start-point average stays inside the primitives of Section~\ref{sec:closure}.
Using $E' = -kE$ and $G' = kE$,
\begin{equation}
  C_1' = -bJ, \qquad C_2' = -J,
  \qquad
  J(u) = \frac{k^2 E}{G^2},
  \qquad
  -\partial_u V^{*}(u,z) = J\,(b-z) > 0.
  \label{eq:bawl-jac}
\end{equation}
The $k \to 0$ limit is exact and elementary: $C_1 \to b/u$, $C_2 \to 1/u$,
$J \to 1/u^2$, recovering \eqref{eq:lba-vstar} term by term.

\subsection{Normal Launch}
With $c(u) = (v - C_1)/s_v$, $m(u) = C_2/s_v$, and
$\mathrm{span} = mA$, the conditional hit probability is again
$\PhiN(c + mz)$, so \eqref{eq:psi} gives
\begin{equation}
  F_{\mathrm{BAwL}}(u)
  = \frac{1}{Z_+}\cdot\frac{\Psi(c+\mathrm{span}) - \Psi(c)}{\mathrm{span}},
  \label{eq:bawl-cdf}
\end{equation}
identical in form to \eqref{eq:lba-cdf} with the LBA coefficients replaced by
the leaky ones. Differentiating through \eqref{eq:bawl-jac},
\begin{equation}
  f_{\mathrm{BAwL}}(u)
  = \frac{1}{Z_+}\cdot\frac{J}{A s_v}\int_0^A \phiN(c+mz)(b-z)\dd z
  = \frac{1}{Z_+}\cdot
    \frac{J\left[(bm+c)\Delta\PhiN + \Delta\phiN\right]}{A\,s_v\,m^2}.
  \label{eq:bawl-pdf}
\end{equation}
At $A = 0$ this collapses to
$f_{\mathrm{BAwL}}(u) = Z_+^{-1}\,C_1C_2\,\phiN(c)/s_v$ with $C_1C_2 = bJ$, and
at $k=0$ every expression reduces identically to \eqref{eq:lba-cdf} and
\eqref{eq:lba-pdf}, since $bm + c \to v/s_v$ and $J/m^2 \to s_v^2$.

The bracket in \eqref{eq:bawl-pdf} is a difference of comparable terms in the
tails, because $\Delta\phiN$ is signed while $(bm+c)\Delta\PhiN$ is not. That is
a conditioning property of the expression rather than a property of the model,
but it matters before the formula is evaluated far from the mode.

\subsection{Positive Launches}
\label{sec:bawl-logn}
Because $V^{*}$ is affine in $z$, all positive launch families are equally
closed. Put
$w = V^{*}(u,z)$, so $\dd z = -\dd w / C_2$, and let
\begin{equation}
  w_{\mathrm{hi}} = V^{*}(u,0) = C_1,
  \qquad
  w_{\mathrm{lo}} = V^{*}(u,A) = C_1 - C_2 A .
\end{equation}
Then
\begin{equation}
  F_{\mathrm{BAwL}}(u)
  = \frac{1}{A}\int_0^A \Gbar\!\left(V^{*}(u,z)\right)\dd z
  = \frac{1}{A\,C_2}\left[C(w_{\mathrm{lo}}) - C(w_{\mathrm{hi}})\right],
  \label{eq:bawl-cdf-logn}
\end{equation}
with $C$ the Black--Scholes form \eqref{eq:bs} for the ordinary lognormal.
For a split-lognormal launch, replace $\Gbar$ and $C$ by
$\Gbar_{\rm sp}$ and $C_{\rm sp}$ from
\eqref{eq:split-survivor-moment}--\eqref{eq:split-stoploss}.
For a Weibull launch, replace them by $\Gbar_W$ and $C_W$ from
\eqref{eq:weibull-launch}--\eqref{eq:weibull-tail-moments}. For the density,
note from
\eqref{eq:bawl-vstar} that $b - z = (w - kb)/C_2$, and from \eqref{eq:bawl-jac}
that $J/C_2^2 = 1/E$, so
\begin{equation}
  f_{\mathrm{BAwL}}(u)
  = \frac{1}{A\,E}\int_{w_{\mathrm{lo}}}^{w_{\mathrm{hi}}} (w - kb)\,g(w)\dd w
  = \frac{1}{A\,E}\Big\{
      \e^{\mu+\sigma^2/2}\,\Delta\PhiN_{d_1}
      - kb\,\Delta\PhiN_{d_2}
    \Big\},
  \label{eq:bawl-pdf-logn}
\end{equation}
where $\Delta\PhiN_{d_j} = \PhiN(d_j(w_{\mathrm{lo}})) - \PhiN(d_j(w_{\mathrm{hi}}))$.
For the split-lognormal launch this becomes, with
$\Delta M_r=M_r(w_{\mathrm{lo}})-M_r(w_{\mathrm{hi}})$,
\begin{equation}
  f_{\mathrm{BAwL,sp}}(u)
  =\frac{\Delta M_1-kb\,\Delta M_0}{A\,E}.
  \label{eq:bawl-pdf-split}
\end{equation}
For a Weibull launch, with
$\Delta M_{q,W}=M_{q,W}(w_{\mathrm{lo}})-M_{q,W}(w_{\mathrm{hi}})$,
\begin{equation}
  f_{\mathrm{BAwL},W}(u)
   =\frac{\Delta M_{1,W}-kb\,\Delta M_{0,W}}{A\,E},
  \label{eq:bawl-pdf-weibull}
\end{equation}
The point-start members are $F = \Gbar(C_1)$ and $f = b\,J\,g(C_1)$, with
$(\Gbar,g)$ replaced by $(\Gbar_{\rm sp},g_{\rm sp})$ or
$(\Gbar_W,g_W)$ for the split-lognormal or Weibull launch, respectively.

The critical launch strength $kb$ appears in \eqref{eq:bawl-pdf-logn} as the
subtracted constant, which is the algebraic trace of the fact established next:
$kb$, and not anything involving $z$, decides whether a trial ever responds.

\subsection{Defectiveness and the Right Tail}
As $u \to \infty$ we have $E \to 0$, $G \to 1$, hence $C_1 \to kb$ and
$C_2 \to 0$. The start-point span vanishes, and
\begin{equation}
  F_{\mathrm{BAwL}}(\infty) = \frac{\Gbar(kb)}{Z_+} = \frac{P(V > kb)}{Z_+}.
  \label{eq:bawl-mass}
\end{equation}
The eventual-hit condition is that the leaky asymptote clears the boundary,
$V/k > b$, and it does not involve $z$ at all. The start point leaks away, so
start-point variability changes the timing of responses but not which trials
produce one.

The approach to that plateau is exponential rather than power-law. Since
$J \sim k^2\e^{-ku}$ and $m \to 0$,
\begin{equation}
  f_{\mathrm{BAwL}}(u) \;\sim\;
  k^2\left(b - \tfrac{1}{2}A\right) g(kb)\,\e^{-ku},
  \qquad u \to \infty,
  \label{eq:bawl-tail}
\end{equation}
for any of the four launch laws, with
$g(kb) = \phiN((kb-v)/s_v)/(s_v Z_+)$ in the normal case and
$g(kb)=g_{\rm sp}(kb)$ for the split launch and $g(kb)=g_W(kb)$ for the
Weibull launch. The right tail decays at rate
exactly $k$, with an amplitude set by the launch density \emph{at} the critical
value. Launch strengths just above $kb$
cross arbitrarily late, which is why the model has no hard right endpoint: the
support is unbounded even though the hit probability is bounded away from one.

The tail rate together with the plateau \eqref{eq:bawl-mass} determines $k$ and
$kb$; because $b$ enters the plateau only through the product $kb$, the plateau
alone cannot separate the two parameters.

\subsection{Correlated Launch Strengths}
If the accumulators in a race are to share trial-level variability, the natural
place to put it is the launch strengths rather than the finishing times. A
one-factor decomposition
\begin{equation}
  V_i = v_i + s_{v,i}\left(\sqrt{\rho_i}\,\xi + \sqrt{1-\rho_i}\,\varepsilon_i\right),
  \qquad
  \xi, \varepsilon_i \stackrel{iid}{\sim} \Normal(0,1),
\end{equation}
makes the racers conditionally independent given $\xi$, with shifted means. The
race likelihood is then a one-dimensional Gauss--Hermite integral over $\xi$ of
a product of the closed forms above. This remains a between-trial drift
correlation and not a dependence between finishing times; the induced rank
correlation of the crossing times is a derived quantity, and under positive
truncation $\rho$ refers to the untruncated draws. It is also not lateral
inhibition: conditional on $\xi$ the accumulators still do not interact.

%======================================================================
\section{Ballistic Accumulator with Drive Decay}
\label{sec:bawd}
%======================================================================
BAwD uses the same outer time transform, so every live and frozen expression
below is evaluated at \(s=c_\eta(u)\) and its physical density uses the single
Jacobian in \eqref{eq:time-warp-law}, as summarized in
\eqref{eq:time-warp-bawd}.
\subsection{The base kernel, trajectory, and required launch}
The drive-decay model is a fixed-index family. Let
\[
  \rho\in\left\{1,2,4,\infty\right\},\qquad
  \gamma\in\left\{0,\tfrac12,\tfrac23,\tfrac34,1\right\},
\]
where $\rho$ controls the temporal shape of the fading drive and $\gamma$
controls how closely the clearance follows that drive. Neither index is a
trialwise random variable. Define
\begin{equation}
 h_\rho(u) =
 \begin{cases}
   (1+ku/\rho)^{-\rho}, & \rho<\infty,\\
   \e^{-ku}, & \rho=\infty,
 \end{cases}
 \qquad
 \dot X(u)=Vh_\rho(u)-\ell h_\rho(u)^\gamma,
 \qquad X(0)=z.
 \label{eq:bawd-generator}
\end{equation}
All members have the same initial logarithmic decay rate $k$. Finite $\rho$
gives a power-law tail, while $\rho=\infty$ is the exponential limit. The
cases $\gamma=0$, $0<\gamma<1$, and $\gamma=1$ mean, respectively, constant
clearance, more slowly fading clearance, and clearance matched to the drive.

Write the two accumulated clocks
\begin{equation}
 Q_\rho(u)=\int_0^u h_\rho(s)\dd s,\qquad
 R_{\rho,\gamma}(u)=\int_0^u h_\rho(s)^\gamma\dd s.
 \label{eq:bawd-clocks}
\end{equation}
For finite $\rho$ and $k>0$,
\begin{align}
 Q_\rho(u)
 &= \frac{\rho}{k(1-\rho)}
    \left[\left(1+\frac{ku}{\rho}\right)^{1-\rho}-1\right],
 && \rho\ne1, \nonumber\\
 Q_1(u)&=\frac{1}{k}\log(1+ku),&& \nonumber\\
 R_{\rho,\gamma}(u)
 &= \frac{\rho}{k(1-\rho\gamma)}
    \left[\left(1+\frac{ku}{\rho}\right)^{1-\rho\gamma}-1\right],
 && \rho\gamma\ne1, \nonumber\\
R_{\rho,\gamma}(u)=\frac{\rho}{k}\log\left(1+\frac{ku}{\rho}\right),
&& \rho\gamma=1.
 \label{eq:bawd-finite-clocks}
\end{align}
The continuous $k\to0$ limits are $Q_\rho(u)=R_{\rho,\gamma}(u)=u$.
For the exponential member,
\begin{equation}
 Q_\infty(u)=\frac{1-\e^{-ku}}{k},\qquad
 R_{\infty,\gamma}(u)=
 \begin{cases}
   u,&\gamma=0,\\
   (1-\e^{-\gamma ku})/(\gamma k),&\gamma>0.
 \end{cases}
 \label{eq:bawd-exponential-clocks}
\end{equation}
Integrating \eqref{eq:bawd-generator} therefore gives
\begin{equation}
 X(u)=z+VQ_\rho(u)-\ell R_{\rho,\gamma}(u),\qquad
 V^*(u,z)=\frac{b-z+\ell R_{\rho,\gamma}(u)}{Q_\rho(u)},
 \label{eq:bawd-traj}
\end{equation}
and
\begin{equation}
 -\partial_u V^*(u,z)
 =\frac{V^*(u,z)h_\rho(u)-\ell h_\rho(u)^\gamma}{Q_\rho(u)}.
 \label{eq:bawd-vstar}
\end{equation}
The original constant-clearance exponential model is the corner
$(\rho,\gamma)=(\infty,0)$.

\subsection{Tangency, critical launch, and the live/frozen split}
\label{sec:bawd-saturation}
Suppose first that $0\le\gamma<1$ and $k,\ell>0$, and put
$\nu=1-\gamma$. The required launch is U-shaped: the trajectory initially
rises, reaches a peak, and eventually falls because the clearance decays more
slowly than the drive. At tangency,
\[
  V_c(z)h_\rho(T_{\rm sat}(z))
  =\ell h_\rho(T_{\rm sat}(z))^\gamma.
\]
For an actual launch with $V>\ell$, the trajectory peak is at
\[
 u_{\rm peak}(V)=
 \begin{cases}
 \dfrac{\rho}{k}\left[
   \left(\dfrac{V}{\ell}\right)^{1/[\rho(1-\gamma)]}-1\right],
   &\rho<\infty,\\[8pt]
 \dfrac{\log(V/\ell)}{k(1-\gamma)},&\rho=\infty.
 \end{cases}
\]
It depends on the launch strength but not on the start point. The saturation
time $T_{\rm sat}(z)$ below is instead the peak time of the launch that is
just sufficient for start point $z$.
For $z<b$ (which holds almost surely unless $A=b$ and the upper endpoint is
selected), the saturation time is the unique positive solution of
\begin{equation}
 \frac{k(b-z)}{\ell}
 = \Theta_{\rho,\gamma}(x),\qquad x=kT_{\rm sat}(z),
 \label{eq:bawd-sat}
\end{equation}
where
\begin{equation}
 \Theta_{\rho,\gamma}(x)
 = k\left[
 h_\rho(x/k)^{-\nu}Q_\rho(x/k)
 -R_{\rho,\gamma}(x/k)\right],
 \qquad
 V_c(z)=\ell h_\rho(x/k)^{-\nu}.
 \label{eq:bawd-tangent-map}
\end{equation}
The critical launch on the frozen branch is therefore
\begin{equation}
 V_{\mathrm{crit}}(z)=V_c(z)
 =\ell h_\rho(T_{\rm sat}(z))^{-(1-\gamma)}.
 \label{eq:bawd-crit}
\end{equation}
For finite $\rho$, with $L=\log(1+x/\rho)$, this map is
\begin{equation}
 \Theta_{\rho,\gamma}(x)
 =\rho\left[
 \e^{\rho\nu L}\frac{\e^{(1-\rho)L}-1}{1-\rho}
 -\frac{\e^{(1-\rho\gamma)L}-1}{1-\rho\gamma}
 \right],
 \label{eq:bawd-theta-finite}
\end{equation}
with the displayed quotients interpreted by continuity when a denominator
vanishes. In the exponential limit,
\begin{equation}
 \Theta_{\infty,\gamma}(x)
 =\e^{\nu x}(1-\e^{-x})
 -\begin{cases}
    x,&\gamma=0,\\
    (1-\e^{-\gamma x})/\gamma,&\gamma>0.
  \end{cases}
 \label{eq:bawd-theta-exp}
\end{equation}
The map is increasing from zero to infinity. For the constant-clearance
exponential corner, \eqref{eq:bawd-sat} reduces to
$\e^x-1-x=k(b-z)/\ell$, the familiar Lambert equation.

The $\rho=2,\gamma=0$ member makes the change of geometry explicit:
\[
 T_{\rm sat}(z)=\sqrt{\frac{2(b-z)}{k\ell}},\qquad
 T_{\max}=\sqrt{\frac{2b}{k\ell}},\qquad
 V_c(z)=\left(\sqrt{\ell}+
       \sqrt{\frac{k(b-z)}{2}}\right)^2.
\]
Thus the same boundary and clearance can imply a different relation between
the response endpoint and the omission threshold than the exponential
member.

The first-passage hit set uses the running minimum of $V^*$:
\begin{equation}
 V^*_{\min}(u,z)=
 \begin{cases}
   V^*(u,z),&u\le T_{\rm sat}(z)\quad\text{(live)},\\
   V_c(z),&u>T_{\rm sat}(z)\quad\text{(frozen)}.
 \end{cases}
 \label{eq:bawd-minsplit}
\end{equation}
The start points still live at time $u$ when $z\le z^*(u)$, where
\begin{equation}
 z^*(u)=b-\frac{\ell}{k}\Theta_{\rho,\gamma}(ku),\qquad
 Z(u)=\min\{\max\{z^*(u),0\},A\}.
 \label{eq:bawd-Z}
\end{equation}
The endpoint is $T_{\max}=T_{\rm sat}(0)$: for
$u\ge T_{\max}$ every start point is frozen, the distribution function is
flat, and the density is zero. The non-finishing mass is intrinsic,
\begin{equation}
  1-F_{\max}
  =\frac{1}{A}\int_0^A P\{V<V_c(z)\}\dd z
 \label{eq:bawd-fmax}
\end{equation}
for $A>0$, with the evident point-start value when $A=0$.

\paragraph{The endpoint chart.} Equation \eqref{eq:bawd-sat} at $z=0$ reads
$\Theta_{\rho,\gamma}(kT_{\max}) = kb/\ell$, and $\Theta_{\rho,\gamma}$ is
given in closed form by \eqref{eq:bawd-theta-finite} and
\eqref{eq:bawd-theta-exp}. The equation has no closed-form solution for
$T_{\max}$, whereas solving it for $\ell$ is explicit:
\begin{equation}
  \ell = \frac{kb}{\Theta_{\rho,\gamma}(kT_{\max})}.
  \label{eq:bawd-ell-from-tmax}
\end{equation}
So for $\gamma<1$ and $k>0$ the pair $(k,\ell)$ and the pair $(k,T_{\max})$ are
two charts on the same model, related by a closed-form bijection, and they
define the same response-time law. In the endpoint chart,
$T_{\max}$ is taken as the coordinate and \eqref{eq:bawd-ell-from-tmax}
gives the corresponding clearance; the rate chart instead solves for
$T_{\max}$.

The distinction is geometric rather than arithmetic: in the likelihood
surface, $T_{\max}$ aligns with the principal endpoint direction while $\ell$
is the complementary coordinate, so the endpoint chart puts the ridge along a
coordinate axis instead of across it.
Two boundaries are inert rather than degenerate. $T_{\max}=\infty$ gives
$\ell=0$, the no-clearance model with unbounded support, and $k=0$ gives a
linear trace with no finite maximum for any $\ell$, so the endpoint carries no
information there. The model is then an LBA in the effective drive $V-\ell$,
so $\ell$ is an additive offset rather than an endpoint parameter. The
$\gamma=1$ corner has no finite endpoint at all and keeps the rate chart.

The matched-decay corner $\gamma=1$ is qualitatively different. In that case
$R_{\rho,1}=Q_\rho$ and
\begin{equation}
 X(u)=z+(V-\ell)Q_\rho(u),\qquad
 V^*(u,z)=\ell+\frac{b-z}{Q_\rho(u)}.
 \label{eq:bawd-matched}
\end{equation}
There is no finite endpoint. Let
\[
 Q_\rho(\infty)=
 \begin{cases}
   \infty,&\rho=1,\\
   \rho/[k(\rho-1)],&1<\rho<\infty,\\
   1/k,&\rho=\infty.
 \end{cases}
\]
Then $V^*(u,z)$ decreases to
\begin{equation}
 V_{\rm crit}(z)=\ell+\frac{b-z}{Q_\rho(\infty)},
 \label{eq:bawd-matched-crit}
\end{equation}
where the second term is zero when $Q_\rho(\infty)=\infty$. Hence the
matched $\rho=1$ member responds eventually exactly when $V>\ell$, whereas
for $\rho>1$ the finite total clock leaves a boundary-dependent omission
threshold.

\subsection{Distribution, density, and the frozen Jacobian}
For $0\le\gamma<1$ and $A>0$, define
\[
 w_{\rm hi}(u)=V^*(u,0),\qquad w_{\rm lo}(u)=V^*(u,Z(u)).
\]
For a truncated normal launch, the following distribution formulas are
divided by $Z_+$; for an untruncated normal, an ordinary or split-lognormal,
or a Weibull launch, $Z_+=1$. On the live range, $V^*$ is affine in $z$, so
\begin{equation}
 \Psi_{\rm live}(u)
 =Q_\rho(u)\left[C_{\mathcal L}(w_{\rm lo})-C_{\mathcal L}(w_{\rm hi})\right],
 \label{eq:bawd-live}
\end{equation}
where $C_{\mathcal L}$ is the appropriate stop-loss transform: the normal
$C$ in \eqref{eq:stoploss}, the ordinary-lognormal $C$ in \eqref{eq:bs},
$C_{\rm sp}$ in \eqref{eq:split-stoploss}, or $C_W$ in
\eqref{eq:weibull-tail-moments}.
The full distribution function is
\begin{equation}
 F_{\rm BAwD}(u)
 =\frac{1}{A}\left\{\Psi_{\rm live}(u)
 +\int_{Z(u)}^A\Gbar(V_c(z))\dd z\right\}.
 \label{eq:bawd-cdf}
\end{equation}
The point-start distribution is the survivor at $w_{\rm hi}$ until it
freezes, and the corresponding density is
$g_{\mathcal L}(w_{\rm hi})[w_{\rm hi}h_\rho-\ell h_\rho^\gamma]/Q_\rho$ on the live
branch and zero after the endpoint, where $g_{\mathcal L}$ is the selected
launch density. For $A>0$,
\begin{equation}
 A f_{\rm BAwD}(u)
 =\int_{w_{\rm lo}}^{w_{\rm hi}}
     g(w)\left[w h_\rho(u)-\ell h_\rho(u)^\gamma\right]\dd w.
 \label{eq:bawd-pdf}
\end{equation}
For the ordinary lognormal launch this becomes
\begin{equation}
 f_{\rm BAwD}(u)
 =\frac{1}{A}\left\{
 h_\rho(u)\e^{\mu+\sigma^2/2}\Delta\PhiN_{d_1}
 -\ell h_\rho(u)^\gamma\Delta\PhiN_{d_2}
 \right\},
 \label{eq:bawd-pdf-logn}
\end{equation}
where $\Delta\PhiN_{d_j}=\PhiN(d_j(w_{\rm lo}))
-\PhiN(d_j(w_{\rm hi}))$. The normal-launch version is the corresponding
weighted normal partial-moment integral and is closed in $\PhiN,\phiN$.
More explicitly, for a normal launch let
\[
 M_1(x)=v\PhiN\!\left(\frac{v-x}{s_v}\right)
       +s_v\phiN\!\left(\frac{x-v}{s_v}\right).
\]
Then, on the live branch,
\begin{equation}
 f_{\mathrm{BAwD}}(u)
 =\frac{1}{A Z_+}\left\{
 h_\rho(u)\left[M_1(w_{\rm lo})-M_1(w_{\rm hi})\right]
 -\ell h_\rho(u)^\gamma
   \left[\PhiN\!\left(\frac{v-w_{\rm lo}}{s_v}\right)
        -\PhiN\!\left(\frac{v-w_{\rm hi}}{s_v}\right)\right]
 \right\}.
\label{eq:bawd-pdf-norm}
\end{equation}

For a split-lognormal launch, define
$\Delta M_r=M_r(w_{\rm lo})-M_r(w_{\rm hi})$ using
\eqref{eq:split-tail-moments}. The live density is
\begin{equation}
 f_{\mathrm{BAwD,sp}}(u)
 =\frac{1}{A}\left\{
 h_\rho(u)\,\Delta M_1
 -\ell h_\rho(u)^\gamma\,\Delta M_0\right\},
 \label{eq:bawd-pdf-split}
\end{equation}
which reduces to \eqref{eq:bawd-pdf-logn} when $\delta=0$.

For a Weibull launch, the corresponding live density is
\begin{equation}
 f_{\mathrm{BAwD},W}(u)
 =\frac{1}{A}\left\{
 h_\rho(u)\,\Delta M_{1,W}
 -\ell h_\rho(u)^\gamma\,\Delta M_{0,W}\right\},
 \qquad
 \Delta M_{q,W}=M_{q,W}(w_{\rm lo})-M_{q,W}(w_{\rm hi}).
 \label{eq:bawd-pdf-weibull}
\end{equation}

The frozen branch is best expressed in the critical launch $w=V_c(z)$.
For $0\le\gamma<1$, its Jacobian is
\begin{equation}
 \left|\frac{\dd z}{\dd w}\right|=
 \begin{cases}
 \dfrac{\rho}{k(\rho-1)}
   \left[1-\left(\dfrac{\ell}{w}\right)^\alpha\right],
 &1<\rho<\infty,\quad
   \alpha=\dfrac{\rho-1}{\rho(1-\gamma)},\\[10pt]
 \dfrac{1}{k(1-\gamma)}\log\left(\dfrac{w}{\ell}\right),
 &\rho=1,\\[10pt]
 \dfrac{1}{k}
   \left[1-\left(\dfrac{\ell}{w}\right)^{1/(1-\gamma)}\right],
 &\rho=\infty.
 \end{cases}
 \label{eq:bawd-frozen-jacobian}
\end{equation}
The $\rho=1$ line is the continuous logarithmic limit of the finite-$\rho$
line. If $w_a=V_c(A)$ and $w_b=V_c(Z(u))$, the frozen contribution is
\begin{equation}
 \Psi_{\rm fr}(u)=\int_{w_a}^{w_b}
   \Gbar(w)\left|\frac{\dd z}{\dd w}\right|\dd w.
 \label{eq:bawd-frozen}
\end{equation}
For a positive launch with closed tail primitives define, for real $\alpha$,
\begin{equation}
  P_\alpha(x)=\int_x^\infty w^{-\alpha}\Gbar(w)\dd w,\qquad
  L(x;\ell)=\int_x^\infty \log(w/\ell)\Gbar(w)\dd w.
  \label{eq:bawd-power-primitives}
\end{equation}
Then the finite-$\rho$ frozen branch is
\begin{equation}
 \Psi_{\rm fr}(u)
 =\frac{\rho}{k(\rho-1)}
 \left\{C(w_a)-C(w_b)
 -\ell^\alpha\left[P_\alpha(w_a)-P_\alpha(w_b)\right]\right\},
 \qquad 1<\rho<\infty,
 \label{eq:bawd-frozen-power}
\end{equation}
the $\rho=1$ branch is
\begin{equation}
 \Psi_{\rm fr}(u)
 =\frac{L(w_a;\ell)-L(w_b;\ell)}{k(1-\gamma)},
 \label{eq:bawd-frozen-log}
\end{equation}
and the exponential branch is \eqref{eq:bawd-frozen-power} in the
$\rho\to\infty$ limit, with prefactor $1/k$ and
$\alpha=1/(1-\gamma)$. For the ordinary lognormal, $P_\alpha$ is a
truncated lognormal moment and $L$ is its logarithmic limit. For the
split-lognormal, use $P_\alpha^{\rm sp}$ and $L^{\rm sp}$ from
\eqref{eq:split-power-primitives}. For the Weibull launch, use
$P_{\alpha,W}$ and $L_W$ from \eqref{eq:weibull-power-primitives}; these are
incomplete-gamma expressions. All three branches are therefore closed for
ordinary and split-lognormal launches through Gaussian-tail primitives and for
Weibull launches through incomplete-gamma primitives. For a normal launch,
these power- or logarithm-weighted frozen integrals generally require
one-dimensional quadrature. Thus the ordinary, split-lognormal, and Weibull
launches close the staggered-saturation branch.
In the displayed frozen formulas, the split case is obtained by replacing
$C$, $P_\alpha$, and $L$ with $C_{\rm sp}$, $P_\alpha^{\rm sp}$, and
$L^{\rm sp}$, respectively.

\subsection{Endpoint order}
For $0\le\gamma<1$ and $k,\ell>0$, the live weight
$w h_\rho-\ell h_\rho^\gamma$ has a simple zero at the lowest-start
tangency. With a point start the density therefore vanishes linearly at
$T_{\max}$; when $A>0$, $Z(u)$ also vanishes linearly and the density
vanishes quadratically. The indices $(\rho,\gamma)$ change the coefficients
through \eqref{eq:bawd-frozen-jacobian}, but not these generic orders.
\subsection{Matched-decay distribution}
When $\gamma=1$, every start point remains on the live, monotone branch.
For $A>0$, put
\[
 w_{\rm hi}(u)=\ell+\frac{b}{Q_\rho(u)},\qquad
 w_{\rm lo}(u)=\ell+\frac{b-A}{Q_\rho(u)}.
\]
The affine start-point average is
\begin{equation}
 F_{\mathrm{BAwD},\gamma=1}(u)
 =\frac{Q_\rho(u)}{A Z_+}
   \left[C(w_{\rm lo})-C(w_{\rm hi})\right],
\label{eq:bawd-matched-cdf}
\end{equation}
and its density is
\begin{equation}
 A f_{\mathrm{BAwD},\gamma=1}(u)
 =\frac{h_\rho(u)}{Z_+}
   \int_{w_{\rm lo}}^{w_{\rm hi}}(w-\ell)g(w)\dd w.
\label{eq:bawd-matched-pdf}
\end{equation}
The point-start formulas follow by taking $A\to0$. These expressions remain
valid as $u\to\infty$ whenever $Q_\rho(\infty)$ is finite, and their limit is
the matched-decay response probability in \eqref{eq:bawd-matched-crit}.
%======================================================================
\section{Ballistic Accumulator with Proportional Clearance}
\label{sec:bawdp}
%======================================================================
For BAwDp, \(q_\lambda\), its running maximum, and the positive/negative
branch split are parent-clock quantities; the physical clock is obtained by
the composition in \eqref{eq:time-warp-bawdp}.
\subsection{Separability and a Universal Freeze Time}
Take the generator \eqref{eq:bawdp-gen} with $k > 0$ and $0 \le \lambda < 1$.
The case $\lambda=0$ is the pure drive-decay limit and has no finite stop
time; the hard-endpoint discussion below assumes $0<\lambda<1$. Integrating
from $X(0)=z$,
\begin{equation}
  X(u) = z + V\,q_{\lambda}(u),
  \qquad
  q_{\lambda}(u) = \frac{1 - \e^{-ku}}{k} - \lambda u = q(u) - \lambda u,
  \label{eq:bawdp-traj}
\end{equation}
with $q=Q_\infty$ as the exponential base clock, so that $q_0=q$ exactly.
The trajectory separates: the start point is static, and the time profile
$q_{\lambda}$ carries no trial-level randomness. For $0<\lambda<1$, the
hard-endpoint construction below applies first to positive launches ($V>0$),
which includes ordinary and split-lognormal launches, Weibull launches, and positive-truncated
normal launches. An
untruncated normal launch has a separate negative-launch branch described
below. For $V>0$, $q_{\lambda}'(u)=\e^{-ku}-\lambda$ makes the trajectory rise
until
\begin{equation}
  u_{*} = \frac{1}{k}\log\frac{1}{\lambda},
  \label{eq:bawdp-ustar}
\end{equation}
and fall thereafter. The positive part of the clock reaches
\begin{equation}
  q_{\lambda}^{\max}=q_{\lambda}(u_{*})
  =\frac{1-\lambda+\lambda\log\lambda}{k}>0,
  \qquad 0<\lambda<1,
  \label{eq:bawdp-mmax}
\end{equation}
while $q_{\lambda}(u)\to-\infty$ as $u\to\infty$ for $0<\lambda<1$. On the
positive-clock branch $q_{\lambda}(u)>0$, the required launch strength is
\begin{equation}
  V^{*}(u,z) = \frac{b-z}{q_{\lambda}(u)},
  V^{*}_{\min}(u,z) = \frac{b-z}{\bar q_{\lambda}(u)},
  \qquad
  \bar q_{\lambda}(u) = q_{\lambda}\!\left(\min(u, u_{*})\right).
  \label{eq:bawdp-vstar}
\end{equation}
For $0<\lambda<1$ and $0<u<u_{*}$, $V^*$ is U-shaped and its running
minimum freezes at $u_{*}$; this describes all hits for positive launches.
The frozen threshold remains affine in $z$ with slope
$-1/q_{\lambda}^{\max}$.
\subsection{The Positive-Launch Likelihood}
For a positive launch law, $\bar q_{\lambda}$ does not depend on $z$, so the
start-point average in Step~3 of Section~\ref{sec:ballistic-idea-recipe} is
the ordinary LBA average evaluated at internal time $\bar q_{\lambda}(u)$.
Writing $F_{\mathrm{LBA}}$ and $f_{\mathrm{LBA}}$ for the distribution function
and density at the same $(A,b)$ and positive launch law,
\begin{equation}
  \boxed{\;
  F_{\mathrm{BAwDp}}(u) = F_{\mathrm{LBA}}\!\left(\bar q_{\lambda}(u)\right),
  \qquad
  f_{\mathrm{BAwDp}}(u) =
  \begin{cases}
    f_{\mathrm{LBA}}\!\left(q_{\lambda}(u)\right) q_{\lambda}'(u),
      & 0 < u < u_{*},\\[3pt]
    0, & u \ge u_{*}.
  \end{cases}\;}
  \label{eq:bawdp-law}
\end{equation}
This is an exact identity for positive launches, not an approximation, and it
holds for every $A\ge0$. Three points about it are worth stating explicitly.

\begin{itemize}
  \item \textbf{The running maximum matters.} $q_{\lambda}$ decreases after
        $u_{*}$, so using $q_{\lambda}$ in place of $\bar q_{\lambda}$ in
        \eqref{eq:bawdp-law} would give a distribution function that runs
        backwards and a density $f_{\mathrm{LBA}}(q_{\lambda})q_{\lambda}'$ that
        is negative. First passage depends on the running maximum of the
        trajectory, and \eqref{eq:bawdp-vstar} implements that.
  \item \textbf{No new primitives are needed.} The right-hand sides of
        \eqref{eq:bawdp-law} are \eqref{eq:lba-cdf} and \eqref{eq:lba-pdf} for a
        normal launch, or the stop-loss forms of Section~\ref{sec:bawl-logn}
        with $k = 0$ for an ordinary, split-lognormal, or Weibull launch, at a warped
        argument. There is no
        frozen branch to integrate, no saturation exponent, no implicit
        inversion, and no quadrature.
  \item \textbf{The response-time map is a monotone inversion.} On $(0, u_{*})$,
        $q_{\lambda}$ is strictly increasing, so if $S$ denotes the LBA
        finishing time in internal units then
        $u = q_{\lambda}^{-1}(S)$ when $S < q_{\lambda}^{\max}$ and the trial is
        an omission otherwise. The inversion is one-dimensional, convex on the
        rising branch, and bracketed by $[0, u_{*}]$.
\end{itemize}

\subsection{Critical Launch, Endpoint, and Defect for Positive Launches}
For positive launch laws, the critical launch strength that arrives exactly as
the positive clock stops is
\begin{equation}
  V_{\mathrm{crit}}(z) = \frac{b - z}{q_{\lambda}^{\max}}
  = \frac{k(b-z)}{1 - \lambda + \lambda\log\lambda},
  \label{eq:bawdp-crit}
\end{equation}
It is affine and decreasing in $z$: a favourable start point genuinely rescues
a trial, exactly as in the positive branch of the drive-decay family and unlike
the leak model.
The positive-launch response probability and defect are single evaluations of
the LBA distribution function at the end of the clock,
\begin{equation}
  F_{\mathrm{BAwDp}}(\infty)=F_{\mathrm{LBA}}\!\left(q_{\lambda}^{\max}\right),
  \qquad
  \text{defect}=1-F_{\mathrm{LBA}}\!\left(q_{\lambda}^{\max}\right),
  \label{eq:bawdp-defect}
\end{equation}
closed for all four positive launch laws (with Weibull terms in incomplete
gamma functions). The positive
response component has bounded support $(t_0,t_0+u_{*}]$, and its endpoint is
\begin{equation}
  T_{\max}=u_{*}=\frac{\log(1/\lambda)}{k},
\end{equation}
free of $b$, of $A$, and of the launch law. The window is a property of the
waveform. In a race of positive-launch components sharing $(k,\lambda)$, all
components share this endpoint.

\paragraph{Endpoint order for positive launches.} Near the endpoint put
$\epsilon=u_{*}-u$. Then $q_{\lambda}'(u)=\e^{-ku}-\lambda
\approx k\lambda\,\epsilon$ vanishes linearly, while
$f_{\mathrm{LBA}}(q_{\lambda}(u))\to f_{\mathrm{LBA}}(q_{\lambda}^{\max})$,
which is generically strictly positive. Thus the positive-launch density
vanishes \emph{linearly} at the endpoint for every $A$,
including $A > 0$. The drive-decay model vanishes linearly at $A = 0$ and
quadratically at $A > 0$ (Section~\ref{sec:bawd}), because there the live
start-point range collapses at the same time as the live weight. This is a
genuine shape restriction: the proportional-clearance model cannot produce the
drive-decay model's softest right shoulder, and $A$ leaves no endpoint
signature here.
\subsection{Untruncated normal launch and sign reversal}
The positive-launch endpoint is not the whole story for an untruncated normal.
Assume $0<\lambda<1$ and let $u_0>u_{*}$ be the second zero of the clock,
\begin{equation}
 q_{\lambda}(u_0)=0,\qquad
 r(u)=-q_{\lambda}(u)>0\quad (u>u_0).
 \label{eq:bawdp-u0}
\end{equation}
Because $q_{\lambda}(u)\to-\infty$, a negative launch first moves away from
the boundary and then reverses direction after $u_0$; every $V<0$ eventually
crosses. For $A>0$, its contribution to the distribution function is
\begin{equation}
 F_-(u)=
 \begin{cases}
 0,&u\le u_0,\\[3pt]
 \dfrac{r(u)s_v}{A}
 \left[
 \Psi\!\left(\dfrac{A-b-r(u)v}{r(u)s_v}\right)
 -\Psi\!\left(\dfrac{-b-r(u)v}{r(u)s_v}\right)
 \right],&u>u_0,
 \end{cases}
 \label{eq:bawdp-negative-cdf}
\end{equation}
where $\Psi$ is the normal survivor antiderivative in \eqref{eq:psi}. For a
point start, replace the second line by
$\PhiN((-b-r(u)v)/(r(u)s_v))$. Thus the complete untruncated-normal law is
\begin{equation}
 F_{\mathrm{BAwDp},\mathrm{normal}}(u)
 =F_{\mathrm{LBA}}\!\left(\bar q_{\lambda}(u)\right)+F_-(u).
 \label{eq:bawdp-normal-full}
\end{equation}
The late density is also explicit. For $A>0$, put
$w_0=-b/r(u)$ and $w_A=-(b-A)/r(u)$; then
\begin{equation}
 f_-(u)=
 \begin{cases}
 0,&u\le u_0,\\[3pt]
 -\dfrac{r'(u)}{A}\left[M_1(w_0)-M_1(w_A)\right],&u>u_0,
 \end{cases}
 \qquad r'(u)=\lambda-\e^{-ku},
 \label{eq:bawdp-negative-pdf}
\end{equation}
with $M_1$ from the normal partial-moment formula above. For a point start, the
corresponding late density is
$f_-(u)=b\,r'(u)g(-b/r(u))/r(u)^2$ for $u>u_0$. In either case, the total
density is the positive-branch density in \eqref{eq:bawdp-law} plus $f_-$.
\begin{equation}
 F_{\mathrm{BAwDp},\mathrm{normal}}(\infty)
 =\PhiN\!\left(-\frac{v}{s_v}\right)
  +F_{\mathrm{LBA}}\!\left(q_{\lambda}^{\max}\right),
 \qquad
 \text{defect}=1-F_{\mathrm{BAwDp},\mathrm{normal}}(\infty).
 \label{eq:bawdp-normal-defect}
\end{equation}
The positive component occupies $(0,u_{*}]$, the negative component occupies
$(u_0,\infty)$, and there is no response density between them. Thus an
untruncated normal removes the global hard endpoint: the late branch has an
unbounded right tail. The $\lambda=0$ limit is singular, because the second
zero and the negative-launch branch recede to infinity.

\subsection{Relation to the Other Models}
\paragraph{It is a constrained drive-decay model.} Comparing
\eqref{eq:bawdp-gen} with \eqref{eq:two-generators}, the substitution is
$\ell \mapsto \lambda V$: the clearance is tied to the trialwise launch
strength rather than being a separate constant. The tie removes one degree of
freedom and, with it, the trial-specific peak time. Under the waveform reading
of \eqref{eq:bawdp-gen} nothing is being ``cleared'' at all: $V$ is the
amplitude of this trial's evidence pulse, $k$ sets how fast its positive part
decays, $\lambda$ sets where it crosses zero, and after $u_{*}$ the momentary
signal opposes what was accumulated earlier.

\paragraph{It nests the pure drive-decay model.} For positive launches, at
$\lambda=0$ the waveform never crosses zero, $u_{*}=\infty$,
$q_{\lambda}=q$, and $q_{\lambda}^{\max}=1/k$, so
$V_{\mathrm{crit}}(z)=k(b-z)$. That is exactly the $\ell\to0$ drive-decay
model of Section~\ref{sec:reductions}: unbounded support, no endpoint, but
still intrinsically defective, with the start point helping. The three models
line up as
\begin{equation}
  \underbrace{\dot X = V}_{\text{LBA}}
  \;\xrightarrow{\ k>0\ }\;
  \underbrace{\dot X = V\e^{-ku}}_{\substack{\text{pure drive decay}\\ u_{\max} = \infty}}
  \;\xrightarrow{\ \lambda>0\ }\;
  \underbrace{\dot X = V(\e^{-ku} - \lambda)}_{\substack{\text{proportional clearance}\\ u_{\max} = \log(1/\lambda)/k}},
\end{equation}
with the drive-decay model of Section~\ref{sec:bawd} branching off the middle
term instead by making the clearance independent of $V$.

\paragraph{It is not the leak model.} The two are easy to conflate, since both
end up as an LBA under a time change. The difference is the start point. Leak
maps $z \mapsto z\e^{-ku}$, so the initial evidence is itself forgotten, the
asymptote is $V/k$, and $V_{\mathrm{crit}} = kb$ is free of $z$; here the start
point never decays and $V_{\mathrm{crit}}$ depends on $z$ through
\eqref{eq:bawdp-crit}. The decaying start point is precisely what breaks the
separability $X = z + V\,q_{\lambda}(u)$, so for $A > 0$ no scalar time change
of an LBA reproduces the leak model, and no reparameterization maps one of
these models onto the other. For positive launches their observable signatures
differ accordingly: latest possible response time $\infty$ against $u_{*}$,
late density $\e^{-ku}$ against exact zero, and start-point variability
irrelevant to omissions against decisive for them. With an untruncated normal,
the proportional-clearance model also has the separate unbounded negative-
launch branch derived above.

\subsection{Closure and Scale}
The $I_2$ obstruction of Section~\ref{sec:closure} does not arise. It was
produced by the interaction of start-point variability with \emph{staggered}
saturation, which forced the saturation exponent as a coordinate and the
multiplicative measure $\dd w / w$ with it. Here saturation is simultaneous, so
the only coordinate ever needed is $z$ itself, $V^{*}$ stays affine in it, and
every start-point average is a stop-loss difference. Consequently the
proportional-clearance model has closed-form density, distribution function,
survivor, and omission mass under a normal launch, an ordinary lognormal, a
split-lognormal launch, or a Weibull launch. This is the one place in this note where the choice of
launch law has no bearing on closure at all.

Scale invariance works differently too. The map
$(V, b, A) \mapsto (cV, cb, cA)$ leaves \eqref{eq:bawdp-traj} and every
crossing time unchanged, and it does not touch $k$ or $\lambda$: both are
dimensionless in the evidence axis, $\lambda$ entirely so and $k$ carrying only
inverse time. So $\lambda$ cannot serve as the anchor the way $\ell = 1$ does
for the drive-decay model. Under a normal launch $s_v = 1$ still fixes the
scale; under an ordinary or split-lognormal launch, or a Weibull launch, it
must be fixed by a boundary convention, of which $b = 1$ (equivalently
$B = 1$ at $A = 0$) is the natural choice, after which $V$, $m_W$, and $k$
read as dimensionless.

%======================================================================
\section{Ballistic Accumulator with Global Fading}
\label{sec:bawf}
%======================================================================
BAwF evaluates its fading kernel and all endpoint expressions at
\(s=c_\eta(u)\), with the density Jacobian specified in
\eqref{eq:time-warp-bawf} and \eqref{eq:time-warp-law}.
\subsection{Trajectory and required launch}
The fading model takes the ordinary LBA trace and attenuates it:
\begin{equation}
  X(u) = h_\rho(u)\left[z + Vu\right],
  \qquad
  h_\rho(u) =
  \begin{cases}
    (1+ku/\rho)^{-\rho}, & \rho<\infty,\\
    \e^{-ku}, & \rho=\infty,
  \end{cases}
  \label{eq:bawf-traj}
\end{equation}
with $z\sim U(0,A)$ and $b=B+A$ as before. The index $\rho$ is fixed rather
than trialwise. Differentiating,
\begin{equation}
  \dot X = V h_\rho(u) - \frac{k}{1+ku/\rho}\,X,
  \qquad
  \dot X = V\e^{-ku} - kX \ \ (\rho=\infty),
  \label{eq:bawf-generator}
\end{equation}
so as a generator this is a leaky integrator driven by a fading drive. The
essential point is the one made in Section~\ref{sec:decay-targets}: the
clearance term is $kX$, which borrows its evidence scale from the state, where
the drive-decay clearance $\ell h_\rho^\gamma$ carries its own. BAwF therefore
has one decay parameter where BAwD has two.

It is convenient to work in the dimensionless phase $x=ku$ and the reciprocal
kernel
\begin{equation}
  H_\rho(x) = \frac{1}{h_\rho(x/k)} = \left(1+\frac{x}{\rho}\right)^{\rho},
  \qquad H_\infty(x)=\e^{x}.
  \label{eq:bawf-H}
\end{equation}
Solving $X(u)=b$ gives $z+Vu = bH_\rho(ku)$ and hence
\begin{equation}
  V^{*}(u,z) = \frac{b\,H_\rho(ku)-z}{u},
  \qquad
  -\partial_u V^{*}(u,z)
  = \frac{V^{*}(u,z) - k\,b\,H_\rho'(ku)}{u}.
  \label{eq:bawf-vstar}
\end{equation}
Two features of \eqref{eq:bawf-vstar} do most of the work below. First,
$V^{*}$ is affine in $z$ with slope $-1/u$, so the start-point Jacobian is
$|\dd z/\dd w| = u$ and the live branch is a stop-loss difference. Second, the
subtracted quantity in $-\partial_uV^{*}$ is $k\,b\,H_\rho'$, which depends on
the boundary and the fading rate but not on the launch strength; it plays
exactly the role that $\ell h_\rho^\gamma$ plays in \eqref{eq:bawd-vstar}, and
it is the reason the two models share their algebraic scaffolding.

\subsection{Tangency, the endpoint, and why $b$ cannot move it}
\label{sec:bawf-endpoint}
$V^{*}$ is U-shaped, and at its minimum
$V^{*}=k\,b\,H_\rho'(x)$ with $x=kT_{\rm sat}(z)$. Substituting
\eqref{eq:bawf-vstar} and cancelling $k/x$ gives the tangency condition in
closed form,
\begin{equation}
  z = b\,G_\rho(x),
  \qquad
  G_\rho(x) = H_\rho(x) - x H_\rho'(x)
            = H_\rho'(x)\left(1 - \frac{x}{x_{\max}}\right),
  \label{eq:bawf-tangency}
\end{equation}
with
\begin{equation}
  x_{\max} =
  \begin{cases}
    \dfrac{\rho}{\rho-1}, & 1<\rho<\infty,\\[6pt]
    1, & \rho=\infty.
  \end{cases}
  \label{eq:bawf-xmax}
\end{equation}
Since $G_\rho'(x) = -xH_\rho''(x) < 0$, $G_\rho$ decreases strictly from
$G_\rho(0)=1$ to $G_\rho(x_{\max})=0$, so \eqref{eq:bawf-tangency} inverts
uniquely and the live start points at time $u$ are those with $z\le z^{*}(u)$,
where
\begin{equation}
  z^{*}(u) = b\,G_\rho(ku),
  \qquad
  Z(u) = \min\{\max\{z^{*}(u),0\},A\}.
  \label{eq:bawf-Z}
\end{equation}
The critical launch on the frozen branch is
\begin{equation}
  V_c(z) = k\,b\,H_\rho'\!\left(x_{\rm sat}(z)\right),
  \qquad
  G_\rho\!\left(x_{\rm sat}(z)\right) = \frac{z}{b},
  \label{eq:bawf-crit}
\end{equation}
and in particular the omission boundary at the lowest start point is
\begin{equation}
  V_c(0) = k\,b\,H_\rho'(x_{\max})
  = k b \left(\frac{\rho}{\rho-1}\right)^{\rho-1}
  \ \longrightarrow\ \e\,kb \quad (\rho\to\infty).
  \label{eq:bawf-vc0}
\end{equation}

The endpoint is where every start point has frozen, $z^{*}(T_{\max})=0$:
\begin{equation}
  \boxed{\ T_{\max} = \frac{x_{\max}}{k}\ }
  \label{eq:bawf-tmax}
\end{equation}
Unlike every other finite-endpoint member of this family,
\eqref{eq:bawf-tmax} contains no boundary. The response window is fixed by the
fading rate and the kernel shape alone, while the omission boundary
\eqref{eq:bawf-vc0} is proportional to $b$ at fixed $k$. The two
quantities therefore have distinct structural roles, whereas in BAwD the endpoint
and the omission boundary are two consequences of the single entangled triple
$(b,k,\ell)$.

\begin{remark}
Why $b$ drops out is worth isolating, because the argument generalizes. Write
the
trace as $a(u)z + q(u)V$ with no additive term. Then $V^{*}=(b-a(u)z)/q(u)$ and
at $z=0$ the tangency condition $\partial_u[b/q(u)]=0$ is $q'(u)=0$: the
endpoint is $\arg\max q$, and $b$ cancels because it multiplies the whole
expression. BAwF has $a=h_\rho$, $q=uh_\rho$, and $q$ peaks at $x_{\max}/k$.
BAwD fails the condition, because its clearance contributes an additive
$-\ell R_{\rho,\gamma}(u)$, and so does BAwR. That is why both have
$b$-dependent windows. BAwDp satisfies it trivially with $a=1$,
$q=q_\lambda$.
\end{remark}

\begin{remark}
The excluded index $\rho=1$ makes the mechanism visible by removing it. There
$H_1(x)=1+x$, so $G_1\equiv1$, no start point ever freezes, and
$V^{*}(u,z)=(b-z)/u+kb$ falls monotonically to $kb$. The model then has no
endpoint, and the eventual-hit condition $V>kb$ is free of $z$: in value the
same omission threshold as the leak model \eqref{eq:bawl-mass}, reached from a
quite different trajectory. Since none of the identification architecture above
survives, $\rho=1$ is treated as a different model and excluded from the
family.
\end{remark}

\subsection{Distribution, density, and the frozen Jacobian}
Write $w_{\rm hi}(u)=V^{*}(u,0)=bH_\rho(ku)/u$ and
$w_{\rm lo}(u)=V^{*}(u,Z(u))$. As in Section~\ref{sec:bawd}, formulas are
divided by $Z_+$ for a truncated normal launch and by one otherwise. For a
split-lognormal launch the survivor and stop-loss are those in
\eqref{eq:split-survivor-moment}--\eqref{eq:split-stoploss}; for a Weibull
launch they are $\Gbar_W$ and $C_W$ from
\eqref{eq:weibull-launch}--\eqref{eq:weibull-tail-moments}. Because
$V^{*}$ is affine in $z$ with Jacobian $u$, the live contribution is a
stop-loss difference,
\begin{equation}
  \Psi_{\rm live}(u) = u\left[C(w_{\rm lo}) - C(w_{\rm hi})\right],
  \qquad
  F_{\rm BAwF}(u) = \frac{1}{A}\left\{
    \Psi_{\rm live}(u) + \int_{Z(u)}^{A}\Gbar\!\left(V_c(z)\right)\dd z
  \right\},
  \label{eq:bawf-cdf}
\end{equation}
and the density is the live weight only, since the frozen branch is constant
in $u$:
\begin{equation}
  A\,f_{\rm BAwF}(u)
  = \int_{w_{\rm lo}}^{w_{\rm hi}}
      g(w)\left[w - k\,b\,H_\rho'(ku)\right]\dd w.
  \label{eq:bawf-pdf}
\end{equation}
For the split-lognormal launch, \eqref{eq:bawf-pdf} is equivalently
$A f_{\rm BAwF}(u)=\Delta M_1-kbH_\rho'(ku)\Delta M_0$, with $\Delta M_r$
as above. In general this is \eqref{eq:bawd-pdf} with $\ell h_\rho(u)^\gamma$ replaced by
$k b H_\rho'(ku)$ and $h_\rho(u)$ replaced by one, so the ordinary lognormal,
split-lognormal, normal, and Weibull evaluations of \eqref{eq:bawd-pdf-logn},
\eqref{eq:bawd-pdf-split}, and \eqref{eq:bawd-pdf-norm} transfer
unchanged under that substitution. With $A=0$ the law is the pushforward
$F(u)=\Gbar_{\mathcal L}(w_{\rm hi}(u))$ on $(0,T_{\max}]$, frozen at
$\Gbar_{\mathcal L}(V_c(0))$ thereafter, with $\mathcal L$ denoting the
selected launch law.
For a Weibull launch, the live density is
\begin{equation}
  A f_{\mathrm{BAwF},W}(u)
   =\Delta M_{1,W}-kbH_\rho'(ku)\,\Delta M_{0,W}.
  \label{eq:bawf-pdf-weibull}
\end{equation}

For the frozen branch, parameterize by the saturation phase $s$ rather than by
$z$. From $w = V_c(s) = kbH_\rho'(s)$ and $z = bG_\rho(s)$,
\[
  \frac{\dd z}{\dd s} = -s\,b\,H_\rho''(s),
  \qquad
  \frac{\dd w}{\dd s} = k\,b\,H_\rho''(s),
\]
so the Jacobian is the phase itself, divided by the rate:
\begin{equation}
  \left|\frac{\dd z}{\dd w}\right| = \frac{s(w)}{k},
  \qquad
  s(w) =
  \begin{cases}
    \rho\left[\left(\dfrac{w}{kb}\right)^{1/(\rho-1)}-1\right],
      & 1<\rho<\infty,\\[10pt]
    \log\dfrac{w}{kb}, & \rho=\infty.
  \end{cases}
  \label{eq:bawf-frozen-jacobian}
\end{equation}
The frozen coordinate is simpler here than in BAwD: one line covers every
admissible $\rho$, where \eqref{eq:bawd-frozen-jacobian} needs three cases. It
induces the same primitives, with one change of sign in the exponent. With $w_a=V_c(A)$, $w_b=V_c(Z(u))$ and the
primitives of \eqref{eq:bawd-power-primitives},
\begin{align}
  \Psi_{\rm fr}(u)
  &= \frac{\rho}{k}\left\{
       (kb)^{-\beta}\left[P_{-\beta}(w_a)-P_{-\beta}(w_b)\right]
       -\left[C(w_a)-C(w_b)\right]\right\},
     \quad \beta=\frac{1}{\rho-1},\ \ \rho<\infty,
     \nonumber\\
  \Psi_{\rm fr}(u)
  &= \frac{1}{k}\left[L(w_a;kb)-L(w_b;kb)\right],
     \qquad \rho=\infty.
  \label{eq:bawf-frozen}
\end{align}
For a split-lognormal launch, the same expression uses
$P_{-\beta}^{\rm sp}$, $L^{\rm sp}$, and $C_{\rm sp}$ from
\eqref{eq:split-stoploss}--\eqref{eq:split-power-primitives}; for a Weibull
launch, use $P_{-\beta,W}$, $L_W$, and $C_W$ from
\eqref{eq:weibull-tail-moments}--\eqref{eq:weibull-power-primitives}.
Note that $P_{-\beta}(x)=\int_x^\infty w^{\beta}\Gbar(w)\dd w$ carries a
\emph{positive} power of $w$, where the drive-decay frozen branch
\eqref{eq:bawd-frozen-power} carries a negative one. Both are truncated
moments closed in $\PhiN,\phiN$ for either ordinary or split-lognormal launch,
and in incomplete gamma functions for a Weibull launch;
they generally require one-dimensional quadrature for a normal launch. The
closure statement of Section~\ref{sec:closure} is therefore unchanged: the
launch law matters only on the staggered frozen branch.

The intrinsic non-finishing mass is \eqref{eq:bawd-fmax} with
\eqref{eq:bawf-crit} substituted; for $A=0$ it is $P\{V<\e kb\}$ at
$\rho=\infty$.

\subsection{Endpoint order}
Evaluate the live weight of \eqref{eq:bawf-pdf} at the lowest start point:
\begin{equation}
  w_{\rm hi}(u) - k b H_\rho'(ku)
  = \frac{kb}{x}\left[H_\rho(x) - xH_\rho'(x)\right]
  = \frac{b\,G_\rho(ku)}{u}
  = \frac{z^{*}(u)}{u}.
  \label{eq:bawf-weight}
\end{equation}
The live weight at the lowest start point is therefore $z^{*}(u)/u$. That
identity holds for any model whose $V^{*}$ is affine in $z$ with Jacobian $u$,
and it makes the endpoint order immediate. $G_\rho$ has a simple
zero at $x_{\max}$, so the weight vanishes linearly at $T_{\max}$; with a point
start the density therefore shuts down linearly. When $A>0$ the integration
range $[w_{\rm lo},w_{\rm hi}]$ also collapses linearly, because $Z(u)=z^{*}(u)$
near the endpoint, and the density vanishes quadratically. These are the same
generic orders as BAwD, so the endpoint shoulder does not distinguish the two.

\subsection{Reductions and scale}
At $k=0$ the kernel is identically one and \eqref{eq:bawf-traj} is the LBA
exactly, with $T_{\max}\to\infty$ and $V_c\to0$; the limit is uniform in $\rho$
and is not merely an approximation, so $k=0$ is a genuine interior reduction.
As $\rho\to\infty$ the finite-$\rho$ kernels converge to the exponential
member and every formula above converges with them.

Scale invariance is simpler than in BAwD and, for that reason, more dangerous.
Multiplying $(V,b,A)$, and hence $z$, by a common factor $c$ multiplies
$X$ and the boundary by $c$ and leaves every crossing time unchanged; $k$,
$\rho$ and $T_{\max}$ are invariant. There is no $\ell$ to pin, so the
convention available for drive decay is not available here and exactly one
evidence-scale quantity must be fixed instead. With a normal launch the usual
$s_v=1$ does this. For an ordinary BAw* lognormal launch, the scale direction
acts on $m_L$ (or its log-mean intercept) while $c_L$ is dimensionless; for a
split-lognormal or BTAwL lognormal launch it acts on $\mu$; for a Weibull launch
it acts on $m_W$. Equivalently, one of these launch-scale coordinates, $B$, or
$A$ must be constrained explicitly.

%======================================================================
\section{Ballistic Accumulator with Ramping Clearance}
\label{sec:bawr}
%======================================================================
The BAwR ramp is a parent-clock function of \(s\) under the eta extension;
the transformed trajectory and required launch are given in
\eqref{eq:time-warp-bawr}.
\subsection{Trajectory and required launch}
The last variant removes the decay from the accumulator entirely and imposes it
from outside. The drive does not fade, the state is not forgotten, and the
opposing term is a function of elapsed time alone:
\begin{equation}
  \dot X(u) = V - \kappa u^{p},
  \qquad
  X(u) = z + Vu - \frac{\kappa u^{p+1}}{p+1},
  \qquad \kappa>0,\ p>0.
  \label{eq:bawr-traj}
\end{equation}
Both $\kappa$ and $p$ are continuous parameters; unlike $\rho$ and
$\gamma$ there is no fixed grid, because nothing in the derivation below
requires a special value. Solving $X(u)=b$,
\begin{equation}
  V^{*}(u,z) = \frac{b-z}{u} + \frac{\kappa u^{p}}{p+1},
  \qquad
  -\partial_u V^{*}(u,z) = \frac{V^{*}(u,z)-\kappa u^{p}}{u}.
  \label{eq:bawr-vstar}
\end{equation}
As in BAwF, $V^{*}$ is affine in $z$ with Jacobian $|\dd z/\dd w|=u$, and the
quantity subtracted in $-\partial_uV^{*}$ is the instantaneous clearance
$c_u=\kappa u^{p}$, free of $V$ and of $z$. The three finite-endpoint models
thus share one algebraic skeleton and differ only in what $c_u$ is:
$\ell h_\rho(u)^\gamma$ for BAwD, $kbH_\rho'(ku)$ for BAwF, and $\kappa u^p$
here. Note that this is the only one of the three that increases in $u$.

\subsection{Tangency, the endpoint, and its dependence on $b$}
$V^{*}$ is the sum of a decreasing and an increasing term and is therefore
U-shaped for any $p>0$. Setting $\partial_uV^{*}=0$,
\begin{equation}
  T_{\rm sat}(z)
  = \left[\frac{(p+1)(b-z)}{\kappa p}\right]^{1/(p+1)},
  \qquad
  V_c(z) = \kappa\,T_{\rm sat}(z)^{p},
  \label{eq:bawr-sat}
\end{equation}
where the second identity follows by substituting
$b-z=\kappa p\,T_{\rm sat}^{p+1}/(p+1)$ into \eqref{eq:bawr-vstar} and
observing that the two terms combine to $\kappa T_{\rm sat}^{p}$ exactly.
Unlike \eqref{eq:bawd-sat}, no transcendental inversion is required: the whole
geometry is available in closed form. The live start points are those with
$z\le z^{*}(u)$, where
\begin{equation}
  z^{*}(u) = b - \frac{\kappa p\,u^{p+1}}{p+1}
           = b\left[1-\left(\frac{u}{T_{\max}}\right)^{p+1}\right],
  \qquad
  Z(u)=\min\{\max\{z^{*}(u),0\},A\},
  \label{eq:bawr-Z}
\end{equation}
and the endpoint is
\begin{equation}
  \boxed{\ T_{\max} = T_{\rm sat}(0)
        = \left[\frac{(p+1)\,b}{\kappa p}\right]^{1/(p+1)}\ }
  \label{eq:bawr-tmax}
\end{equation}

\begin{remark}
$T_{\max}$ \emph{does} depend on $b$, and the contrast with
\eqref{eq:bawf-tmax} is instructive rather than incidental. The remark in
Section~\ref{sec:bawf-endpoint} requires the trace to have the multiplicative
form $a(u)z+q(u)V$; \eqref{eq:bawr-traj} has an additive
$-\kappa u^{p+1}/(p+1)$ that survives at $z=0$, so $b$ does not cancel and the
argument does not apply. Concretely, doubling $b$ multiplies $T_{\max}$ by
$2^{1/(p+1)}$: by $\sqrt2$ at $p=1$, and an eightfold increase in $b$ doubles
the window at $p=2$. Increasing $b$ therefore increases the response window;
BAwF and BAwDp have no analogous dependence.

One consequence is that $\kappa$ and $b$ cannot both be given window-free
interpretations. Taking $\kappa$ as the coordinate and deriving $T_{\max}$
keeps $B$ as a pure boundary offset; taking the window as the coordinate and
deriving $\kappa$ couples $B$ to the endpoint.
\end{remark}

\subsection{Distribution, density, and the frozen Jacobian}
With $w_{\rm hi}(u)=V^{*}(u,0)$, $w_{\rm lo}(u)=V^{*}(u,Z(u))$ and $Z_+$ as
before, the live branch is again a stop-loss difference with Jacobian $u$ and
the density is the live weight alone:
\begin{align}
  F_{\rm BAwR}(u)
  &= \frac{1}{A}\left\{u\left[C(w_{\rm lo})-C(w_{\rm hi})\right]
     + \int_{Z(u)}^{A}\Gbar\!\left(V_c(z)\right)\dd z\right\},
  \label{eq:bawr-cdf}\\
  A\,f_{\rm BAwR}(u)
  &= \int_{w_{\rm lo}}^{w_{\rm hi}}
       g(w)\left[w-\kappa u^{p}\right]\dd w.
  \label{eq:bawr-pdf}
\end{align}
For a split-lognormal launch, the $C$ and $\Gbar$ in the first line are
$C_{\rm sp}$ and $\Gbar_{\rm sp}$, respectively; for a Weibull launch they
are $C_W$ and $\Gbar_W$.
For an ordinary lognormal launch \eqref{eq:bawr-pdf} evaluates as in
\eqref{eq:bawd-pdf-logn},
\begin{equation}
  f_{\rm BAwR}(u)
  = \frac{1}{A}\left\{
      \e^{\mu+\sigma^2/2}\,\Delta\PhiN_{d_1}
      - \kappa u^{p}\,\Delta\PhiN_{d_2}\right\},
  \label{eq:bawr-pdf-logn}
\end{equation}
and for a normal launch as in \eqref{eq:bawd-pdf-norm} with $h_\rho\mapsto1$
and $\ell h_\rho^\gamma\mapsto\kappa u^{p}$.
For a split-lognormal launch, with $\Delta M_r$ from
\eqref{eq:split-tail-moments},
\begin{equation}
  f_{\rm BAwR,sp}(u)
  =\frac{1}{A}\left[\Delta M_1-\kappa u^p\Delta M_0\right].
  \label{eq:bawr-pdf-split}
\end{equation}
For a Weibull launch,
\begin{equation}
  f_{\rm BAwR,W}(u)
   =\frac{1}{A}\left[\Delta M_{1,W}-\kappa u^p\Delta M_{0,W}\right].
  \label{eq:bawr-pdf-weibull}
\end{equation}

The frozen branch is the simplest of the three. Inverting
\eqref{eq:bawr-sat}, the saturation time of the start point whose critical
launch is $w$ is $T_{\rm sat}=(w/\kappa)^{1/p}$, and
$\dd z/\dd T_{\rm sat}=-\kappa p T_{\rm sat}^{p}$ while
$\dd w/\dd T_{\rm sat}=\kappa p T_{\rm sat}^{p-1}$, so
\begin{equation}
  \left|\frac{\dd z}{\dd w}\right| = T_{\rm sat}(w)
  = \kappa^{-1/p}\,w^{1/p},
  \qquad
  \Psi_{\rm fr}(u) = \kappa^{-1/p}
    \left[P_{-1/p}(w_a)-P_{-1/p}(w_b)\right],
  \label{eq:bawr-frozen}
\end{equation}
with $w_a=V_c(A)$, $w_b=V_c(Z(u))$ and $P_{-1/p}$ the positive-power primitive
of \eqref{eq:bawd-power-primitives}. As in BAwF this is a positive power, so it
is a truncated moment, closed in $\PhiN,\phiN$ for either ordinary or
split-lognormal launch and in incomplete gamma functions for a Weibull launch,
and generally quadrature for a normal one. There is no logarithmic case: the
exponent $1/p$ is never zero for finite $p$, so the pathology that the
$\rho=1$ drive-decay member introduces has no analogue here.

The non-finishing mass is \eqref{eq:bawd-fmax} with $V_c$ from
\eqref{eq:bawr-sat}, and the point-start value is $P\{V<V_c(0)\}$ with
$V_c(0)=\kappa T_{\max}^{p}$.

\subsection{Endpoint order}
The identity \eqref{eq:bawf-weight} applies verbatim, since $V^{*}$ is affine
in $z$ with Jacobian $u$:
\[
  w_{\rm hi}(u)-\kappa u^{p}
  = \frac{b}{u}-\frac{\kappa p\,u^{p}}{p+1}
  = \frac{z^{*}(u)}{u},
\]
and $z^{*}$ from \eqref{eq:bawr-Z} has a simple zero at $T_{\max}$ for every
$p>0$. The density therefore shuts down linearly with a point start and
quadratically when $A>0$, matching BAwD and BAwF; the shoulder is a property of
the start-point distribution rather than of the decay mechanism.

\subsection{Reductions, limits, and scale}
\begin{itemize}
  \item \textbf{$\kappa\to0$.} Exactly the LBA, for every $p$: the clearance
        vanishes identically, $T_{\max}\to\infty$, $V_c\to0$. This is an
        interior reduction, not an approximation.
  \item \textbf{$p\to0$.} The generator becomes $\dot X=V-\kappa$, an LBA in
        the effective drive $V-\kappa$, and $T_{\max}$ diverges like
        $(\kappa p)^{-1}$. The endpoint and the omission mechanism both
        disappear, and $\kappa$ confounds additively with the launch location
        exactly as $\ell$ does in the drive-decay $k\to0$ limit. The model is
        defined for $p>0$ and this is its degenerate edge.
  \item \textbf{$p\to\infty$ at fixed $T_{\max}$.} Since
        $\kappa/(p+1)=b/(p\,T_{\max}^{p+1})$, the boundary of
        Section~\ref{sec:clocks} is
        $b\left[1+p^{-1}(u/T_{\max})^{p+1}\right]$, which converges to the flat
        boundary $b$ on $[0,T_{\max})$ and diverges beyond it. The limit is an
        LBA with a hard external deadline at $T_{\max}$: responses are
        undistorted up to the deadline and absent after it. This is a
        legitimate limiting model rather than a degeneracy, but it is a flat
        direction in which the likelihood becomes insensitive to $p$ once the
        distortion within the window is negligible.
  \item \textbf{Scale.} As in BAwF, multiplying $(V,b,A,\kappa)$ by a common
        factor leaves every crossing time exactly unchanged, along with
        $T_{\max}$ and $p$, so one of the ordinary BAw* mean $m_L$ (or the
        log-scale location for a split/BTAwL launch), $b$, $A$ and $\kappa$
        must be fixed. Note that $\kappa$ carries evidence units here and can serve
        as the anchor, where BAwF has no such parameter at all.
\end{itemize}

\subsection{Relation to the other variants}
BAwR is not a member of any family above, and the reason is structural rather
than parametric. In BAwD and BAwF the opposing term is $\ell h_\rho^\gamma$ or
$kX$: in both cases it is bounded and eventually vanishing, and the endpoint
arises because the \emph{drive} loses its ability to make progress. In BAwR the
drive never weakens; the endpoint arises because the opposing term overtakes
it. To obtain $\kappa u^p$ from \eqref{eq:bawd-generator} one would need
$h_\rho\equiv1$, which forces $k=0$ and hence a \emph{constant} clearance
$\ell$, giving the $p=0$ degenerate corner rather than any $p>0$ member. Nor is
it BAwL or BAwF: neither the state nor the start point is attenuated, so $z$
contributes undiminished at every $u$ and $\partial X/\partial z=1$ throughout.
Its nearest relative is BAwDp, in that both leave the accumulation itself
untouched and impose a common temporal structure on all trials. But BAwDp does
so through a clock that halts, which keeps the boundary flat and $T_{\max}$
free of $b$, where BAwR does so through a rising boundary in ordinary time,
and that is exactly why its window is not free of $b$.

%======================================================================
\section{Ballistic Transient/Sustained Local Race with Leak}
\label{sec:btawl}
%======================================================================
The full BTAwL race and both component models share one subject-level clock:
apply \eqref{eq:time-warp-btawl-components}--\eqref{eq:time-warp-btawl-race}
after forming the parent component race.
The full BTAwL model is a \emph{local race}, not a
coactivation model.  Each accumulator carries two independent subraces,
transient ($T$) and sustained ($S$), and the response is the first
of the two to reach the common boundary.  The transient and sustained models
expose the two subraces separately, and the full model combines them through
the local-race equations below.

\subsection{Subrace generators and launch allocation}
\label{sec:btawl-origin}
For $j\in\{T,S\}$ let $z_j\sim U(0,A)$ and let $V_j$ be independent of
$z_j$ and of the other subrace.  The two drive profiles are
\begin{equation}
  \eta_T(u)=\frac{u}{\tau_t}\e^{-u/\tau_t},
  \qquad
  \eta_S(u)=1-\e^{-u/\tau_s}.
  \label{eq:btawl-profiles}
\end{equation}
The subrace dynamics are
\begin{equation}
  \dot X_j=V_j\eta_j(u)-kX_j,
  \qquad X_j(0)=z_j,
  \qquad j\in\{T,S\}.
  \label{eq:btawl-local-generator}
\end{equation}
The transient profile is the two-stage low-pass impulse response used by the
Smith detector; the sustained profile is its step response.

The scalar parameter $\pi\in[0,1]$ allocates the base launch law between the
two independent channels.  At an interior value of $\pi$, the allocation is
written in BTAwL's direct log-scale coordinates $(\mu,\sigma)$ for a
lognormal base launch (the ordinary BAw* models use the mean--CV pair above):
\begin{align}
  V_T&\sim\Normal((1-\pi)v,[(1-\pi)s_v]^2),&
  V_S&\sim\Normal(\pi v,(\pi s_v)^2)
  &&\text{(normal)},\\
  \log V_T&\sim\Normal(\mu+\log(1-\pi),\sigma^2),&
  \log V_S&\sim\Normal(\mu+\log\pi,\sigma^2)
  &&\text{(lognormal)},\\
  V_T&\sim\mathcal W((1-\pi)m_W,r),&
  V_S&\sim\mathcal W(\pi m_W,r)
  &&\text{(Weibull)}.
  \label{eq:btawl-allocation}
\end{align}
With channel-specific launch parameters, the Weibull subraces instead use
independent laws $\mathcal W(m_S,r_S)$ and $\mathcal W(m_T,r_T)$.
For a normal launch, positive-drift truncation is applied separately to each
channel.  For a split-lognormal launch, the two
log locations in \eqref{eq:btawl-allocation} are shifted by the same amounts
and the split-shape parameter $\delta$ is unchanged.  At $\pi=0$ the
sustained launch is identically zero and the transient launch is the base
law; at $\pi=1$ the converse holds.  Thus the endpoint values are exact
nested submodels, not limits obtained by sending either time constant to
zero.  In particular, the launches are \emph{independent}; there is no
single trialwise $V$ shared by the two channels.

For either subrace write $\overline G_j$ and $g_j$ for the survivor and
density of its allocated launch law.  All response-time formulas below are in
accumulation time $u=T-t_0$.

\subsection{Trajectories and required launch strengths}
Put $E(u)=\e^{-ku}$.  Variation of constants gives
\begin{equation}
  X_j(u)=z_jE(u)+V_jH_j(u),
  \qquad
  H_j(u)=\int_0^u\e^{-k(u-\theta)}\eta_j(\theta)\dd\theta .
  \label{eq:btawl-traj}
\end{equation}
For the transient member, with $d_t=k-1/\tau_t$,
\begin{equation}
  H_T(u)=
  \frac{(d_tu-1)\e^{-u/\tau_t}+\e^{-ku}}{\tau_t d_t^2},
  \qquad
  H_T'(u)=\eta_T(u)-kH_T(u).
  \label{eq:btawl-K}
\end{equation}
The coincident and leak-free values are
\begin{equation}
  H_T(u)=
  \begin{cases}
    \tau_t[1-(1+u/\tau_t)\e^{-u/\tau_t}],& k=0,\\[2pt]
    u^2\e^{-u/\tau_t}/(2\tau_t),& k=1/\tau_t .
  \end{cases}
  \label{eq:btawl-K-special}
\end{equation}
For the sustained member,
\begin{equation}
  H_S(u)=\frac{1-\e^{-ku}}{k}
       -\frac{\e^{-u/\tau_s}-\e^{-ku}}{k-1/\tau_s},
  \qquad
  H_S'(u)=\eta_S(u)-kH_S(u),
  \label{eq:btawl-sustained}
\end{equation}
with $H_S(u)=u-\tau_s(1-\e^{-u/\tau_s})$ at $k=0$ and
$H_S(u)=(1-\e^{-ku})/k-u\e^{-ku}$ at $k=1/\tau_s$.

For either channel, the launch required to hit at $(u,z)$ is
\begin{equation}
  W_j(u,z)=\frac{b-zE(u)}{H_j(u)}.
  \label{eq:btawl-vstar}
\end{equation}
The start point enters affinely, so every live start-point integral is a
stop-loss or normal-probit difference.  For the transient channel,
\begin{equation}
  \partial_u W_T(u,z)
   =\frac{zkE(u)H_T(u)-[b-zE(u)]H_T'(u)}{H_T(u)^2}.
  \label{eq:btawl-vstar-derivative}
\end{equation}
The sustained kernel has $H_S'(u)>0$ for finite $u$, hence
$W_S$ decreases monotonically and has no finite endpoint.  If $k>0$,
$H_S(u)\to1/k$ and $W_S(u,z)\to kb$; if $k=0$, $H_S(u)\to\infty$ and
$W_S(u,z)\to0$.

\subsection{Transient endpoint and frozen start points}
\label{sec:btawl-endpoint}
For $k>0$, $H_T$ has one maximum after the input peak.  The transient endpoint
is
\begin{equation}
  H_T'(T_{\max})=0,
  \qquad T_{\max}=\tau_t\,\psi(k\tau_t),
  \label{eq:btawl-tmax}
\end{equation}
and is independent of $b$ and $A$.  The endpoint chart takes this
$T_{\max}$ as its coordinate and solves this equation for $\tau_t$; the rate
chart uses $\tau_t$ directly.  At $k=0$, $T_{\max}=\infty$.  In the full local
race, the transient endpoint remains a component property, while the full-race
endpoint is $\infty$ whenever $\pi>0$ because of the sustained subrace.

The live/frozen seam is obtained from $\partial_uW_T=0$, equivalently
\begin{equation}
  z_t(u)=\frac{bH_T'(u)}{E(u)\eta_T(u)}
  \quad\text{because}\quad kH_T(u)+H_T'(u)=\eta_T(u).
  \label{eq:btawl-tangency}
\end{equation}
The equivalent numerator equation
$zkE H_T=[b-zE]H_T'$ is solved on the descending branch
$u\in[\tau_t,T_{\max}]$ and clips $z_t$ to $[0,A]$.  Let
$T_{\rm sat}(z)$ denote that root for a start point that freezes and define
\begin{equation}
  V_{c,T}(z)=W_T(T_{\rm sat}(z),z)
            =\frac{kb}{\eta_T(T_{\rm sat}(z))}.
  \label{eq:btawl-crit}
\end{equation}
The equality follows by substituting the tangency condition into
\eqref{eq:btawl-vstar}.  Since $\eta_T\le\e^{-1}$, every transient omission
threshold on this frozen branch is at least $\e kb$.

For $A>0$, put
\[
  Z_T(u)=\min\{A,\max\{0,z_t(u)\}\}.
\]
The transient component CDF is
\begin{equation}
  F_T(u)=\frac1A\left[
    \int_0^{Z_T(u)}\overline G_T(W_T(u,z))\dd z+
    \int_{Z_T(u)}^A\overline G_T(V_{c,T}(z))\dd z
  \right],
  \label{eq:btawl-cdf}
\end{equation}
where the second integral is present only after the seam has entered the
start-point interval.  For $A=0$, use
$F_T(u)=\overline G_T(b/H_T(u))$ before $T_{\max}$ and
$F_T(u)=\overline G_T(V_{c,T}(0))$ at and after the endpoint.  The transient
density on the live branch is
\begin{equation}
  f_T(u)=\frac1A\int_0^{Z_T(u)}
     g_T(W_T(u,z))[-\partial_uW_T(u,z)]\dd z ,
  \label{eq:btawl-pdf}
\end{equation}
with no contribution from the frozen integral.  At $A=0$ this becomes
$g_T(b/H_T)\,bH_T'/H_T^2$.  The live formulas close with the normal,
lognormal, split-lognormal, and Weibull partial moments; the frozen integral is
the one-dimensional quadrature for all four launch laws.
The Lambert representation
\[
  T_{\rm sat}=-\tau_tW_{-1}\!\left(-\frac{kb}{V_{c,T}}\right)
\]
is useful for evaluating the quadrature, but it does not turn its measure into
an additive or power-law primitive.

\subsection{The sustained component}
For $A>0$ the sustained CDF and density are
\begin{align}
  F_S(u)&=\frac1A\int_0^A\overline G_S(W_S(u,z))\dd z,
  \label{eq:btawl-sustained-cdf}\\
  f_S(u)&=\frac1A\int_0^A
      g_S(W_S(u,z))[-\partial_uW_S(u,z)]\dd z .
  \label{eq:btawl-sustained-pdf}
\end{align}
The point-start versions are obtained by evaluating the integrands at $z=0$.
There is no frozen branch.  Its eventual response mass is
\begin{equation}
  F_S(\infty)=
  \begin{cases}
    \overline G_S(kb),&k>0,\\
    \overline G_S(0),&k=0.
  \end{cases}
  \label{eq:btawl-sustained-limit}
\end{equation}

\subsection{The local race}
The two subraces are independent conditional on their row parameters, so the
full BTAwL distribution is not a mixed trajectory.  Its survivor, CDF, and
density are
\begin{align}
  S_{\mathrm{BTAwL}}(u)
    &=S_T(u)S_S(u),\\
  F_{\mathrm{BTAwL}}(u)
    &=1-[1-F_T(u)][1-F_S(u)],\\
  f_{\mathrm{BTAwL}}(u)
    &=f_T(u)[1-F_S(u)]+f_S(u)[1-F_T(u)].
  \label{eq:btawl-local-race}
\end{align}
The two terms in the density are the transient and sustained winning
subdensities.  At $\pi=0$, these equations reduce exactly to the transient
component; at $\pi=1$, they reduce exactly to the sustained component.  For every
interior $\pi>0$ with $k>0$, the sustained subrace removes the transient hard endpoint:
the full race can respond after $T_{\max}$ whenever the sustained channel has
not yet finished.  Its omission probability is
\begin{equation}
  1-F_{\mathrm{BTAwL}}(\infty)
  =[1-F_T(\infty)][1-F_S(\infty)].
  \label{eq:btawl-local-defect}
\end{equation}
Thus a finite endpoint is a property of the pure transient submodel, not of
the full BTAwL local race.

\subsection{Reductions, charts, and scale}
At $k=0$ the transient trajectory is $X_T=z_T+V_TH_T(u)$, with
$H_T(u)\uparrow\tau_t$ and
\begin{equation}
  V_{c,T}(z)=\frac{b-z}{\tau_t}.
  \label{eq:btawl-k0}
\end{equation}
This is an LBA in a clock that saturates at $\tau_t$ and remains defective
because the transient delivers only a finite total impulse.  The sustained
member at $k=0$ has an unbounded clock $H_S(u)$ and reduces to an LBA in that
clock, so positive launches respond eventually.  The full local race inherits
the product survival in \eqref{eq:btawl-local-defect}.

Multiplying all launch strengths, $b$, and $A$ by a common positive constant
leaves every crossing time unchanged.  The shape parameters $\sigma$ and
$\delta$ (or Weibull shape $r$), the allocation $\pi$, and the time parameters
$k,\tau_t,\tau_s$ are
dimensionless or time-scale parameters and are not evidence-scale anchors.
One location/scale quantity must be fixed, such as $v$, $s_v$, $B$, a log-mean
intercept for an ordinary BAw* lognormal, an intercept of $\mu$ for a
split-lognormal or BTAwL lognormal, or an intercept of $\log m_W$. At a fixed
$\pi$, the inactive channel's time
constant is inert and must be held fixed.

\subsection{Relation to the two-channel parent model}
The profiles in \eqref{eq:btawl-profiles} are sustained and transient
low-pass profiles.
The model retains the race architecture: the two channels have independent
launches and start points and race locally.  It is distinct from the
alternative coactivation construction
$X=z+V[\pi H_S+(1-\pi)H_T]$; that shared-drive trajectory is not BTAwL.

\section{Reductions and Limits}
\label{sec:reductions}
%======================================================================
\subsection{The Degenerate Case: No Dynamics At All}
Set $k = 0$ and $A = 0$ (and, for the drive-decay model, $\ell = 0$). Then
\eqref{eq:bawl-vstar} and \eqref{eq:bawd-vstar} both collapse to
\eqref{eq:lba-vstar} with $z = 0$, so the crossing time is
\begin{equation}
  u = \frac{b}{V}.
  \label{eq:degenerate}
\end{equation}
Every accumulation mechanism has been switched off. The response-time law is
the pushforward of the launch law through a fixed, deterministic, strictly
decreasing map, and the model retains exactly one random variable per trial.
This is the sense in which the models reduce to their drift distribution plus a
constant: the constant is $t_0$, and the shape is whatever the launch law
becomes under $V \mapsto b/V$.

What that shape is depends on the launch law, and the four cases are not
symmetric.

\begin{itemize}
  \item \textbf{Lognormal launch.} The lognormal family is closed under
        reciprocals and under multiplication by a constant, so
        \begin{equation}
          \log u = \log b - \log V \sim \Normal\!\left(\log b - \mu,\ \sigma^2\right),
        \end{equation}
        and $T = t_0 + \mathrm{Lognormal}(\log b - \mu, \sigma^2)$ exactly. The
        reduced model is the shifted lognormal race model, with median
        $t_0 + b\e^{-\mu}$. For an ordinary BAw* launch, the same median is
        $t_0+b\sqrt{1+c_L^2}/m_L$ after applying
        \eqref{eq:lognormal-mean-cv}. Note also that $b$ and the log-scale
        location enter only through $\log b - \mu$ (equivalently $b/m_L$ at
        fixed $c_L$), so the boundary is not separately identified in this
        limit.
  \item \textbf{Split-lognormal launch.} Reflection about the median preserves
        the split-normal form on the log scale. Thus $\log u=\log b-Y$ has
        median $\log b-\mu$, split point $\log b-c$, left/right widths
        $(s_R,s_L)$, and asymmetry $-\delta$. The reduced response is a
        shifted split-lognormal law, with median $t_0+b\e^{-\mu}$.
  \item \textbf{Weibull launch.} With the parameterisation in
        \eqref{eq:weibull-launch}, $u=b/V$ has the reciprocal-Weibull
        pushforward
        \begin{equation}
          F_T(t)=\exp\!\left[-\left(\frac{b}{(t-t_0)\theta_W}\right)^r\right],
          \qquad t>t_0,
        \end{equation}
        and density $g_W(b/u)b/u^2$. The response law is closed in elementary
        Weibull functions but is not itself Weibull on the time scale.
  \item \textbf{Normal launch, truncated to $V>0$.} The truncated normal family
        is \emph{not} closed under reciprocals. The exact statement is that
        $1/(T - t_0)$ is a scaled positively truncated normal, that is,
        $b/(T-t_0) \sim \Normal(v, s_v^2)$ conditioned on being positive. The
        response-time density is then
        \begin{equation}
          f_T(t) = \frac{1}{Z_+}\cdot\frac{b}{(t-t_0)^2 s_v}\,
                   \phiN\!\left(\frac{b/(t-t_0) - v}{s_v}\right),
        \end{equation}
        a reciprocal truncated normal rather than a truncated normal. It has
        the right-hand tail $f_T(t) \sim b\,\phiN(v/s_v)/(Z_+ s_v t^2)$, a
        power law, which no truncated normal has.
\end{itemize}

\begin{remark}
The general form of this reduction matters more than these instances,
because it is the reason the ballistic family is analytically tractable at all.
With $A=0$, every sign-preserving model in this note has a closed-form
crossing time $u=h(V)$ for a model-specific monotone $h$, so its response-time
law is a pushforward and never requires an integral. A sign-changing waveform
with an untruncated launch, such as BAwDp for $0<\lambda<1$, must be treated
piecewise: positive launches can be omitted or cross on the early branch,
whereas negative launches cross on the late branch. Start-point variability is
what turns a pushforward into an integral, and the rest of this document is
about when that integral closes.
\end{remark}

\subsection{Model Limits}
\begin{itemize}
  \item \textbf{BAwL, $k \to 0$.} Exactly the LBA, in every expression, with
        $C_1 \to b/u$, $C_2 \to 1/u$, and $J \to 1/u^2$. The defect
        \eqref{eq:bawl-mass} becomes $P(V > 0)$, so an untruncated normal
        launch stays defective by $\PhiN(-v/s_v)$ while a truncated normal,
        ordinary lognormal, split-lognormal, or Weibull launch becomes proper.
  \item \textbf{BAwD, $k \to 0$.} $h_\rho\to1$,
        $Q_\rho,R_{\rho,\gamma}\to u$, and
        $V^*=(b-z)/u+\ell$ for every $\gamma$. Thus the family becomes an
        LBA in the effective drive $V-\ell$, with no saturation or endpoint.
  \item \textbf{BAwD, $\ell\to0$.} The clock is $Q_\rho$ and the critical
        launch is $(b-z)/Q_\rho(\infty)$ when that limit is finite (and zero
        otherwise). The endpoint disappears; finite $\rho$ and the
        exponential limit differ through the total available clock.
  \item \textbf{BAwD, $\gamma\to1$.} Clearance is proportional to the base
        kernel, $R_{\rho,1}=Q_\rho$, so the model has no finite endpoint and
        its critical launch is \eqref{eq:bawd-matched-crit}.
  \item \textbf{BAwD, $\rho\to\infty$.} The finite-rho kernels converge to the
        exponential kernels and all formulas above converge to the
        exponential member, including the constant-clearance corner.
  \item \textbf{BAwD, $\ell=0$ and $A=0$.} In the exponential corner this
        coincides with the point-start leak model; finite-rho members have
        the same pushforward structure but a different clock. With $A>0$,
        the drive-decay critical launch retains its dependence on $z$.
  \item \textbf{BAwDp, $\lambda \to 0$.} On the positive-launch branch, the
        proportional-clearance model approaches the pure exponential drive
        clock, with $u_{*}\to\infty$, $q_\lambda^{\max}\to1/k$, and
        $V_{\mathrm{crit}}(z)\to k(b-z)$. This is the
        $\rho=\infty,\ell=0$ member of BAwD, not every finite-rho member. For
        an untruncated normal, the negative-launch branch recedes to infinite
        time, so the eventual defect has a nonuniform $\lambda\downarrow0$
        limit.
  \item \textbf{BAwDp, $k \to 0$.} For $\lambda<1$,
        $q_{\lambda}(u) \to (1-\lambda)u$, so the model is an LBA in the
        effective drive $(1-\lambda)V$, with $u_{*}\to\infty$. Only the
        product is identified: $\lambda$ confounds exactly with the launch
        scale. For split-lognormal launches it shifts
        $\mu \mapsto \mu+\log(1-\lambda)$ and the split point $c$ by the
        same amount, while leaving $\sigma,\delta$ unchanged. BTAwL's
        ordinary lognormal launch uses the same direct $\mu$ shift with
        $\sigma$ unchanged. For an ordinary BAw* lognormal launch it is equivalently
        $m_L\mapsto(1-\lambda)m_L$ with $c_L$ unchanged. The same
        multiplicative scaling applies to the Weibull mean
        $m_W\mapsto(1-\lambda)m_W$,
        $(v,s_v)\mapsto(1-\lambda)(v,s_v)$ under a normal one. Note
        that this is a \emph{multiplicative} confound, where the drive-decay
        model's $k\to0$ limit gives the additive confound $V-\ell$.
  \item \textbf{BAwDp, $\lambda \to 1$.} For positive launches, the positive
        part of the waveform disappears, $u_{*}\to0$ and
        $q_{\lambda}^{\max}\to0$, so no positive-launch trial responds. The
        model is defined on $0\le\lambda<1$ and this is the degenerate edge of
        the positive-launch branch. With an untruncated normal, negative
        launches retain a sign-reversing response branch, so the eventual
        response mass does not vanish; the limit is nonuniform.
  \item \textbf{BAwF, $k \to 0$.} $h_\rho\to1$ and $H_\rho\to1$, so
        \eqref{eq:bawf-traj} is the LBA exactly, with $T_{\max}\to\infty$ and
        $V_c\to0$. The reduction is uniform in $\rho$. There is no separate
        ``no clearance'' limit to take, because there is no independent
        clearance to remove: BAwF has a single decay parameter.
  \item \textbf{BAwF, $\rho \to \infty$.} The finite-$\rho$ kernel and its
        derivatives converge to $\e^{x}$, $x_{\max}\to1$, and the frozen
        power primitive $P_{-\beta}$ with $\beta=1/(\rho-1)\to0$ degenerates
        into the logarithmic primitive $L$, matching \eqref{eq:bawf-frozen}.
  \item \textbf{BAwF, $\rho \to 1$.} $G_1\equiv1$, no start point freezes, and
        $V^{*}$ falls monotonically to $kb$. The endpoint disappears and the
        omission threshold becomes free of $z$, coinciding in value with
        the leak model's $kb$ despite a different trajectory. This is a
        different model, and the family is defined on $\rho>1$.
  \item \textbf{BTAwL transient, $k \to 0$.} The transient trajectory
        separates, $X_T=z_T+V_TH_T(u)$, so it is an LBA in a clock that
        saturates at $\tau_t$.  The endpoint recedes to infinity but the
        threshold \eqref{eq:btawl-k0} remains positive and affine in $z_T$;
        finite total transient impulse therefore leaves a defect for positive
        launch laws.
  \item \textbf{BTAwL transient, $k\tau_t \to \infty$.} The transient window
        collapses onto its input peak and the accumulator becomes a follower of
        $V_T\eta_T(u)/k$, a level test on the filtered input.
  \item \textbf{BTAwL sustained, $k \to 0$.} The sustained clock $H_S(u)$ is
        unbounded, so the sustained subrace reduces to an LBA in that clock
        and positive launches respond eventually.
  \item \textbf{BTAwL local race, $\pi$.} At $\pi=0$ the full kernel is the
        pure transient component; at $\pi=1$ it is the pure sustained component.
        For $0<\pi<1$, the CDF and survivor are the product formulas
        \eqref{eq:btawl-local-race}; the sustained subrace removes the
        transient hard endpoint.
  \item \textbf{BAwR, $\kappa \to 0$.} Exactly the LBA, for every $p$.
  \item \textbf{BAwR, $p \to 0$.} An LBA in the effective drive $V-\kappa$,
        with $T_{\max}\to\infty$. This is the same \emph{additive} confound
        that the drive-decay model produces at $k\to0$, and for the same
        reason: a clearance that is constant in $u$ is indistinguishable from
        a shift of the launch location. The model is defined on $p>0$.
  \item \textbf{BAwR, $p \to \infty$ at fixed $T_{\max}$.} An LBA truncated by
        a hard external deadline at $T_{\max}$: undistorted before it, absent
        after it. Unlike the two limits above this is not a reduction to a
        member of the family, and it is the direction along which the
        likelihood becomes least sensitive to $p$.
  \item \textbf{Any model, $A \to 0$.} For sign-preserving launch laws, the
        pushforward form of Section~\ref{sec:reductions} applies: all
        start-point integrals disappear and every distribution function is a
        single survivor evaluation. The untruncated-normal BAwDp law instead
        retains the positive and late negative branches in
        \eqref{eq:bawdp-normal-full}.
  \item \textbf{Any model, $s_v \to 0$ or $\sigma \to 0$.} The launch strength
        is deterministic, so with $A = 0$ the response time is deterministic
        and with $A > 0$ it is the pushforward of a uniform start point through
        the model-specific crossing map; it is not generally uniform in time.
        These limits are degenerate and are noted only because they mark the
        edge of the parameter space.
\end{itemize}

\subsection{Summary of Reductions}
\begin{center}
\small
\begin{tabular}{L{0.075\textwidth}L{0.185\textwidth}L{0.27\textwidth}L{0.29\textwidth}}
\toprule
Model & Reduces to the LBA when & Other named limits & Confound in the LBA
corner \\
\midrule
BAwL & $k=0$ & --- & none; the reduction is exact \\
\addlinespace
BAwD & $k=0$ and $\ell=0$, i.e.\ $T_{\max}=\infty$ &
$\ell=0$: drive clock $Q_\rho$, no endpoint. $\gamma=1$: matched decay, no
endpoint, threshold \eqref{eq:bawd-matched-crit}. $\rho\to\infty$: the
exponential member. &
$k=0$ alone gives an LBA in $V-\ell$: additive, so $\ell$ confounds with the
launch location \\
\addlinespace
BAwDp & $\lambda=0$ and $k=0$ &
$\lambda=0$: the $\rho=\infty,\ \ell=0$ drive-decay member. $\lambda\to1$: no
positive-launch trial responds &
$k=0$ alone gives an LBA in $(1-\lambda)V$: multiplicative, so $\lambda$
confounds with the launch \emph{scale} \\
\addlinespace
BAwF & $k=0$ &
$\rho\to\infty$: the exponential member. $\rho=1$: no endpoint, threshold
$kb$ as in BAwL &
none; there is no second decay parameter to confound \\
\addlinespace
BAwR & $\kappa=0$ &
$p\to0$: LBA in $V-\kappa$. $p\to\infty$ at fixed $T_{\max}$: LBA with a hard
deadline &
$p\to0$ gives the same additive confound as BAwD's $k\to0$ \\
\addlinespace
BTAwL transient & never exactly; $k=0$ gives an LBA in the saturating clock
$H_T$, which stays defective &
$k=0$: threshold $(b-z)/\tau_t$, no endpoint.
$k\tau_t\to\infty$: a level test on the filtered input &
$k\tau_t\to\infty$ leaves only $kb$ identified \\
\addlinespace
BTAwL sustained & $k=0$ gives an LBA in the unbounded clock $H_S$ &
$k>0$: asymptotic threshold $kb$ and no endpoint &
$k$ and $b$ enter the sustained plateau through $kb$ \\
\addlinespace
BTAwL local race & $\pi=0$ and $\pi=1$ give the two component limits &
$0<\pi<1$: product survivor in \eqref{eq:btawl-local-race}; no hard endpoint &
$\pi$ allocates independent launch scales \\
\bottomrule
\end{tabular}

\smallskip
Two further reductions are common to the family. With $A=0$ and a
sign-preserving launch law, every distribution function is the pushforward
$u=h(V)$ of Section~\ref{sec:reductions}. With $A=0$ and all decay parameters
zero, an ordinary lognormal launch gives a shifted lognormal response time, a
split-lognormal launch gives a shifted split-lognormal response time, a Weibull
launch gives the reciprocal-Weibull pushforward in \eqref{eq:degenerate}, and a
truncated normal launch gives the reciprocal truncated normal of
\eqref{eq:degenerate}. The last of these excludes the transient BTAwL component,
whose crossing map is $u=H_T^{-1}(b/V_T)$ rather than $b/V_T$ even after the
leak is switched off: the transient shape has no ``off'' setting.  The full
BTAwL local race is the minimum of this transient crossing time and the
sustained crossing time, so it is not a single pushforward.
\end{center}

%======================================================================
\section{Defective Components}
\label{sec:defect}
%======================================================================
\subsection{What Makes a Ballistic Law Defective}
For a sign-preserving launch law, a ballistic accumulator fails to respond
exactly when its launch strength falls below the smallest value that would ever
carry it to the boundary. Define
\begin{equation}
  V_{\mathrm{crit}}(z) = \inf_{u > 0} V^{*}(u,z)
  = \lim_{u \to \infty} V^{*}(u,z)
    \ \text{ or }\
    \min_{u} V^{*}(u,z),
\end{equation}
according to whether $V^{*}$ is monotone or U-shaped. Then
\begin{equation}
  F(\infty) = \frac{1}{Z_+}\,E_z\!\left[\Gbar\!\left(
      V_{\mathrm{crit}}(z)\right)\right],
  \qquad
  1 - F(\infty) = \text{defect},
  \label{eq:generic-defect}
\end{equation}
and the entire question of defectiveness reduces to identifying
$V_{\mathrm{crit}}$. A sign-changing waveform with an untruncated launch,
such as proportional clearance, must instead be decomposed into its
sign-specific branches as in \eqref{eq:bawdp-normal-defect}. This is a much
simpler object than the corresponding question for a diffusion accumulator,
where the non-hitting probability comes from the image term of the
first-passage law.

The models place $V_{\mathrm{crit}}$ in seven different places. Take the
conventions requested for each: an \emph{untruncated} normal launch for the
LBA, and a positive launch (ordinary, split-lognormal, or Weibull) with only
the leak or decay mechanism active for the others. Note that only the LBA's is
a property of the launch law alone;
everywhere else the accumulation parameters set the defect.

\subsection{LBA with an Untruncated Normal Launch}
Here $V^{*}(u,z) = (b-z)/u \to 0$, so
\begin{equation}
  V_{\mathrm{crit}}(z) = 0
  \qquad\text{for every } z,
\end{equation}
and
\begin{equation}
  F_{\mathrm{LBA}}(\infty) = P(V > 0) = \PhiN\!\left(\frac{v}{s_v}\right),
  \qquad
  \text{defect} = \PhiN\!\left(-\frac{v}{s_v}\right).
  \label{eq:lba-defect}
\end{equation}
The defect is a \emph{sign failure}. An accumulator with a negative launch
strength travels away from the boundary forever; one with a positive launch
strength always arrives, however slowly. Three features follow.

\begin{itemize}
  \item The defect does not depend on $b$, on $A$, or on the start point. It is
        a property of the launch law alone. Raising the boundary delays every
        response but converts no responder into a non-responder.
  \item The surviving mass has unbounded support with a heavy right tail. Since
        $u = (b-z)/V$ and the density of $V$ is bounded away from zero at
        $V = 0^{+}$,
        \begin{equation}
          f_{\mathrm{LBA}}(u) \;\sim\;
          \frac{\left(b - \tfrac{1}{2}A\right)}{s_v\,u^{2}}\,
          \phiN\!\left(\frac{v}{s_v}\right),
          \qquad u \to \infty,
          \label{eq:lba-tail}
        \end{equation}
        a $u^{-2}$ power law. The mean response time is infinite. This is the
        practical reason positive truncation is usually imposed: truncation
        removes the defect and, at the same time, removes the launch density at
        $V = 0$, which is what generated the power-law tail.
  \item In a race the defect compounds multiplicatively,
        $\prod_j \PhiN(-v_j/s_{v,j})$ by \eqref{eq:race-omission}, so an
        $n$-accumulator LBA with untruncated drifts predicts omissions only
        when every accumulator draws a negative launch strength. With several
        accumulators and moderate drifts this is a very small number, which is
        why the untruncated LBA is a poor model of omissions even though it is
        formally defective.
\end{itemize}

\subsection{Leak Model with a Positive Launch}
Here $V^{*}(u,z) \to kb$ from above, independently of $z$, so
\begin{equation}
  V_{\mathrm{crit}}(z) = kb
  \qquad\text{for every } z,
\end{equation}
and
\begin{equation}
  F_{\mathrm{BAwL}}(\infty) = \PhiN\!\left(\frac{\mu - \log(kb)}{\sigma}\right),
  \qquad
  \text{defect} = \PhiN\!\left(\frac{\log(kb) - \mu}{\sigma}\right).
  \label{eq:bawl-defect}
\end{equation}
For a split-lognormal launch, the same two quantities are
$F_{\mathrm{BAwL}}(\infty)=\Gbar_{\rm sp}(kb)$ and
$\text{defect}=1-\Gbar_{\rm sp}(kb)$, with the survivor in
\eqref{eq:split-survivor-moment}.
For a Weibull launch, they are
$F_{\mathrm{BAwL}}(\infty)=\Gbar_W(kb)$ and
$\text{defect}=1-\Gbar_W(kb)$.
The defect is an \emph{asymptote failure}: the accumulator's steady state $V/k$
lies below the boundary. Because the ordinary lognormal, split-lognormal, and
Weibull launches are positive by construction, no truncation is involved and the defect
is entirely produced by
the leak.

\begin{itemize}
  \item The defect depends on $k$ and $b$ only through the product $kb$, and on
        the launch law through the single survivor evaluation at $kb$. It is again free
        of $A$: the start point leaks away, so it cannot rescue a trial. This
        is the sharpest structural difference from the drive-decay model.
  \item The defect is a smooth, monotone function of $k$ that is zero at
        $k = 0$ and approaches one as $k \to \infty$. Unlike the LBA, this
        model can therefore produce any omission rate at all, and the rate is a
        deterministic consequence of the accumulation parameters rather than a
        free contaminant probability.
  \item The surviving mass has unbounded support with an exponential right
        tail at rate $k$, by \eqref{eq:bawl-tail}. The mean response time is
        finite conditional on responding. This is the opposite behaviour from
        \eqref{eq:lba-tail}: leak both creates omissions and lightens the tail
        of the responses that do occur, because the trials that would have been
        very slow are exactly the ones converted into omissions.
  \item The defect and the tail rate are governed by overlapping parameters,
        $\log(kb)$ and $k$ respectively. Observing both an omission rate and a
        tail shape is what separates $k$ from $b$; either alone does not.
\end{itemize}

\subsection{Drive-Decay Model with a Positive Launch}
For $A>0$ and $V_{\mathrm{crit}}(z)>0$,
\begin{equation}
 F_{\mathrm{BAwD}}(\infty)
 =\frac{1}{A}\int_0^A
 \PhiN\!\left(\frac{\mu-\log V_{\mathrm{crit}}(z)}{\sigma}\right)\dd z,
 \qquad
 V_{\mathrm{crit}}(z)=
 \begin{cases}
 V_c(z),&\gamma<1,\\
 \ell+(b-z)/Q_\rho(\infty),&\gamma=1.
 \end{cases}
 \label{eq:bawd-defect-logn}
\end{equation}
For $A=0$, this is understood as the point-start value
$F_{\mathrm{BAwD}}(\infty)=
\PhiN((\mu-\log V_{\mathrm{crit}}(0))/\sigma)$. The $\ell=0$ cases with
$V_{\mathrm{crit}}=0$ are taken by continuity: either positive log-scale or
Weibull launch has response probability one and defect zero.
For a split-lognormal launch, replace each displayed Gaussian tail by
$\Gbar_{\rm sp}(V_{\mathrm{crit}}(z))$; the start-point integral is unchanged.
For a Weibull launch, the corresponding expression is
\begin{equation}
  F_{\mathrm{BAwD},W}(\infty)
   =\frac{1}{A}\int_0^A \Gbar_W\!\left(V_{\mathrm{crit}}(z)\right)\dd z,
  \label{eq:bawd-defect-weibull}
\end{equation}
with the integral evaluated by the stop-loss and incomplete-gamma primitives
above.
For $0\le\gamma<1$, the defect is a tangency failure: the trajectory peaks
below the boundary and falls away. The support is bounded by
$T_{\max}=T_{\rm sat}(0)$. For $\gamma=1$ there is no finite endpoint; the
defect is an asymptotic failure whenever the total clock is finite, and is the
launch-law mass below $\ell$ when $\rho=1$.
\begin{itemize}
  \item For $\gamma<1$, or for $\gamma=1$ with finite
        $Q_\rho(\infty)$, the defect depends on $(A,b,k,\ell)$, on
        $(\rho,\gamma)$, and on the launch law. Unlike the leak model,
        start-point variability changes it: $V_{\mathrm{crit}}(z)$ is
        decreasing in $z$, so a high start point genuinely rescues trials.
  \item At the matched $\rho=1,\gamma=1$ corner, $V_{\mathrm{crit}}(z)=\ell$
        is independent of $z$, $A$, and $b$; the defect is simply the launch
        mass below $\ell$.
  \item For $\gamma<1$ the surviving mass has bounded support
        $u\in(0,T_{\max}]$, with density vanishing linearly ($A=0$) or
        quadratically ($A>0$) at the endpoint. The saturation equation links
        the omission mass and the response window.
  \item For $\gamma=1$ the support is unbounded. When $1<\rho<\infty$,
        $Q_\rho(\infty)$ is finite and the approach to the limiting threshold
        is power-law; when $\rho=\infty$ it is exponential; when $\rho=1$
        the total clock diverges and only launches below $\ell$ omit.
  \item In a race, \eqref{eq:race-omission} gives the product of the
        component defects. For $\gamma<1$ the race has bounded support at
        the largest component endpoint; matched-decay components instead
        contribute unbounded response times.
\end{itemize}

\subsection{Proportional Clearance for Positive Launches}
Here $V^{*}$ is U-shaped but the freeze is simultaneous, and
$V_{\mathrm{crit}}(z) = (b-z)/q_{\lambda}^{\max}$ from
\eqref{eq:bawdp-crit} is affine in $z$. For $A>0$,
\begin{equation}
  F_{\mathrm{BAwDp}}^{+}(\infty)
  = \frac{1}{A}\int_0^{A}\Gbar\!\left(
      \frac{b-z}{q_{\lambda}^{\max}}\right)\dd z
  = F_{\mathrm{LBA}}\!\left(q_{\lambda}^{\max}\right),
\end{equation}
with the point-start value obtained by setting $z=0$. This is a single
stop-loss difference, closed for all four positive launch laws. The failure mode
is again a tangency failure, but of a degenerate kind: all trajectories peak at
the same instant, so the model asks one question, whether
$Vq_{\lambda}^{\max}$ clears $b-z$, and the positive-launch defect is just the
LBA's own non-finishing mass evaluated at a finite internal time.
For an untruncated normal, the eventual response probability instead includes
the late negative-launch branch and is given by
\eqref{eq:bawdp-normal-defect}; it is not a single frozen-branch evaluation.

\begin{itemize}
  \item Under the positive-launch convention, the defect depends on
        $(A,b,k,\lambda)$ only through
        $q_{\lambda}^{\max}=(1-\lambda+\lambda\log\lambda)/k$ together with
        $(A,b)$. Start-point variability reduces it, as in the drive-decay
        model, because a high start point shortens the distance the clock has to
        cover.
  \item The positive-launch support is bounded, $u\in(0,u_{*}]$, with density
        vanishing linearly at the endpoint for every $A$. Since $u_{*}$ and
        $q_{\lambda}^{\max}$ are different functions of $(k,\lambda)$ --- the
        first is $\log(1/\lambda)/k$ and the second
        $(1-\lambda+\lambda\log\lambda)/k$ --- the window length and omission
        rate are not tied to each other as tightly as the drive-decay model's
        \eqref{eq:bawd-sat} ties them. At fixed $u_{*}$, varying $\lambda$
        still moves the defect.
  \item In a race of positive-launch components the omission mass is
        $\prod_j(1-F_{\mathrm{LBA}}^{(j)}
        (q_{\lambda}^{\max,(j)}))$, and components sharing $(k,\lambda)$ have
        support exactly $(0,u_{*}]$, independent of which accumulator wins.
        The positive-branch endpoint is exactly
        $\log(1/\lambda)/k$.
  \item With an untruncated normal, add the negative-launch defect and late
        branch from \eqref{eq:bawdp-normal-defect}: the support is
        $(0,u_{*}]\cup(u_0,\infty)$, not bounded, and the limit
        $\lambda\downarrow0$ is nonuniform because $u_0\to\infty$.
\end{itemize}

\subsection{Global Fading and Ramping Clearance}
Both new variants have a staggered tangency failure of exactly the BAwD kind:
$V^{*}$ is U-shaped, the freeze time depends on the start point, and
$V_{\mathrm{crit}}(z)=V_c(z)$ is decreasing in $z$, so a high start point
rescues trials that a low one would lose. For $A>0$ the defect is
\eqref{eq:bawd-defect-logn} with $V_c$ taken from \eqref{eq:bawf-crit} or
\eqref{eq:bawr-sat}, and for $A=0$ it is the point-start value
$P\{V<V_c(0)\}$, with
\begin{equation}
  V_c^{\rm BAwF}(0) = k\,b\left(\frac{\rho}{\rho-1}\right)^{\rho-1}
  \ \ \left(=\e kb\ \text{at}\ \rho=\infty\right),
  \qquad
  V_c^{\rm BAwR}(0) = \kappa\,T_{\max}^{p}.
  \label{eq:new-defects}
\end{equation}
What differs between them is how the defect is coupled to the response window.

For BAwF the two are separated. The window is $x_{\max}/k$ and the threshold is
proportional to $kb$, so at a fixed window a change in $b$ moves the
omission rate and nothing else, and at a fixed omission rate a change in $k$
moves the window and nothing else. No other model in this note has that
property: BAwD ties the two through the saturation equation
\eqref{eq:bawd-sat}, BAwR through \eqref{eq:bawr-tmax}, and BAwDp leaves them
loosely coupled through two different functions of $(k,\lambda)$.

For BAwR the coupling is tight and explicit. Substituting
\eqref{eq:bawr-tmax} into \eqref{eq:new-defects},
\begin{equation}
  V_c(0) = \frac{(p+1)\,b}{p\,T_{\max}},
  \label{eq:bawr-defect-window}
\end{equation}
so at a given window and boundary the point-start omission threshold is
determined up to the factor $(p+1)/p$, which ranges over $(1,\infty)$ but is
already within $10\%$ of its limit by $p=10$. Thus $p$ is determined almost
entirely by the shape of the density inside the
window rather than by the omission rate, which is the $p\to\infty$ limit of
Section~\ref{sec:reductions} seen from the defect side.

In a race, \eqref{eq:race-omission} applies unchanged to both, and the support
of the race is bounded by the largest component endpoint. For BAwF that
endpoint is common to any set of accumulators sharing $k$ and $\rho$, whatever
their boundaries; for BAwR it is not, since components with different $b$ have
different windows.

\subsection{BTAwL components}
The transient BTAwL component has the staggered defect
$V_{c,T}(z)=kb/\eta_T(T_{\rm sat}(z))$ from \eqref{eq:btawl-crit}; its
threshold is never below $\e kb$ for $k>0$ and its density ends at
$T_{\max}$.  At $k=0$ it has the affine threshold
\eqref{eq:btawl-k0}, an unbounded support, and a finite-impulse defect.
The sustained component has no finite endpoint and, for $k>0$, its omission
threshold is the start-point-free plateau $kb$.

For the full local race the component defects multiply:
\[
  1-F_{\mathrm{BTAwL}}(\infty)
  =[1-F_T(\infty)][1-F_S(\infty)].
\]
Consequently the interior model has unbounded support even though its
transient component has a hard endpoint.  The endpoint shoulder and omission
mass of the full model therefore cannot be read from the transient component
alone; they depend on $\pi$ and on both allocated launch laws.

\subsection{Comparison}
\begin{center}
\small
\begin{tabular}{L{0.075\textwidth}L{0.125\textwidth}L{0.09\textwidth}L{0.095\textwidth}L{0.10\textwidth}L{0.10\textwidth}L{0.085\textwidth}L{0.06\textwidth}}
\toprule
Model & $V_{\mathrm{crit}}(z)$ & Failure mode & Support & Right tail &
Endpoint order & Defect uses $A,b$? & Cond.\ mean \\
\midrule
LBA$^{\dagger}$ & $0$ & wrong sign & $(0,\infty)$ & $u^{-2}$ & --- &
no & infinite \\
\addlinespace
BAwL & $kb$, free of $z$ & asymptote below $b$ & $(0,\infty)$ &
$\e^{-ku}$ & --- & through $kb$ only & finite \\
\addlinespace
BAwD, $\gamma<1$ & $V_c(z)$, eq.~\eqref{eq:bawd-crit} & staggered peak below
$b$ & $(0,T_{\max}]$ & exact zero & linear at $A=0$, quadratic at $A>0$ &
yes & finite \\
\addlinespace
BAwD, $\gamma=1$ & $\ell+(b-z)/Q_\rho(\infty)$ & asymptote below $b$ &
$(0,\infty)$ & $\rho$-dependent: power-law, exponential or log-clock & --- &
yes, except at $\rho=1$ & param.\ dep. \\
\addlinespace
BAwDp$^{\ddagger}$ & $(b-z)/q_\lambda^{\max}$, affine in $z$ &
peak below $b$, simultaneous & $(0,u_{*}]$ & exact zero & linear at any $A$ &
yes & finite \\
\addlinespace
BAwF & $kbH_\rho'(x_{\rm sat})$ & staggered peak below $b$ &
$(0,x_{\max}/k]$ & exact zero & linear at $A=0$, quadratic at $A>0$ &
yes, but not the endpoint & finite \\
\addlinespace
BAwR & $\kappa T_{\rm sat}(z)^{p}$ & staggered peak below $b$ &
$(0,T_{\max}]$ & exact zero & linear at $A=0$, quadratic at $A>0$ &
yes & finite \\
\addlinespace
BTAwL transient, $k>0$ & $kb/\eta_T(T_{\rm sat}(z))\ge\e kb$ &
staggered peak below $b$ & $(0,T_{\max}]$ & exact zero &
linear at $A=0$, quadratic at $A>0$ & yes, but not the endpoint & finite \\
\addlinespace
BTAwL sustained & $kb$ for $k>0$; $0$ for $k=0$ & asymptotic plateau &
$(0,\infty)$ & no hard endpoint & --- & through $kb$ & finite \\
\addlinespace
BTAwL local race, $0<\pi<1$ & product of component defects &
independent local subraces & $(0,\infty)$ & sustained tail &
not applicable & both components & finite \\
\addlinespace
BTAwL transient, $k=0$ & $(b-z)/\tau_t$, affine in $z$ & drive exhausted &
$(0,\infty)$ & $u\,\e^{-u/\tau_t}$ & --- & yes & finite \\
\bottomrule
\end{tabular}

\smallskip
{\footnotesize $^{\dagger}$ untruncated normal launch. \quad
$^{\ddagger}$ positive launches only.}
\end{center}
The key distinction between these models is the difference between a
tail and a hard endpoint.  The full local race is unbounded whenever its
sustained allocation is positive, while the pure transient component retains the
finite endpoint.  The other finite-endpoint variants have exact zero density
after their respective windows.

Among the finite-endpoint models, the endpoint shoulder separates BAwDp
(linear at any $A$, since its freeze is simultaneous) from the staggered
models (quadratic once $A>0$).  The transient BTAwL component is one of those
staggered models; the full local race has no endpoint for $0<\pi<1$.

For an untruncated normal, the proportional-clearance late negative-launch
branch has unbounded support and an infinite mean, and the untruncated-normal
LBA is defective by $\PhiN(-v/s_v)$ for reasons that have nothing to do with
decay. None of BAwF, BAwR and either BTAwL subrace has a sign-changing
waveform, so their defects are governed by \eqref{eq:generic-defect}, with the
full local race combining them through \eqref{eq:btawl-local-defect}.

%======================================================================
\section{Where the Launch Laws Differ}
\label{sec:closure}
%======================================================================
Affine start-point geometry produces the additive primitive $I_1$, while
\emph{staggered} frozen saturation can induce power-law or logarithmic
measures:
\begin{equation}
 I_1=\int \Gbar(w)\dd w,\qquad
 I_2=\int \Gbar(w)\frac{\dd w}{w}.
 \label{eq:two-integrals}
\end{equation}
The generalized frozen branch also uses
\begin{equation}
 I_\alpha(a,b)=\int_a^b w^{-\alpha}\Gbar(w)\dd w,\qquad
 I_{\log}(a,b;\ell)=\int_a^b\log(w/\ell)\Gbar(w)\dd w.
 \label{eq:bawd-closure-primitives}
\end{equation}
Here $I_2$ is the special case $\alpha=1$. For the finite-$\rho$ frozen
BAwD branch with $0\le\gamma<1$,
$\alpha=(\rho-1)/[\rho(1-\gamma)]$, while $\rho=1$ gives the logarithmic
limit. Which primitive appears is decided by the geometry of the start-point
coordinate, not by the launch law.

\begin{itemize}
  \item Whenever $V^{*}$ is \emph{affine in $z$}, the substitution $w = V^{*}$
        has constant Jacobian, so $\dd z \propto \dd w$ and only $I_1$ appears.
        This covers the LBA, the whole leak model, and the live branch of every
        decaying model in this note, since all of them have a $V^{*}$ that is
        affine in $z$.
  \item On the frozen branch of the drive-decay family the coordinate is
        $T_{\rm sat}$ and the critical launch is
        $V_c(z)=\ell h_\rho(T_{\rm sat})^{\gamma-1}$. Its induced measure
        is model-dependent; in the exponential constant-clearance corner it
        includes $\dd w/w$, producing the $I_2$ obstruction.
  \item The positive-launch proportional-clearance branch also has a frozen
        branch, but every start point freezes at the same instant $u_{*}$, so
        the frozen critical launch $(b-z)/q_{\lambda}^{\max}$ is still affine
        in $z$. The saturation exponent is never needed as a coordinate, and
        only $I_1$ appears. Saturation as such is not what breaks closure;
        \emph{staggered} saturation is. Its untruncated-normal late branch is
        instead handled by the normal partial integrals in
        \eqref{eq:bawdp-negative-cdf}.
  \item The frozen branches of global fading and ramping clearance are
        staggered, so both need a power coordinate, and both take
        $\alpha<0$: $\alpha=-1/(\rho-1)$ from
        \eqref{eq:bawf-frozen-jacobian} and $\alpha=-1/p$ from
        \eqref{eq:bawr-frozen}. Negative $\alpha$ makes no difference to
        closure. $I_\alpha$ is a truncated ordinary- or split-lognormal
        moment for any real $\alpha$ and closes; for a Weibull launch it is the
        incomplete-gamma primitive $P_{\alpha,W}$; for a normal launch it closes for no $\alpha$
        except the trivial one. Global fading recovers $I_{\log}$ in its
        exponential limit $\rho\to\infty$, where $\alpha\to0$; ramping
        decay never does, since $1/p$ is bounded away from zero. For a
        split-lognormal launch, the same power primitives are the explicit
        $P_\alpha^{\rm sp}$ and $L^{\rm sp}$ in
        \eqref{eq:split-power-primitives}; for a Weibull launch they are
        $P_{\alpha,W}$ and $L_W$ in \eqref{eq:weibull-power-primitives}.
\end{itemize}

$I_1$ is the stop-loss transform \eqref{eq:stoploss}, and it is closed in
$\PhiN$ and $\phiN$ for any launch law with a closed-form partial expectation.
The normal, ordinary lognormal, split-lognormal, and Weibull launch laws all
qualify (the last through incomplete gamma functions). Thus the
launch law is irrelevant to closure for the LBA, the leak model, and every
live branch. On a staggered frozen branch (BAwD with $\gamma<1$, BAwF, or
BAwR), the ordinary and split-lognormal launches close through the generalized
power and logarithmic primitives, and the Weibull launch closes through
incomplete-gamma primitives, whereas the normal generally requires
one-dimensional quadrature. The matched-decay $\gamma=1$ branch and the proportional-clearance
branch have no staggered frozen integral at all.


\begin{remark}
The contrast is narrower than is often stated. With leak, a truncated normal
launch has closed-form start-point averaging just as an ordinary,
split-lognormal, or Weibull one does;
equations \eqref{eq:bawl-cdf} and \eqref{eq:bawl-cdf-logn} are exact. The
launch law decides closure on the staggered frozen branch of the drive-decay
family, including its non-finishing mass
\eqref{eq:bawd-fmax} or \eqref{eq:bawd-defect-logn}; the live and matched
$\gamma=1$ branches remain closed for all four launch laws.
\end{remark}

\subsection{Additive Versus Multiplicative Geometry}
\begin{proposition}[informal]
$I_2$ is closed in $\PhiN$ and $\phiN$ whenever the launch survivor is a
fixed (possibly piecewise) Gaussian-tail function of $\log w$. This includes
the ordinary lognormal and the split-lognormal for fixed $\delta$; both are
scale families generated by location families on the log scale.
\end{proposition}

The reason is a change of measure. Write $x = (\log w - \mu)/\sigma$, so that
$\dd w / w = \sigma \dd x$. The measure $\dd w / w$ is the invariant (Haar)
measure of the multiplicative group on $(0,\infty)$, and $x$ is the coordinate
in which that group acts by translation. So $I_2$ is an integral of the
survivor against \emph{translation-invariant} measure in $\log w$.

For the lognormal, the survivor in that coordinate is
\begin{equation}
  \Gbar(w) = \Phibar(x),
\end{equation}
a fixed function of $x$ with no residual dependence on $w$. Define
\begin{equation}
 D(x)=\phiN(x)-x\Phibar(x),\qquad D'(x)=-\Phibar(x).
 \label{eq:Dfun}
\end{equation}
Hence
\begin{equation}
 I_2=\sigma\int\Phibar(x)\dd x=-\sigma D(x)+\text{const}.
\end{equation}
The lognormal is a location family in $\log w$ and $I_2$ integrates over
$\log w$; the two match, and the integral is the same object as $I_1$ viewed
in the other coordinate.

For the split-lognormal, use $x=(\log w-c)/s_L$ on the left of the split and
$x=(\log w-c)/s_R$ on the right. The survivor is the corresponding piecewise
Gaussian-tail expression from \eqref{eq:split-tail-moments}, and integrating
over $\dd w/w=\dd(\log w)$ gives endpoint terms in $\PhiN$ and $\phiN$ on
each side. The asymmetry parameter $\delta$ changes the two pieces and the
location of the join but does not introduce a new quadrature.

For a Weibull launch, put $x=(w/\theta_W)^r$. Then
$\dd w/w=\dd x/(rx)$ and the survivor is $\e^{-x}$, so on a finite interval
\begin{equation}
  I_{2,W}(a,b)
   =\frac{1}{r}\left[\Gamma\!\left(0,(a/\theta_W)^r\right)
                    -\Gamma\!\left(0,(b/\theta_W)^r\right)\right].
  \label{eq:weibull-i2}
\end{equation}
Thus the multiplicative frozen coordinate is also closed for Weibull launches,
through the upper incomplete gamma function of order zero.

For the normal, the survivor in the same coordinate is
\begin{equation}
 \Gbar(w)=
 \PhiN\!\left(\frac{v-\e^{\sigma x+\mu}}{s_v}\right).
\end{equation}
A Gaussian tail is now composed with an exponential map. The normal is a
location family in $w$, while $I_2$ integrates over $\log w$. Concretely, its
finite-interval form is
\begin{equation}
 I_2^{\rm N}(a,b)
 =\int_a^b\PhiN\!\left(\frac{v-w}{s_v}\right)\frac{\dd w}{w},
 \label{eq:bad-integral}
\end{equation}
which is not expressible in error functions and elementary operations. It
belongs to the family of exponential integrals, and no reparameterization of
the model removes it, because the mismatch is between the group structure of
the launch law and the group structure of the coordinate imposed by the
saturation equation.

\subsection{Where the Multiplicative Coordinate Comes From}
The multiplicative coordinate is not a choice. It is forced by the saturation
equation \eqref{eq:bawd-sat}. For finite $\rho$ this is a power-law
coordinate, while $\rho=\infty,\gamma=0$ gives $\e^y-1-y$ and the
multiplicative critical launch $V_c=\ell\e^{ku}$. Thus the familiar
$\dd w/w$ obstruction is a special corner, not a claim about every member.

Three corollaries follow.
\begin{itemize}
  \item The density \eqref{eq:bawd-live} differentiates the live branch only,
        so it is closed for all four launch laws. Distribution and omission
        integrals can require the induced frozen-coordinate measure.
  \item The obstruction does not appear in the leak model at any parameter
        setting, because leak never produces a frozen branch: $V^{*}$ is
        monotone in $u$ for all $k$, so the multiplicative coordinate never
        arises.
  \item The obstruction is absent from the positive-launch proportional-
        clearance branch. With $\ell=\lambda V$ the tangency condition loses
        its dependence on $V$ altogether, so there is no per-trial saturation
        time to parameterize. Proportional clearance has no staggered
        saturation coordinate, and its positive-branch distribution closes
        directly. An untruncated normal adds the sign-specific late branch,
        whose normal partial integrals remain closed but whose support is
        unbounded.
  \item At $A=0$ there is no start-point integral, so the frozen-coordinate
        obstruction does not arise: the surviving quantity
        $\Gbar(V_c(0))$ is a single survivor evaluation and is trivially closed
        for all four launch laws. The obstruction requires both start-point
        variability and staggered saturation.
\end{itemize}

\subsection{A Second Property of Positive Launches}
Two further properties of the positive launches are structural and separate
from the closure argument.

First, $V > 0$ holds by construction for ordinary, split-lognormal, and
Weibull launches, so no truncation normalizer exists and
none of the ambiguity about what a truncated parameter means arises. Under a
truncated normal launch, $v$ and $s_v$ describe the untruncated law, and the
mean and variance of the launch strength actually used are neither of them.

Second, the log-scale and Weibull launches are closed under the scale
transformation that the model itself is invariant to. The evidence axis is defined only up to a
scale: the
map $(V, \ell, b, A) \mapsto (cV, c\ell, cb, cA)$ leaves
\eqref{eq:bawd-traj} and every crossing time unchanged, so exactly one of these
quantities must be fixed. For an ordinary BAw* lognormal launch, rescaling $V$
by $q>0$ maps $m_L\mapsto q m_L$ and leaves $c_L$ unchanged. For a
split-lognormal or BTAwL lognormal launch, it maps $\mu\mapsto\mu+\log q$
(and $c\mapsto c+\log q$ for the split form), while leaving the shape
coordinates unchanged. For a Weibull launch, the same rescaling maps
$m_W\mapsto q m_W$ and leaves $r$ unchanged. Consequently $\sigma$, $\delta$,
and $r$ are dimensionless, and the usual $s_v = 1$ convention has no
positive-launch analogue: fixing a shape parameter fixes nothing. In the endpoint chart the anchor
is a log-mean intercept for an ordinary BAw* launch, an intercept of $\mu$ for
a split/BTAwL lognormal launch, $\log m_W$, $B$ or $A$; in the rate chart it is $\ell=1$, after
which $V$ and $kb$ read as dimensionless. See \S\ref{sec:identifiability}.

%======================================================================
\section{Decay Mechanisms: Clocks and Boundaries}
\label{sec:clocks}
%======================================================================
All six decaying models are ``an accumulator that decays''. What separates
them needs to be
stated precisely, because each can be rewritten as an ordinary LBA
racing a \emph{moving boundary} in a \emph{warped clock}. The mechanistic
content of each model is therefore not the presence of decay but the pair
(clock, boundary) it induces, and, across conditions, which of those objects a
manipulation is allowed to change.

The six fall into three groups under this description. Leak and the transient
drive warp the clock \emph{and} rescale the state, which is why they are the
models that forget what they have already accumulated. Drive decay and
proportional clearance warp the clock alone. Global fading and ramping
clearance do the opposite: they leave the clock untouched and move the
boundary, one multiplicatively and the other additively.

\subsection{Leak: a Straight Boundary in an Expanding Clock}
Multiply \eqref{eq:bawl-traj} by $\e^{ku}$ and define the rescaled state and
clock
\begin{equation}
  Y(u) = \e^{ku}X(u),
  \qquad
  r(u) = \int_0^{u}\e^{k\theta}\dd\theta = \frac{\e^{ku}-1}{k}.
\end{equation}
Then $Y = z + Vr$ exactly: a static start point and a constant drift, which is
an LBA. The boundary inherits the rescaling,
\begin{equation}
  B_{\mathrm{leak}}(r) = b\,\e^{ku(r)} = b\left(1 + kr\right),
  \label{eq:bawl-boundary}
\end{equation}
so \textbf{the leak model is an LBA against a straight boundary of slope $kb$,
in a clock that expands exponentially.} Two straight lines race, and a line
beats a line if and only if its slope is larger, so the eventual-hit condition
is $V > kb$ with no dependence on $z$. That recovers \eqref{eq:bawl-mass}
geometrically and explains why start-point variability affects the timing of
responses but not which trials produce one. The clock is unbounded,
$r(\infty) = \infty$, so there is no endpoint; and because late crossings
correspond to near-parallel lines meeting at large $r$, while $r$ grows like
$\e^{ku}/k$, the real-time density inherits the $\e^{-ku}$ tail of
\eqref{eq:bawl-tail}.

A clock change alone does not suffice here. The gauge factor $\e^{ku}$ applied
to the state is essential, precisely because leak attacks the start point as
well as the accumulated evidence.

\subsection{Transient and sustained BTAwL clocks}
Each BTAwL subrace uses the leak gauge.  Under $\e^{ku}$ its trace becomes
$z_j+V_jr_j$ with
\[
  r_j(u)=\int_0^u\e^{k\theta}\eta_j(\theta)\dd\theta,
  \qquad
  B_j(r_j)=b\e^{ku(r_j)}.
\]
The boundary slope is $kb/\eta_j(u)$.  For the transient profile this slope
has a minimum at the pulse peak and the clock reaches its maximum at
$T_{\max}$, producing the frozen branch in \eqref{eq:btawl-crit}.  For the
sustained profile the slope tends to $kb$ and the clock continues, so there is
no finite endpoint.  The full local race combines the two independent
survivors as in \eqref{eq:btawl-local-race}; it does not create a third,
shared-drive clock.

\subsection{Drive Decay: a boundary in the accumulated-drive clock}
For every $(\rho,\gamma)$ take the clock $s=Q_\rho(u)$, with inverse
$u=u_\rho(s)$ on its range. Then \eqref{eq:bawd-traj} is the ordinary line
$z+Vs$ against
\begin{equation}
 B_{\rho,\gamma}(s)=b+\ell R_{\rho,\gamma}(u_\rho(s)).
\label{eq:bawd-boundary}
\end{equation}
The boundary is independent of $V$ and $z$. For $\gamma<1$ its tangency
with a line gives \eqref{eq:bawd-sat} and the finite endpoint
$T_{\max}=T_{\rm sat}(0)$; for $\gamma=1$, $B=b+\ell s$ is a straight
boundary and the clock may continue indefinitely, so no finite endpoint
occurs. When $\ell=0$ the boundary is flat and the model reduces to an LBA
in clock $Q_\rho$. The exponential constant-clearance boundary
$b-(\ell/k)\log(1-ks)$ is the special case $(\rho,\gamma)=(\infty,0)$.

\subsection{Proportional Clearance: a Flat Boundary on the Positive Clock}
For $0<\lambda<1$ and positive launches, take the clock to be the accumulated
waveform on its rising branch,
\begin{equation}
  s=q_{\lambda}(u)\in[0,q_{\lambda}^{\max}],\qquad 0\le u\le u_{*},
  \qquad \dot s=\e^{-ku}-\lambda.
\end{equation}
Then \eqref{eq:bawdp-traj} is $z+Vs$ against the \emph{original, flat}
boundary $b$. Thus the positive-launch proportional-clearance branch is an
ordinary LBA, with an ordinary boundary, in a clock that runs forward,
decelerates, and stops at $u_{*}$. All of this branch's content is in the
clock; the boundary carries none of it, where in the leak model the boundary
tilts and in the drive-decay model it diverges. The raw clock continues below
zero after $u_{*}$ when $\lambda>0$; that sign reversal is precisely what
creates the late negative-launch branch for an untruncated normal.
\begin{itemize}
  \item A positive-launch line $z+Vs$ that has not reached $b$ by clock time
        $q_{\lambda}^{\max}$ never will. That is the intrinsic non-finishing
        mass, and it is the LBA's own defect evaluated at a finite time rather
        than a new mechanism.
  \item There is no tangency, because there is no curvature to be tangent to.
        The positive-branch critical launch is the slope that just reaches
        $b-z$ within the available clock, \eqref{eq:bawdp-crit}, which is why
        it is affine in $z$ and why nothing multiplicative ever appears.
  \item The positive-branch endpoint is a property of the clock alone, so it
        is common to every trial and to every accumulator sharing $(k,\lambda)$,
        and it does not move when $b$ or $A$ is manipulated. An untruncated
        normal is the explicit exception globally because its negative branch
        responds after the second zero.
\end{itemize}
Read backwards, the clock distinguishes the models that warp it: it expands
without bound under leak, saturates but keeps running under drive decay, and
halts outright under proportional clearance. The last of these is the only one
that leaves the boundary alone. Symmetrically, the two variants of
Sections~\ref{sec:bawf} and \ref{sec:bawr} are the only ones that leave the
clock alone.

\subsection{Global Fading: a Multiplicative Boundary in Ordinary Time}
Divide \eqref{eq:bawf-traj} by the kernel. The clock is not warped at all:
\begin{equation}
  Y(u) = \frac{X(u)}{h_\rho(u)} = z + Vu,
  \qquad
  B_{\rm fade}(u) = b\,H_\rho(ku).
  \label{eq:bawf-boundary}
\end{equation}
So \textbf{the fading model is an LBA in ordinary time against a boundary that
recedes multiplicatively.} The gauge factor is applied to the state, as under
leak, which is why the start point is attenuated in both; but the clock is the
identity, which is why the two have different tails. The endpoint follows from
the geometry in one line: the line $z+Vu$ can touch $bH_\rho(ku)$ only while
the boundary's chord from the origin is still falling, and
$\dd/\dd u\left[bH_\rho(ku)/u\right]=0$ at $ku=x_{\max}$ regardless of $b$,
because $b$ multiplies the whole boundary and cancels from the tangency
condition. That is \eqref{eq:bawf-tmax}, and it is the geometric content of the
remark in Section~\ref{sec:bawf-endpoint}.

The excluded index is equally clear here: $H_1(x)=1+x$ is a straight boundary,
two straight lines meet at most once, and there is no tangency to freeze.

\subsection{Ramping Clearance: an Additive Boundary in Ordinary Time}
Move the exogenous drain to the other side of \eqref{eq:bawr-traj}:
\begin{equation}
  Y(u) = X(u) + \frac{\kappa u^{p+1}}{p+1} = z + Vu,
  \qquad
  B_{\rm poly}(u) = b + \frac{\kappa u^{p+1}}{p+1}.
  \label{eq:bawr-boundary}
\end{equation}
So \textbf{the ramping model is an LBA in ordinary time against a boundary
that rises additively.} No gauge factor is applied to the state, which is
exactly why the start point survives undiminished: $Y$ and $X$ differ by a
quantity free of $z$ and $V$. Comparing \eqref{eq:bawf-boundary} with
\eqref{eq:bawr-boundary} isolates the whole difference between the two new
variants, $b\,H_\rho$ against $b+\kappa u^{p+1}/(p+1)$, and with it the
difference in their endpoints. A multiplicative boundary cancels $b$ from the
tangency condition; an additive one does not.

The comparison with drive decay is the other one to make, since
\eqref{eq:bawr-boundary} and \eqref{eq:bawd-boundary} are both additively
rising boundaries. They differ in the clock: BAwD's rising boundary
$b+\ell R_{\rho,\gamma}$ lives in the saturating clock $Q_\rho$, so its
divergence is measured against a clock that is itself running out, and the two
effects are what \eqref{eq:bawd-sat} entangles. BAwR's rises in ordinary time
against a clock that never slows, so the endpoint is a pure statement about
when a straight line last touches a convex curve.

\subsection{Consequences, and a Warning About Diffusion}
\begin{center}
\small
\begin{tabular}{L{0.075\textwidth}L{0.155\textwidth}L{0.16\textwidth}L{0.14\textwidth}L{0.16\textwidth}L{0.095\textwidth}}
\toprule
Model & Transform & Clock range & Boundary & Failure to finish & Peak time \\
\midrule
BAwL & clock $r=(\e^{ku}-1)/k$ \emph{and} state rescaling $\e^{ku}$ &
$[0,\infty)$, expanding & straight, slope $kb$ &
slope race lost, $V\le kb$ & none \\
\addlinespace
BAwD & clock $s=Q_\rho(u)$ alone &
$[0,Q_\rho(\infty)]$; finite for $\rho>1$ and $\rho=\infty$, infinite for
$\rho=1$ &
$b+\ell R_{\rho,\gamma}$, diverging; straight if $\gamma=1$ &
peak below $b$ if $\gamma<1$; asymptotic threshold if $\gamma=1$ &
$u_{\rm peak}(V)$ \\
\addlinespace
BAwDp & clock $s=q_\lambda(u)$ alone &
$[0,q_\lambda^{\max}]$, halting at $u_{*}$; $\lambda=0$ approaches $1/k$
without a finite halt &
flat at $b$, unchanged &
clock runs out first; an untruncated normal adds a late negative branch &
common, $\log(1/\lambda)/k$ \\
\addlinespace
BAwF & state rescaling $h_\rho(u)^{-1}$ alone; the clock is untouched &
$[0,\infty)$, ordinary time & $b\,H_\rho(ku)$, receding multiplicatively &
staggered peak below $b$ & $u_{\rm peak}(V)$ \\
\addlinespace
BAwR & an additive shift $+\kappa u^{p+1}/(p+1)$; the clock is untouched &
$[0,\infty)$, ordinary time & $b+\kappa u^{p+1}/(p+1)$, rising additively &
staggered peak below $b$ & $(V/\kappa)^{1/p}$ \\
\addlinespace
BTAwL transient & clock $r_T=\int_0^u\e^{k\theta}\eta_T$ \emph{and} state
rescaling $\e^{ku}$ &
$[0,r_T(\infty)]$; finite when $k\tau_t<1$, unbounded otherwise &
$b\e^{ku}$, slope $kb/\eta_T(u)$ &
staggered peak below $b$; endpoint $T_{\max}$ &
$u_{\rm peak}(V_T)$ \\
\addlinespace
BTAwL sustained & clock $r_S=\int_0^u\e^{k\theta}\eta_S$ \emph{and} state
rescaling $\e^{ku}$ &
$[0,r_S(\infty)]$, unbounded for every $k\ge0$ &
$b\e^{ku}$, slope tending to $kb$ &
asymptotic plateau, no endpoint &
$u_{\rm peak}(V_S)$ \\
\bottomrule
\end{tabular}
\end{center}

Whether the induced start-point average closes in $\PhiN$ and $\phiN$ is a
separate question, answered model by model in Section~\ref{sec:closure} and
tabulated in Section~\ref{sec:closure-summary}. The mechanism table continues:

\begin{center}
\small
\begin{tabular}{L{0.062\textwidth}L{0.20\textwidth}L{0.185\textwidth}L{0.185\textwidth}L{0.185\textwidth}}
\toprule
Model & Decays & Decay depends on & Right tail & Start point \\
\midrule
BAwL & the accumulated state, start point included &
how much evidence is held & exponential, rate $k$ &
irrelevant to \emph{whether} a trial finishes \\
\addlinespace
BAwD & the input and its clearance; the state is a perfect integrator &
time since onset and the fixed $(\rho,\gamma)$ &
hard zero if $\gamma<1$; unbounded and $\rho$-dependent if $\gamma=1$ &
shifts both timing and finishing \\
\addlinespace
BAwDp & the input; the state is a perfect integrator &
time since onset & hard zero at $u_{*}$; unbounded for the signed normal
branch & shifts both timing and finishing \\
\addlinespace
BAwF & the readout; nothing is lost from the accumulation &
time since onset and the fixed $\rho$ & hard zero at $x_{\max}/k$ &
attenuated, and shifts both timing and finishing \\
\addlinespace
BAwR & nothing; an exogenous drain opposes the drive &
time since onset alone & hard zero at $T_{\max}$ &
shifts both timing and finishing \\
\addlinespace
BTAwL transient & the accumulated state, start point included, \emph{and} the
transient stimulus drive &
how much evidence is held, and time since onset &
hard zero at $\tau_t\psi(k\tau_t)$; $u\e^{-u/\tau_t}$ when $k=0$ &
forgotten as under leak; the threshold is staggered \\
\addlinespace
BTAwL sustained & the accumulated state, start point included, and the
sustained stimulus drive &
how much evidence is held, and time since onset &
unbounded, with asymptotic threshold $kb$ for $k>0$ &
forgotten as under leak; the plateau is start-point-free \\
\addlinespace
BTAwL local race & both independent subraces &
the two channel profiles and $\pi$ &
unbounded for $0<\pi\le1$; transient endpoint only at $\pi=0$ &
the winner is the minimum of the two subrace times \\
\bottomrule
\end{tabular}
\end{center}

Two consequences follow. First, within a single parameter setting these representations are
exact identities, so ``decaying gain'', ``expanding boundary'', and ``warped
clock'' are the \emph{same model}; they are equivalent descriptions of the
same response-time law. What distinguishes the descriptions mathematically is the induced
pair (clock, boundary): leak expands its clock, the positive proportional-
clearance branch halts its clock, a transient drive expands the clock and then
stalls it, drive decay either supplies a finite
accumulated clock or an unbounded one according to $\rho$, and the two
variants of Sections~\ref{sec:bawf} and \ref{sec:bawr} keep ordinary time and
move the boundary instead. The corresponding observable signatures range from
exponential tails to hard endpoints or matched-decay unbounded tails. Across
parameter settings, the choice of which quantity is held fixed determines
whether a change in $k$ in the leak parameterization is equivalent to a change
in $b$. The sharpest distinction is whether changing the boundary level moves
the response window. Within a single setting, window and boundary are two
descriptions of one law; across settings they separate BAwF from BAwR.

Second, this interchangeability is a property of \emph{ballistic} accumulation,
in which the only randomness is the launch strength and the start point. It
does not extend to diffusion models. A time change
$t \mapsto \int_0^t h(\theta)\dd\theta$ applied to a Brownian path also
rescales its quadratic variation, so a time-dependent gain on the drift induces
a time-dependent diffusion coefficient as well. Gain and boundary are then no
longer interchangeable, which is the familiar reason that urgency signals and
collapsing boundaries are distinguishable in diffusion models but degenerate
here. Conclusions drawn from \eqref{eq:bawl-boundary} or
\eqref{eq:bawd-boundary} must not be transported to diffusion accumulators.

%======================================================================
\section{Identifiability}
\label{sec:identifiability}
%======================================================================
Three separate identifiability issues arise, and they should be kept apart.

\paragraph{Scale.} The evidence axis has no units, so exactly one of
$(V, \ell, b, A)$ must be fixed. With a normal launch the usual $s_v = 1$
convention does this, because $s_v$ carries the units of the evidence axis.
With an ordinary BAw* lognormal launch it does not: $c_L$ is dimensionless
and rescaling shifts only the arithmetic mean $m_L$. With a split-lognormal or
BTAwL lognormal launch, rescaling shifts $\mu$ (and the split point for the
split form) while leaving the shape coordinates unchanged. With a Weibull
launch, the shape $r$ is dimensionless and rescaling shifts only the arithmetic
mean $m_W$. Fix $b$ for the leak model and read the remaining quantities as
dimensionless.

The drive-decay model may fix $\ell = 1$ in the rate chart. In the endpoint
chart of \eqref{eq:bawd-ell-from-tmax}, $\ell$ is derived instead: the map
$(V,\ell,b,A)\mapsto c\,(V,\ell,b,A)$ leaves $T_{\max}$ and $k$ invariant,
so the scale gauge acts through the internal log-scale location, $B$, and $A$;
for an ordinary BAw* launch the former is the log-mean coordinate implied by
$m_L$. One of those coordinates must therefore be fixed; fixing an intercept of $B$ preserves condition-wise
changes in the boundary offset.

The proportional-clearance model has no $\ell$ to fix, and $\lambda$ is
dimensionless, so it must use the boundary convention too: fix $b$, or $B$ when
$A = 0$. The same applies to BAwF, whose only decay parameter $k$ is a rate
and carries no evidence units. In all three cases the anchor has to be stated
explicitly, because there is nothing in the parameter set that does it
silently. BAwR is the exception among the new variants: $\kappa$ has the units
of drive per unit $u^{p}$, so fixing it works exactly as fixing $\ell$ once did
for drive decay.

The invariance is exact. Scaling $(V,\ell,b,A)$, $(V,\kappa,b,A)$, or
$(V,b,A)$ as appropriate leaves the response-time law unchanged, so one
evidence-scale coordinate must be fixed to remove the corresponding gauge.

\paragraph{Dynamical identifiability in the drive-decay family.} For
$\gamma<1$, the observable window $T_{\max}=T_{\rm sat}(0)$ is the natural
summary of the coupled rates $(k,\ell,\rho,\gamma)$. The dimensionless
saturation coordinate and the launch location relative to $V_c(0)$ provide
separate shape and omission information. For $\gamma=1$ there is no endpoint;
the limiting critical launch $\ell$ replaces the window as the relevant
scale. The finite grid of $(\rho,\gamma)$ should therefore be treated as
model indices, while continuous rates and launch parameters are identified
only through the kernels $Q_\rho,R_{\rho,\gamma}$ and the critical equation
\eqref{eq:bawd-sat}. Scale invariance still permits fixing one of
$(V,\ell,b,A)$.

There is a second, milder conditioning issue in the same family. With $\ell$
fixed, $k$ and the boundary jointly determine $T_{\max}$ through
\eqref{eq:bawd-sat}; the endpoint chart uses $T_{\max}$ as the direct
coordinate and leaves $\ell$ derived. The two charts describe the same
likelihood and differ only in coordinates.

\paragraph{Global fading.} The response window $x_{\max}/k$ depends on $k$
alone, while the omission boundary $V_c(0)$ is proportional to $kb$. The two
quantities therefore play separate structural roles, unlike the coupled
$(b,k,\ell)$ geometry of drive decay. The scale gauge still requires one of
the launch-scale location ($m_L$ for an ordinary BAw* launch or $\mu$ for a
split/BTAwL launch), $b$, or $A$ to be fixed.

\paragraph{Ramping clearance.} The parameter $\kappa$ sets the window through
\eqref{eq:bawr-tmax} and the omission threshold through
\eqref{eq:bawr-defect-window}. The exponent $p$ controls the curvature of the
boundary \eqref{eq:bawr-boundary}; as $p\to\infty$ at fixed $T_{\max}$, the
model approaches an LBA with a hard deadline. Thus $p$ primarily controls the
within-window shape rather than the endpoint itself.

\paragraph{BTAwL local race.} The transient subrace has the clean endpoint
summary $T_{\max}=\tau_t\psi(k\tau_t)$ and
$V_{c,T}(0)=kb/\eta_T(T_{\max})$, both free of $A$ and with the latter
proportional to $b$.  The sustained subrace has no endpoint and contributes
the asymptotic threshold $kb$ when $k>0$.  Consequently the full local race
has an endpoint only at $\pi=0$ when $k>0$; for $0<\pi\le1$ its shape and tail
depend on $\pi$, $\tau_t$, and $\tau_s$ through the product survival
\eqref{eq:btawl-local-race}, while the omission probability is the product
\eqref{eq:btawl-local-defect}.  In the fast-leak corner the transient window
approaches $\tau_t$ and its threshold approaches $\e kb$, but this does not
remove the sustained tail when $\pi>0$.  The evidence-scale gauge remains:
one of the launch location/scale and boundary/start-point scales must be
fixed.

\paragraph{Waveform degeneracy in the proportional-clearance model.} The same
issue appears in a milder form on the positive-launch branch. As $k \to 0$ the
model becomes an LBA with drive $(1-\lambda)V$, so $\lambda$ and the launch
scale confound exactly in that corner, and near it they confound badly. What
breaks the degeneracy is the positive-branch endpoint:
$u_{*} = \log(1/\lambda)/k$ is the positive-branch endpoint, and
$q_{\lambda}^{\max}$ enters its omission rate through
\eqref{eq:bawdp-defect}. These are the natural endpoint and defect
coordinates, analogous to $T_{\max}$ for the drive-decay model. With an
untruncated normal there is no global
endpoint: the late negative-launch component has support $(u_0,\infty)$ and its
mass must be included through \eqref{eq:bawdp-normal-defect}.

They do not, however, invert to $(k,\lambda)$ on their own. Their ratio
\begin{equation}
  \frac{q_{\lambda}^{\max}}{u_{*}}
  = \frac{1 - \lambda + \lambda\log\lambda}{\log(1/\lambda)}
\end{equation}
depends on $\lambda$ alone but is not monotone: it vanishes at both
$\lambda \to 0$ and $\lambda \to 1$ and peaks at $0.298$ near
$\lambda = 0.166$, so each attainable ratio below that peak is produced by two
values of $\lambda$. The two branches are separated by the shape of the clock
between its endpoints --- a small $\lambda$ with a small $k$ spends most of the
window near-linear, a large $\lambda$ with a large $k$ front-loads the
accumulation --- so the likelihood distinguishes them, but the endpoint pair
alone does not. Unlike the drive-decay model,
nothing here contaminates $b$: the endpoint is free of the boundary, so a
window estimate and a threshold estimate do not compete.

\paragraph{Rate and boundary coupling in the leak model.} The plateau
\eqref{eq:bawl-mass} depends on $k$ and $b$ only through $kb$, so the omission
rate alone cannot separate them. The tail rate \eqref{eq:bawl-tail} depends on
$k$ alone, so the two together do. Thus the leak parameter is separated from
the plateau by the slow-response tail; removing that tail weakens the
separation.

\paragraph{Degenerate limits.} In the fully reduced model of
Section~\ref{sec:reductions}, $b$ and the log-scale location enter a lognormal
launch only through $\log b - \mu$. For an ordinary BAw* launch this is
equivalently the ratio $b/m_L$ at fixed $c_L$; BTAwL and split-lognormal
launches use their directly fitted $\mu$. A Weibull launch uses the analogous
ratio $b/m_W$, and $\ell$ and $v$ enter a normal-launch drive-decay
model at $k=0$ only through $v - \ell$. At those corners the likelihood has a
flat direction; the exact invariance is visible directly from the reduced
equations.

%======================================================================
\section{Closure Summary}
\label{sec:closure-summary}
%======================================================================
\begin{center}
\small
\begin{tabular}{p{0.42\textwidth}p{0.24\textwidth}p{0.24\textwidth}}
\toprule
Case & Density & Distribution function \\
\midrule
LBA, $A = 0$, any launch law & pushforward, closed & pushforward, closed \\
LBA, $A > 0$, normal launch & closed & closed \\
LBA, $A > 0$, ordinary, split-lognormal, or Weibull launch & closed & closed (stop-loss difference) \\
Leak, $A \ge 0$, $k \ge 0$, normal launch & closed & closed \\
Leak, $A \ge 0$, $k \ge 0$, ordinary, split-lognormal, or Weibull launch & closed & closed \\
Leak with one-factor correlated launches & one-dimensional Gauss--Hermite & one-dimensional Gauss--Hermite \\
Drive decay, $A = 0$, any $(\rho,\gamma)$ and launch law & closed & closed \\
Drive decay, $A > 0$, live branch, any $(\rho,\gamma)$ and launch law & closed & closed \\
Drive decay, $A > 0$, $\gamma<1$, ordinary, split-lognormal, or Weibull launch &
closed (live and frozen) & closed (power/log or incomplete-gamma primitives) \\
Drive decay, $A > 0$, $\gamma<1$, normal launch &
closed (live part only) & one-dimensional quadrature on the frozen branch \\
Drive decay, $A > 0$, $\gamma=1$, any launch law &
closed & closed (affine start-point average) \\
Proportional clearance, positive launches, any $A \ge 0$, $0<\lambda<1$ &
closed, LBA density at $q_{\lambda}(u)$ times $q_{\lambda}'(u)$ &
closed, LBA distribution function at $\bar q_{\lambda}(u)$ \\
Proportional clearance, positive-branch omission mass &
--- & closed, $1-F_{\mathrm{LBA}}(q_{\lambda}^{\max})$ \\
Proportional clearance, untruncated normal &
closed, positive density plus the late density in
\eqref{eq:bawdp-negative-pdf} &
closed, the two-branch law \eqref{eq:bawdp-normal-full} \\
Proportional clearance, untruncated-normal defect &
--- & closed, \eqref{eq:bawdp-normal-defect} \\
Global fading, $A = 0$, any $\rho$ and launch law & closed & closed \\
Global fading, $A > 0$, live branch, any $\rho$ and launch law &
closed & closed \\
Global fading, $A > 0$, ordinary, split-lognormal, or Weibull launch &
closed (live and frozen) & closed (positive-power/log or incomplete-gamma primitives) \\
Global fading, $A > 0$, normal launch &
closed (live part only) & one-dimensional quadrature on the frozen branch \\
BTAwL transient, $A = 0$, any launch law & closed & closed \\
BTAwL transient, $A > 0$, live branch, any launch law &
closed, \eqref{eq:btawl-pdf} & closed \\
BTAwL transient, $A > 0$, frozen branch, any launch law &
--- & one-dimensional quadrature; the Lambert coordinate closes for no
launch law \\
BTAwL transient, $k = 0$, any $A \ge 0$, any launch law &
closed & closed; there is no frozen branch \\
BTAwL sustained, any $A \ge 0$, any launch law &
closed & closed; the start-point average is affine \\
BTAwL local race, $0<\pi<1$ &
component densities combined as in \eqref{eq:btawl-local-race} &
component CDFs combined as in \eqref{eq:btawl-local-race} \\
Ramping clearance, $A = 0$, any $p$ and launch law & closed & closed \\
Ramping clearance, $A > 0$, live branch, any $p$ and launch law &
closed & closed \\
Ramping clearance, $A > 0$, ordinary, split-lognormal, or Weibull launch &
closed (live and frozen) & closed (positive-power or incomplete-gamma primitive) \\
Ramping clearance, $A > 0$, normal launch &
closed (live part only) & one-dimensional quadrature on the frozen branch \\
\bottomrule
\end{tabular}
\end{center}

The rule behind the table is the one established in Section~\ref{sec:closure}.
Closed forms appear while the start-point integral stays an integral of the
launch survivor against $\dd w$, which is the stop-loss transform and is
available for all four launch laws. In the staggered frozen branch, finite $\rho$
instead induces combinations of $\dd w$ and $w^{-\alpha}\dd w$, while the
$\rho=1$ limit induces $\log(w/\ell)\dd w$. These power and logarithmic
primitives close for the ordinary and split-lognormal launches through Gaussian
tails and for Weibull launches through incomplete gamma functions, but generally
require quadrature for the normal launch. The familiar $\dd w/w$ obstruction is only the
$\alpha=1$ exponential constant-clearance corner; the scale-family statement
belongs to that special multiplicative coordinate, not to the whole family.

The transient BTAwL component marks the edge of the rule. Its frozen coordinate
is available through the Lambert inverse, but as a
$W$ of $w$ rather than as a power of it, so the induced measure
$\e^{ks(w)}K(s(w))\dd w$ is neither $\dd w$ nor $w^{\pm\alpha}\dd w$ and
none of the four launch laws integrates it.  The sustained component and the
full local race inherit the component closure through the product formulas in
\eqref{eq:btawl-local-race}.  This is a useful boundary case because it shows
that positive launch laws close the frozen branch only when the induced
coordinate is additive or a power of the launch.

The other two new variants extend the rule without changing it, and they make
its scope clearer. The sign of the exponent turns out to be irrelevant: BAwF's
frozen Jacobian \eqref{eq:bawf-frozen-jacobian} induces $w^{+1/(\rho-1)}\dd w$
and BAwR's \eqref{eq:bawr-frozen} induces $w^{+1/p}\dd w$, both positive powers
where drive decay produces negative ones, and all of them are closed truncated
moments for either log-scale or Weibull launch. What matters is not whether the frozen
coordinate is a positive or a negative power of $w$ but whether it is a power
of $w$ at all, because that is what the normal launch cannot integrate in
closed form and the log-scale and Weibull launches can. BAwR is the extreme case in one
respect: since $1/p$ is never
zero for finite $p$, it has no logarithmic member, and its frozen branch is a
single power primitive for every parameter value.

%======================================================================
\section{A Derivation Checklist}
\label{sec:checklist}
%======================================================================
For deriving a new member of this family, the following order works.
\begin{enumerate}
  \item Write the generator and solve it for $X(u)$ given $(V,z)$. If the
        solution is not available in closed form, the model is not ballistic in
        the sense used here and nothing below applies.
  \item Check whether the solution \emph{separates}, $X(u) = z + V\,c(u)$ for
        some deterministic $c$ free of $(V,z)$. If it does, the model is an LBA
        in the clock $c$ against the original boundary, and everything below is
        already answered: the endpoint is where $c$ stops increasing, the
        distribution function is $F_{\mathrm{LBA}}(\bar c)$ with $\bar c$ the
        running maximum, the density carries the extra factor $c'$, and closure
        is inherited from the LBA for any launch law. Separation fails as soon
        as the decay attacks the state, as in \eqref{eq:bawl-gen}, or as soon as
        a term in the generator is free of $V$, as in \eqref{eq:two-generators}.
  \item Solve $X(u) = b$ for the required launch strength $V^{*}(u,z)$.
  \item Check whether $V^{*}$ is monotone in $u$. If it is, the hit set is an
        upper tail of the launch law at $V^{*}(u,z)$ and the support is
        unbounded. If it is U-shaped, the hit set freezes at
        $V_c(z) = \min_u V^{*}(u,z)$, which produces a hard right endpoint, a
        flat distribution function beyond it, and an intrinsic non-finishing
        mass.  The transient BTAwL subrace has one such minimum and freezes
        after its endpoint; the full BTAwL local race instead combines that
        component with the sustained survivor through
        \eqref{eq:btawl-local-race}, so it has no shared mixed trajectory.
  \item If there is an endpoint, ask whether the boundary appears in it. When
        the trace has the multiplicative form $a(u)z+q(u)V$ with no additive
        term, the endpoint is $\arg\max q$ and $b$ cancels; any additive term
        surviving at $z=0$ breaks the cancellation and the window becomes a
        function of the boundary level. This identifies whether the endpoint is
        boundary-dependent.
  \item Identify $V_{\mathrm{crit}}(z)$ and evaluate the defect
        \eqref{eq:generic-defect} before doing anything else. It is usually a
        one-line calculation and it tells you what kind of model you have.
  \item Check whether $V^{*}$ is affine in $z$. If it is, the start-point
        average is a stop-loss difference and closes for any launch law with a
        closed-form partial expectation. If it is not, find the coordinate in
        which the critical launch strength is simple, and check what measure
        that coordinate induces on launch strengths. An additive coordinate
        gives $\dd w$ and closes generically. A finite power-law coordinate
        gives combinations of $\dd w$ and $w^{-\alpha}\dd w$, while its
        $\rho=1$ limit gives $\log(w/\ell)\dd w$; these close for ordinary,
        split-lognormal, and Weibull launches but generally require quadrature
        for a normal launch. Anything
        outside those two shapes, such as the Lambert coordinate of the
        transient BTAwL frozen branch, closes for no launch law, and quadrature
        is intrinsic to the model's transformed coordinate.
  \item Differentiate for the density, remembering that a frozen branch is
        constant in $u$ and contributes nothing. The density often closes when
        the distribution function does not.
  \item Check the reductions: $A \to 0$ should give a pushforward, and each
        decay parameter going to zero should give back the LBA or a named
        special case. A limit that does not reduce cleanly usually indicates an
        algebra error in $V^{*}$.
  \item Check the scale invariance and fix exactly one quantity. Confirm which
        parameters appear only in combinations, and state those combinations.
\end{enumerate}

%======================================================================
\section{Further Reading}
%======================================================================
The components used above are standard. For the ballistic idea and the original
closed-form derivation, the entry point is the
\href{https://en.wikipedia.org/wiki/Linear_ballistic_accumulator}{linear
ballistic accumulator}, which was introduced as a tractable simplification of
the leaky competing accumulator. For the full leaky competing system, including
lateral inhibition and rectification that we deliberately omit, see the
literature on
\href{https://en.wikipedia.org/wiki/Leaky_integrate-and-fire}{leaky
integration} and competing accumulator models of choice. For the stochastic parent of
Section~\ref{sec:btawl}, the sources are Smith's (1995) two-channel model of
visual simple reaction time, the racing change and level detectors of Burbeck
and Luce (1982) that it formalizes, and the single-accumulator alternative of
Ratcliff and Van Dongen (2011); the psychophysics behind the sustained and
transient profiles is the linear-systems treatment of visual temporal
sensitivity, for which
\href{https://en.wikipedia.org/wiki/Temporal_contrast_sensitivity_function}{temporal
contrast sensitivity} and the impulse response of a
\href{https://en.wikipedia.org/wiki/Low-pass_filter}{low-pass filter} cascade
are the entry points. The stop-loss
transform \eqref{eq:stoploss} is the
\href{https://en.wikipedia.org/wiki/Black-Scholes_model}{Black--Scholes} call
price under an ordinary lognormal; the split-lognormal counterpart is
\eqref{eq:split-stoploss}, while the Weibull counterpart is the incomplete-gamma
form in \eqref{eq:weibull-tail-moments}. The general object is the partial
expectation or first-order lower partial moment of a distribution. The saturation condition
\eqref{eq:bawd-sat} is inverted by a general $\Theta_{\rho,\gamma}$ map. Only
in the special exponential constant-clearance corner
$(\rho,\gamma)=(\infty,0)$ does it reduce to
\href{https://en.wikipedia.org/wiki/Lambert_W_function}{Lambert $W$} on its
lower branch.  The same branch inverts the transient omission threshold through
$s=-\tau_tW_{-1}(-kb/w)$.  For the invariance argument in Section~\ref{sec:closure}, the
relevant background is the
\href{https://en.wikipedia.org/wiki/Haar_measure}{Haar measure} of the
multiplicative group on $(0,\infty)$ and the notion of a
\href{https://en.wikipedia.org/wiki/Location-scale_family}{location-scale
family}. Diffusion counterparts of these models, where the interchangeability
of gain and boundary fails, are treated through
\href{https://en.wikipedia.org/wiki/Hitting_time}{first-passage} theory and the
\href{https://en.wikipedia.org/wiki/Inverse_Gaussian_distribution}{inverse
Gaussian} law.

\end{document}
